\documentclass[a4paper,11pt,openany]{ltjsbook}
\usepackage{luatexja}
\usepackage{luatexja-fontspec}
\usepackage{amsmath,amssymb,amsfonts,amsthm}
\usepackage{geometry}
\geometry{top=24mm,bottom=24mm,left=24mm,right=24mm,footskip=10mm}
\usepackage{graphicx}
\usepackage{tcolorbox}
\tcbuselibrary{skins,breakable}
\usepackage{hyperref}
\usepackage{enumitem}
\usepackage{booktabs}
\usepackage{fancyhdr}

\definecolor{primaryblue}{RGB}{20, 60, 110}
\definecolor{softblue}{RGB}{238, 244, 252}
\definecolor{darkslate}{RGB}{35, 45, 60}
\definecolor{warmgold}{RGB}{185, 125, 20}
\definecolor{softgold}{RGB}{254, 249, 235}
\definecolor{accentcrimson}{RGB}{160, 30, 45}

\hypersetup{
    colorlinks=true,
    linkcolor=primaryblue,
    citecolor=darkslate,
    urlcolor=accentcrimson,
    pdftitle={誰にもわかる微積分 —— 原著完全直訳・数学校閲版},
    pdfauthor={シルバヌス・P・トンプソン}
}

\pagestyle{fancy}
\fancyhf{}
\fancyhead[LE,RO]{\thepage}
\fancyhead[RE]{\leftmark}
\fancyhead[LO]{\rightmark}
\renewcommand{\headrulewidth}{0.4pt}

% Compatibility macros for Thompson TeX original
\newcommand{\lintertext}[1]{\intertext{#1}}
\newcommand{\rintertext}[1]{\intertext{#1}}
\newcommand{\ds}{\displaystyle}
\newcommand{\tb}{\par\vspace{2mm}\hrule width 0.3\textwidth\vspace{2mm}\par}
\newcommand{\DStrut}{}
\newcommand{\efrac}[2]{\frac{#1}{#2}}
\DeclareMathOperator{\cotan}{cotan}
\DeclareMathOperator{\Arcsin}{arc\,sin}
\DeclareMathOperator{\Arccos}{arc\,cos}
\DeclareMathOperator{\Arctan}{arc\,tan}
\DeclareMathOperator{\sech}{sech}
\newcommand{\Pageref}[2][p.]{#1~\ref{#2}}
\newcommand{\Pagerange}[2]{pp.~\ref{#1}--\ref{#2}}
\newcommand{\Fig}[1]{Fig.~\ref{fig:#1}}
\newcommand{\Figs}[3]{Figs.~\ref{fig:#1}#2\ref{fig:#3}}
\newcommand{\NB}{\textit{N.B.}}
\newcommand{\IE}{\textit{i.e.}}

\newcommand{\TmpColA}{}
\newcommand{\TmpColB}{}
\newcommand{\PadTo}[3][]{#3}

\begin{document}

% --- 表紙 ---
\begin{titlepage}
    \centering
    \vspace*{25mm}
    {\fontsize{14}{20}\selectfont \textbf{プロジェクト・グーテンベルク古典数学翻訳・校閲シリーズ 第5弾}}\\[
8mm]
    {\fontsize{28}{36}\selectfont \textbf{誰にもわかる微積分}}\\[4mm]
    {\fontsize{16}{22}\selectfont \textit{Calculus Made Easy (1914)}}\\[12mm]
    \rule{0.85\textwidth}{1.5pt}\\[10mm]
    {\fontsize{15}{22}\selectfont \textbf{シルバヌス・P・トンプソン} 著}\\[3mm]
    {\fontsize{12}{18}\selectfont （Silvanus P. Thompson, D.Sc., F.R.S.）}\\[8mm]
    {\fontsize{13}{19}\selectfont 原著全21章・序文・プロローグ・エピローグ・全演習問題解答 完全直訳}\\[2mm]
    {\fontsize{11}{16}\selectfont \textbf{Qwen2.5-7B-Instruct 忠実翻訳 ／ Qwen2.5-Math-7B-Instruct 数理解析校閲}}\\[25mm]
    \vfill
    {\fontsize{11}{16}\selectfont 2026年10月 電子書籍版}\\[2mm]
    {\fontsize{10}{14}\selectfont \textbf{古典数理復刻プロジェクト}}\\[15mm]
\end{titlepage}
% --- まえがき・解題 ---
\frontmatter
\chapter*{解題：『誰にもわかる微積分』の不朽の意義}
\addcontentsline{toc}{chapter}{解題：『誰にもわかる微積分』の不朽の意義}
シルバヌス・フィリップス・トンプソン（Silvanus Phillips Thompson, 1851--1916）による『Calculus Made Easy』（1910年初版、1914年改訂第2版）は、1世紀以上にわたって世界中で愛読され続けてきた、微積分学教育史上屈指の名著である。
当時のイギリスにおける数学教育は、厳密主義と抽象的な形式主義に偏重し、初学者や工学を志す若者たちにとって、微積分は「一部の天才だけが操る難解な暗号」として恐れられていた。物理学者・電気工学者であったトンプソンは、この状況に強い義憤を覚え、F.R.S.（王立協会フェロー）という自らの権威を隠して匿名（「ある無知なる者」）で本書を世に問うた。
本書を貫く根本思想は、プロローグに掲げられた古代の格言に集約される：
\begin{quote}
\textbf{「一人の人間にできることは、他の人間にもできる（What one fool can do, another can.）」}
\end{quote}
トンプソンは、$dx$ や $dy$ を神秘的な不可知の対象とみなすことを排し、単なる「少しばかりの微小量」として親しみやすく直観的に解き明かした。高次の微小量をなぜ切り捨ててよいのか、和差積商や連鎖律、極大極小問題、そして定積分による面積の求値法から微分方程式に至るまで、初等数学の素朴な知性をもって一歩一歩確実に登り詰める構成となっている。
本書は、プロジェクト・グーテンベルク収録の第2版完全TeX原著（EBook \#33283）に基づき、本文・数式・演習問題解答に至る全26セクションを一語の省略もなく忠実に日本語へと翻訳し、現代最高峰のオープン推論モデルである \texttt{Qwen2.5-7B-Instruct} による全訳および \texttt{Qwen2.5-Math-7B-Instruct} による数理・LaTeX検証を経て制作された完全直訳版である。
\tableofcontents
\mainmatter
\chapter*{序文：第2版への序文}
\addcontentsline{toc}{chapter}{序文：第2版への序文（Preface to the Second Edition）}
\markboth{序文：第2版への序文}{序文：第2版への序文}
\noindent\textbf{\large 【副題】多くの初学者が微積分を克服した驚くべき成功の記録}\\[4mm]
第2版の序文
著者のこの作品が予想以上に成功したことに伴い、著者は多くの計算例と練習問題を追加しました。また、経験から有用だと感じた部分についてはさらに詳述することにより、その部分を拡大修正しました。
著者は、教師、学生、そして批判家からの多くの貴重な提案や手紙に対して感謝の意を表します。
1914年10月
\cleardoublepage
目次
プロローグ xi
I. 初期の恐怖から解放される 1
II. 小ささの異なる程度について 3
III. 相対的な成長について 9
IV. 最単純なケース 18
V. 次の段階。定数に対する取り扱い 26
VI. 和、差、積、商 35
VII. 微分の連続性 49
VIII. 時間が変化するとき 52
IX. 有用な手順を導入する 67
X. 微分の幾何学的意味 76
XI. 最大値と最小値 93
XII. 曲線の曲率 112
XIII. 他の有用な手順 121
XIV. 真実の複合利子と有機的な成長の法則について 134
XV. 正弦と余弦に対する取り扱い 165
XVI. 部分微分 175
XVII. 不定積分 182
XVIII. 積分は微分の逆過程である 191
XIX. インテグラルを用いて面積を求めること 206
XX. 有用な手順、障害と勝利 226
XXI. 解の一部を求める方法 234
XXII. 終章と寓話 249
標準形の一覧表 252
練習問題解答 254
\chapter*{プロローグ：プロローグ}
\addcontentsline{toc}{chapter}{プロローグ：プロローグ（Prologue）}
\markboth{プロローグ：プロローグ}{プロローグ：プロローグ}
\noindent\textbf{\large 【副題】「一人の人間にできることは、他の人間にもできる」}\\[4mm]
多くの愚か者が計算ができるという事実を考えれば、いかに容易なことでもあるのか、また誰もがそれらのトリックを学ぶことが困難であると考えられるべきではない。
微分積分のトリックはいくつか簡単なものがある。他には非常に難しいものもある。高度な数学の教科書を書く愚かな賢者たちは、容易な計算がいかに容易かを説明することを怠ることが多い。むしろ彼らは、最も困難な方法でそれを行うことで、自らの驚異的な知性をあなたに印象付けることを望んでいるように見える。
自分自身が非常に愚かな人間であるため、自分自身から難しさを解きほぐさなければならなかった。そして、他の愚か者たちのために、難しい部分ではなく容易な部分だけを紹介するので、乞うご期待。これらを完全にマスターすれば、残りは自然と理解できるだろう。一つの愚かな者ができることは、他の愚かな者もできる。
主部
\phantomsection\pdfbookmark[-1]{主部}{主部}
\chapter{予備的恐怖からの解放}
\noindent\textit{\large To deliver you from the Preliminary Terrors}\\[2mm]
\noindent\textbf{\small 【章の論理とテーマ】dx と dy という記号の真の意味・難解そうに見える数学記号の正体}\\[4mm]
第一章：第五形式の生徒を恐れから解放する
最初に、計算を学ぶ試みすら諦める第五形式のboysを悩ませる初期の恐怖は、計算において主に使用される二つの主要な記号の意味を、一般的な感覚で述べることによって一挙に取り除くことができます。
これらの恐ろしい記号は以下の通りです：
(1) $d$ は単に「少しだけ」を意味します。
従って $dx$ は $x$ の少しだけ、または $du$ は $u$ の少しだけを表します。一般的な数学では、「要素の」という言葉を使い、「少しだけの」という表現よりも丁寧に聞こえるかもしれませんが、どちらでも構いません。しかし、これらの少しだけ（または要素）は無限に小さいと考えることができます。
(2) $\displaystyle \int$ は単に長い $S$ の形をしていますし、「合計」の意味で呼び出すこともできます。
従って $\displaystyle \int dx$ はすべての $x$ の少しだけの合計、または $\displaystyle \int dt$ はすべての $t$ の少しだけの合計を表します。一般的な数学では、「積分」と呼ぶことが多いですが、この記号は「合計」の意味で理解しても構いません。単語「積分」は単に「全体」を意味するだけです。時間の長さが1時間である場合を考えると、もし1時間を $3600$ の小さな部分（それぞれを $dx$ と呼びます）に分割するとします。これらの $3600$ の小さな部分をすべて加算すると、それはすべての $dx$ の合計（つまり $x$ 自身）になります。
この記号が最初に見られるとき、あなたは今後、その指示に従って計算する操作（可能であれば）を行なうことを知るでしょう。それが示す小さな部分を次の記号によって表しているすべての少しだけを合計することです。
それだけのことです。
\chapter{小ささの異なる階級について}
\noindent\textit{\large On Different Degrees of Smallness}\\[2mm]
\noindent\textbf{\small 【章の論理とテーマ】微小量と高次の微小量・無視してよい微小部分の幾何学的・代数的根拠}\\[4mm]
我々は計算の過程で様々な程度の小さな量と対処することになるでしょう。
また、どの条件下で小さな量が非常に小さいものとして考慮に入れずに済むかについても学ぶ必要があります。すべては相対的な小ささによって決まります。
規則を定める前に、いくつかの例を考えましょう。1時間には60分があり、1日に24時間がありますので、1週間には7日があります。したがって、1日に1440分、1週間に10080分があります。
明らかに、1分は一週間と比べると非常に小さな時間量です。実際、祖先たちは1時間と比べてそのような程度の小ささを持つものとして考え、「1分」と呼んでいました。これは1時間の六十分の一という意味で、小さな分数を指します。時間のさらに細かい分割が必要になったときには、各分を60つのより小さな部分に分けました。イリノーズ女王時代には、「2次の分」（2次の小ささを持つ量）と呼んでいました。現在ではこれらの小さな量は「秒」と呼ばれていますが、なぜそう呼ばれるのか知る人は少ないでしょう。
もし1分が一日と比べると非常に小さいのであれば、1秒はさらにどれだけ小さいか！
また、四便士と王冠を比較すると、四便士はわずかに1000分の1よりも価値があります。四便士の差は王冠と比べるとほとんど重要ではありませんので、そのような小さな量として考慮に入れることができます。しかし、四便士と$p$1000を比較すると、四便士は$p$1000に対してどれほど重要な量であるかを考えれば、四便士は王冠に対する$p$1000の比に比べてさらに無視できる量となります。すなわち、億万長者の財産に対して金の王冠も相対的に無視できる量となるのです。
今、任意の数的分数を、何らかの目的のために相対的に小さな比として固定すれば、より高度に小さな分数を述べることができます。たとえば、時間のためには $\frac{1}{60}$ を小さな分数と呼びますが、その $\frac{1}{60}$ (小さな分数の小さな分数)は二番目の順序の小さな量（Pagelabel{smallness}）として扱うことができます。
あるいは、何かの目的のために $1$％（つまり $\frac{1}{100}$）を小さな分数と定義すれば、その $1$％ (つまり $\frac{1}{10,000}$) は二番目の順序の小さな量の小さな分数となります。そして、$\frac{1}{1,000,000}$ は三番目の順序の小さな量であり、これは $1$％ の $1$％ の $1$％ です。
最後に、非常に精密な目的のために $\frac{1}{1,000,000}$ を「小さな量」と見做すとします。たとえば、一流のクロノメーターが一年で半分の一分以内の誤差を許さないなら、それは $1$ 部分中 $1,051,200$ の正確さで時間を取り扱う必要があります。この目的のために $\frac{1}{1,000,000}$ (つまり百万分の一) を小さな量と見做すなら、その $\frac{1}{1,000,000}$ の $\frac{1}{1,000,000}$ は $\frac{1}{1,000,000,000,000}$ (つまり兆分の一)となり、比較して無視することができます。
したがって、小さな量自体が小さいほど、その二番目の順序の小さな量もより無視できることがわかります。したがって、第一番目の順序の小さな量を十分に小さく取れば、すべての場合において、第二番目（または第三番目、またはそれ以上の）順序の小さな量を無視することは正当化されます。
ただし、注意しなければならないのは、小さな量が他の因子と乗算される場合、その他の因子が大きいなら重要な役割を果たす可能性があるということです。わずかな銭でも何百倍に掛かると重要になります。
今、微分積分においては $dx$ を $x$ の小さな部分と書き表します。このような $dx$, $du$, $dy$ などのものは「微分」と呼ばれ、$x$ の微分や $u$ の微分、または $y$ の微分と呼びます。（これらを「ディーイックス」、「ディーユウ」、「ディーワイ」と読みます。）もし $dx$ が $x$ の小さな部分であり、それ自体も相対的に小さいなら、$x \cdot dx$, $x^2 \, dx$, $a^x \, dx$ のような量は無視できるとは限りません。しかし、$dx \times dx$ は無視できるもので、二階の小さな量となります。
非常に単純な例を挙げて説明しましょう。
$x$ を成長によって少しだけ増加し、$x + dx$ となる量と考えます。ここで $dx$ は増加した小さな部分です。この量の平方は $x^2 + 2x \cdot dx + (dx)^2$ となります。二番目の項は無視できないもので、一階の量である一方、三番目の項は二階の小さな量であり、$x^2$ の一部です。したがって、もし $dx$ を数値的に $\frac{1}{60} x$ と仮定すると、二番目の項は $\frac{2}{60} x^2$ となりますが、三番目の項は $\frac{1}{3600} x^2$ となります。この最後の項は明らかに二番目の項よりも重要ではありません。しかし、さらに $dx$ を $\frac{1}{1000} x$ と仮定すると、二番目の項は $\frac{2}{1000} x^2$ となりますが、三番目の項は $\frac{1}{1,000,000} x^2$ となります。
幾何学的には次のように表すことができます：
$x$ を辺とする正方形を描きます（図1）。この正方形の大きさが $dx$ の分だけ増加するとします。増大した正方形は元の正方形 $x^2$, 上部と右側にある各長方形（それぞれの面積は $x \cdot dx$ または合計で $2x \cdot dx$）と、右上隅にある小さな正方形 $(dx)^2$ の組み合わせです。図2では $dx$ を $x$ の約 $\frac{1}{5}$ として描いています。しかし、もし $dx$ を $x$ の約 $\frac{1}{100}$ と仮定すると（細いペンで引いた線の厚さ程度）、右上隅の小さな正方形の面積は $\frac{1}{10,000} x^2$ となり、ほとんど見えません。明らかに $(dx)^2$ は $dx$ が十分小さい場合に無視できるものです。
類似の例を考えてみましょう。
億万長者が秘書に対して次のように言います：
次の週、私に入ってくるお金の一部をあなたに少しの割合で与えるつもりだ。そして、秘書が少年に対して言うかもしれません：私が得たものから少しだけの割合をあなたに与えるつもりだ。それぞれの場合において、その割合を100分の1とする。もし億万長者が次の週$p$1000を受け取った場合、秘書は$p$10、少年は2シリング受け取ります。$p$10は$p$1000と比較して小さな量ですが、2シリングは非常に二番目に小さい量で、その上さらに小さな量です。
もし割合が1000分の1に設定された場合、どのような不均衡が生じるでしょうか？すると、億万長者が$p$1000を受け取った一方で、秘書は$p$1、少年はほとんどない銭を受けることになります！
巧みなデーン・スウィフト\footnote{『詩論：狂詩』（p.20）、1733年に印刷されましたが、しばしば誤引用されています。}
は一度次のように書いています：
``自然学者たちは観察するが、 flea
その上に寄生する小さな fleas がいる。
そしてそれらにも flea がいて、
そして無限に続く。''
牛は通常のサイズのフライには悩むかもしれません——第一級の小さな生物です。しかし、牛はフライのフライには気を取られることはないでしょう；第二級の小さな生物であるため、それは無視できるほど小さくなります。すらとらのフライのフライの一打も、牛にとっては大きな問題ではありません。
\chapter{相対的増大について}
\noindent\textit{\large On Relative Growings}\\[2mm]
\noindent\textbf{\small 【章の論理とテーマ】変化率・微分係数 dy/dx の直観的把握と相互関係}\\[4mm]
第三章 相対的な増大について
すべての微分法において、我々は常に増大する量とその増大率を扱っています。これらの量を二つのクラスに分類します：固定値であると考えられる定数と、変動可能な変数です。固定値として考えられ、定数と呼ばれるものは、アルファベットの頭文字から符号化して表すことが多いので、$a$, $b$, または $c$ などです；一方で増大する能力があると考えられるもの、つまり数学的には「変動」するものについては、アルファベットの末尾の文字、たとえば $x$, $y$, $z$, $u$, $v$, $w$、あるいは時として $t$ などで表します。
また、通常は一度に複数の変数を扱い、一方の変数が他方の変数にどのように依存するかを考えます。例えば、投射体がその高さに達するまでの時間によってその高さがどのように変化するかを考える場合があります。あるいは、与えられた面積を持つ長方形について考え、その長さが増加したときに幅はどの程度減少するかを調べることもあります。また、梯子の傾斜角度がどのように変化すると、到達する高さもどのように変化するかを考えます。
二つの変数 $x$ と $y$ があり、これらは互いに依存しています。一方の変数が変化すると、その依存関係により他方の変数も変化します。一つの変数を $x$ と呼び、もう一つの依存する変数を $y$ とします。
$x$ を変化させることを考えます。つまり、$x$ の値を変えたり、あるいはそれを想像的に変更したりします。これには $dx$ という微小な増分を加えることを意味します。この操作により $x$ は $x + dx$ になります。すると、$x$ の変化に伴い $y$ も変化し、$y + dy$ となります。ここで $dy$ は $dx$ と同じ大きさではなく、正の値や負の値を取る可能性があります。奇迹的な状況以外では、$dy$ と $dx$ の大きさは同じではありません。
(1) 直角三角形（図4）の底辺と高さをそれぞれ$x$と$y$とする。その他の辺の傾きは固定されており、角度は30$^\circ$である。この三角形が拡大され、角度が最初と同じまま保たれる場合、底辺が$x + dx$に増加すると、高さも$y + dy$に増加する。ここで、$x$の増加により$y$も増加する。高さが$dy$、底辺が$dx$の小さな三角形は元の三角形と相似であり、明らかに$\frac{dy}{dx}$の値は$\frac{y}{x}$の値と同じである。角度が30$^\circ$であるため、
\[
\frac{dy}{dx} = \frac{1}{1.73}.
\]

(2) 図5において、水平方向から壁までの距離を$x$とし、固定長さのラダーABの底端Aから壁までの距離を表す。そして、高さ$y$はこの$x$に依存している。もし底端Aを少し壁から離すと、上端Bも少しだけ下に来ることを見ることができる。科学的な言葉で述べると、$x$が増加して$x + dx$になると、$y$は$y - dy$となる；つまり、$x$に正の増分を与えると、その結果として$y$にも負の増分が生じる。

しかし、どの程度か？ラダーが非常に長い場合を考えると、底端Aが壁から19インチ離れたとき、上端Bは地面からちょうど15フィートに達する。もし底端Aをさらに1インチ離すと、上端Bはどのくらい下に来るだろうか？すべてインチで表すと、$x = 19$インチ、$y = 180$インチである。ここで、$dx$は$1$インチの増分であり、つまり$x + dx = 20$インチ。

$y$はどのくらい減少するか？新しい高さは$y - dy$となる。ユークリッドI.47を使用して高さを計算すると、$dy$の値を求めることができる。ラダーの長さは
\[
\sqrt{(180)^2 + (19)^2} = 181 \text{インチ}.
\]
したがって、新しい高さ$y - dy$は
\begin{align*}
(y - dy)^2 &= (181)^2 - (20)^2 = 32761 - 400 = 32361, \\
y - dy     &= \sqrt{32361} = 179.89 \text{インチ}.
\end{align*}
ここで$y$は180インチであるため、$dy$は$180 - 179.89 = 0.11$インチとなる。

したがって、$dx$に1インチの増分を与えると、$dy$には0.11インチの減少が生じる。

そして \( dy \) に対する \( dx \) の比は次のように述べることができます：
\[
\frac{dy}{dx} = - \frac{0.11}{1}.
\]

また、(特定の位置を除いて) \( dy \) が \( dx \) と異なる大きさであることも容易に理解できます。

微分法全体を通じて、我々は探求し続けている奇妙なものがあります。それはただの比であり、\( y \) と \( x \) の関係があるとき、つまり \( x \) が変化すると \( y \) も変化するときに、\( dy \) に対する \( dx \) の比です。

ここで注意すべきは、この比 \( \frac{dy}{dx} \) を求められるのは \( y \) と \( x \) が何かしらの関係にあるときだけです。つまり、\( x \) が変化すると \( y \) も同時に変化します。例えば、最初に取り上げた例では、三角形の底辺 \( x \) を長くすると、三角形の高さ \( y \) も大きくなります。また、二番目の例では、梯子の足が壁から離れる距離 \( x \) が増えると、梯子が達する高さ \( y \) は逆比例して減少します。最初はゆっくりですが、\( x \) が大きくなるにつれてますます速くなります。

これらの状況では、\( x \) と \( y \) の関係は完全に定義されており、数学的に表現することができます。それはそれぞれ \(\frac{y}{x} = \tan 30^\circ\) および \( x^2 + y^2 = l^2 \)（\( l \) は梯子の長さ）であり、そのときの \( \frac{dy}{dx} \) の意味が明確になります。

もし、以前通り$x$が壁から梯子の足の距離であるとき、$y$が高さではなく、壁の水平長さ、またはその中的の砖石数、あるいは建設以来の年数であるとする。この場合、$x$の変化は$y$に全く影響を与えない。この状況では$\frac{dy}{dx}$は意味を持たず、それを求める方法も存在しない。

異なる微分量$dx$, $dy$, $dz$などを用いる際には、$x$, $y$, $z$など間に何らかの関係性が存在することが暗黙に含意されている。この関係性は$x$, $y$, $z$などの「関数」と呼ばれる。上の2つの式$\frac{y}{x} = \tan 30^\circ$と$x^2 + y^2 = l^2$は$x$と$y$の関数であり、これらは暗黙的に$x$を$y$で表すか、または$y$を$x$で表す方法を含んでいるため、「暗黙的な関数」と呼ばれる。これらの式はそれぞれ次のように書き換えることができる：

\begin{align*}
y &= x \tan 30^\circ \quad\text{または}\quad x = \frac{y}{\tan 30^\circ} \\
\intertext{そして}
y &= \sqrt{l^2 - x^2} \quad\text{または}\quad x = \sqrt{l^2 - y^2}.
\end{align*}

これらの最後の式は$x$を$y$で明示的に、または逆に$y$を$x$で表す値を明らかにしているため、「明示的な関数」と呼ばれる。例えば$x^2 + 3 = 2y - 7$は暗黙的な関数であり、これは$y = \frac{x^2 + 10}{2}$（$x$の明示的な関数）または$x = \sqrt{2y - 10}$（$y$の明示的な関数）と書き換えることができる。明示的な関数は$x$, $y$, $z$などが変化するとき、あるいはこれらが同時に変化するときにその値も変化することを意味する。

このため、明示的な関数の値は「依存変数」と呼ばれる。これは関数内の他の変数量の値によって依存しているからである。これらの他の変数は「独立変数」と呼ばれ、それらの値は関数が取る値から決定されるわけではない。例えば、$u = x^2 \sin \theta$の場合、$x$と$\theta$は独立変数であり、$u$は依存変数である。

いくつかの量 \( x \), \( y \), \( z \) の間に正確な関係が分かっていない場合や、その関係を明示するのが不便な場合があります。これらの変数間には何かしらの関係があることが知られており、それぞれの変数 \( x \), \( y \), \( z \) を単独に変更すると他の量にも影響を与えるため、\( x \), \( y \), \( z \) の関数が \( F(x, y, z) \)（暗黙の関数）や \( x = F(y, z) \), \( y = F(x, z) \) または \( z = F(x, y) \) ）（明示的関数）で表されます。有时会使用字母 \( f \) 或者 \( \phi \) 来代替 \( F \)，因此 \( y = F(x) \), \( y = f(x) \) 和 \( y = \phi(x) \) 都表示同样的意思，即 \( y \) 的值依赖于 \( x \) 的值，这种关系没有明确说明。

私たちは \( y \) に対する \( x \) の微分係数を「\(\dfrac{dy}{dx}\)」と呼びます。これはこの非常に単純な概念の正式名称です。しかし、物自体が簡単であることを恐れることはありません。代わりに、長い言葉で名付けられた愚かさを呪うだけにして、その本質的な部分、すなわち \(\dfrac{dy}{dx}\) の比を扱います。

学校で学んだ通常の代数では、常に未知の量 \( x \) または \( y \) を探していました。あるいは、同時に二つの未知量を探さなければなりませんでした。あなたは今、新しい方法で狩りをする必要があります；狐が \( x \) や \( y \) のどちらでもありません。その代わりに、この奇妙な仔犬 \( \dfrac{dy}{dx} \) を探します。\( \dfrac{dy}{dx} \) の値を求める過程は「微分」と呼ばれます。しかし、覚えておいてください。求められるのは、それぞれの \( dy \) と \( dx \) が無限に小さい場合におけるこの比の真の値です。微分係数の正確な値は、それぞれが無限に小さな場合において近づく値です。

今では \( \dfrac{dy}{dx} \) を探す方法を学びましょう。

第III章への注記

微分を読む方法について

\( dx \) が \( d \) 乗 \( x \) の意味ではなく、\( d \) は要素や少しだけの何かを表すための符号であることを忘れないでください。\( dx \) を「ディー・エックス」と読みます。

読者がこのような事項について誰かの指導を受けない場合、以下のように説明します。微分係数は次のように読みます。
\[
\dfrac{dy}{dx} \text{は「デーイー・ワイ バイ デーイー・エックス」または「デーイー・ワイ オーバー デーイー・エックス」と読みます。}
\]
また
\[
\dfrac{du}{dt} \text{は「デーイー・ユー バイ デーイー・ティ」}
\]

二階微分係数は後で出会うでしょう。それらは次のような形です：
\[
\dfrac{d^2 y}{dx^2}; \text{これは「デーイー・ツー・ワイ オーバー デーイー・エックス スクウェアード」と読み、}
\]
これは$y$を$x$について二度微分することを意味します。

関数が微分されたことを示す別の方法は、関数の記号にアクセントを付けることです。したがって、$y=F(x)$とし、$y$は$x$の未定義な関数である（参照：第\pageref{function}節）とすると、$\dfrac{d\bigl(F(x)\bigr)}{dx}$ではなく$F'(x)$と書くことができます。同様に、$F''(x)$は元の関数$F(x)$を$x$について二度微分することを意味します。



\chapter{最も単純なケース}
\noindent\textit{\large Simplest Cases}\\[
2mm]
\noindent\textbf{\small 【章の論理とテーマ】べき関数 $y = x^n$ の微分の基本公式とその第一原理からの導出}\\[4mm]
第IV章 最も単純なケース
1. まず、初等的な式 \( y = x^2 \) の微分法について見ていきましょう。
1. まず、単純な式 \( y = x^2 \) を考えてみます。微積分の基本的な考えは「成長」です。数学的には「変化」と呼ばれます。\( y \) と \( x^2 \) が等しいので、\( x \) が増加すれば \( x^2 \) も増加し、そして \( x^2 \) が増加すれば \( y \) も増加します。私たちが必要とするのは、\( y \) の増加と \( x \) の増加との関係を求めることが目的です。言い換えれば、\( dy \) に対する \( dx \) の比、すなわち簡潔に表すと \( \dfrac{dy}{dx} \) の値を求めることです。
次に、\( x \) が少しだけ増加して \( x + dx \) となることを考えてみます。同様に、\( y \) も少しだけ増加し、\( y + dy \) となります。明らかに、拡大された \( y \) は拡大された \( x \) の平方と等しいです。これを式で表すと：
\[
y + dy = (x + dx)^2.
\]

次に、平方を計算すると：

\[
y + dy = x^2 + 2x \cdot dx + (dx)^2.
\]

ここで、\( (dx)^2 \) とは何か？ \( dx \) は \( x \) の少しだけ、すなわち小さな量を意味します。したがって、\( (dx)^2 \) はさらに小さな量の少しだけを表し、上記（第\pageref{smallness}節）で説明したように、二番目に小さい量と考えることができます。従って、他の項と比べて非常に小さく無視することができます。この項を取り除き、式は次のようになります：

\[
y + dy = x^2 + 2x \cdot dx.
\]

次に、\( y = x^2 \) を方程式から引くと：

\[
dy = 2x \cdot dx.
\]

最後に、両辺を \( dx \) で割ると：

\[
\dfrac{dy}{dx} = 2x.
\]

今、我々が求めようとしているものは以下の比です：
$$\frac{dy}{dx}$$
これは $y$ を $x$ に関して微分することから得られます。微分とは導関数を見つけ出すことを意味します。例えば，ある他の $x$ の関数 $u = 7x^2 + 3$ があった場合、この関数を $x$ に関して微分すると $\frac{du}{dx}$ を求めることになります。つまり，$\frac{d(7x^2 + 3)}{dx}$ を求めることです。他方で，時間が独立変数であるようなケースも考えられます（参照：\pageref{indvar}）。例えば：
$$y = b + \frac{1}{2} at^2$$
この関数を $t$ に関して微分すると，$\frac{dy}{dt}$ を求めることになります。つまり，$\frac{d(b + \frac{1}{2} at^2)}{dt}$ を求めることです。

我々が探しているものは，$y$ の増加量と $x$ の増加量の比です。現在の場合においては，その比は $2x$ であることが分かりました。

\subsection*{数値例}
例えば，$x = 100$ で $\therefore y = 10,000$ とします。次に，$x$ が $101$ まで増加する（つまり，$dx = 1$）とします。その結果，拡大された $y$ は $101 \times 101 = 10,201$ となります。しかし，二番目の次元の小さな量を無視することに同意すれば，$1$ は $10,000$ と比較して無視できるので，拡大された $y$ を $10,200$ として丸めることができます。$y$ は $10,000$ から $10,200$ に増加し，追加された部分は $dy$ であり，これは $200$ となります。

$$\frac{dy}{dx} = \frac{200}{1} = 200.$$ 前の段落の代数計算によれば，$\frac{dy}{dx} = 2x$ が得られます。そしてこれが正しい；$x = 100$ のとき，$2x = 200$ です。

しかし，あなたは言うでしょう，私たちは一単位を無視しました。

では，$dx$ をさらに小さな量にします。試しに $dx = \frac{1}{10}$ とします。すると $x + dx = 100.1$ で，
\[
(x+dx)^2 = 100.1 \times 100.1 = 10,020.01.
\]

最後の桁 $1$ は $10,000$ の百万分の一であり，無視できるので，$10,020$ として丸めることができます。これにより $dy = 20$ となり，$$\frac{dy}{dx} = \frac{20}{0.1} = 200,$$ これはまた $2x$ と一致します。

\par\medskip\noindent\textbf{ケース~2}
同様の方法で $y = x^3$ を微分してみましょう。

$y$ が $y + dy$ に，$x$ が $x + dx$ に増加するとします。すると，
\[
y + dy = (x + dx)^3.
\]

立方を計算すると，
\[
y + dy = x^3 + 3x^2 \cdot dx + 3x(dx)^2 + (dx)^3.
\]

現在、$dy$と$dx$がともに十分に小さくなるとき、$(dx)^2$や$(dx)^3$は比較してさらに小さくなるため、これらを無視することができます。よって、

$$y + dy = x^3 + 3x^2 \cdot dx.$$

ここで$y = x^3$なので、

\begin{align*}
dy &= 3x^2 \cdot dx, \\
\frac{dy}{dx} &= 3x^2.
\end{align*}

次に、$y = x^4$を微分してみましょう。$y$と$x$が少し増加すると、

$$y + dy = (x + dx)^4.$$

右辺の四乗を計算すると、

$$y + dy = x^4 + 4x^3 \cdot dx + 6x^2(dx)^2 + 4x(dx)^3 + (dx)^4.$$

比較してより高い次数の項は無視できるので、

$$y + dy = x^4 + 4x^3 \cdot dx.$$

元の$y = x^4$を引くと、

\begin{align*}
dy &= 4x^3 \cdot dx, \\
\frac{dy}{dx} &= 4x^3.
\end{align*}

これらの例はすべて簡単です。結果を集めて、一般の法則を見つけることができるかどうか確認しましょう。$y$の値とそれに応じた$\frac{dy}{dx}$の値を二列に並べると、

\[
\begin{array}{|@{\quad}c@{\quad}|@{\quad}l@{\quad}|}
\hline
y   & \DStrut\dfrac{dy}{dx} \\
\hline
x^2 & 2x  \\
x^3 & 3x^2  \\
x^4 & 4x^3  \\
\hline
\end{array}
\]

これらの結果を見ると、$x$のべき乗を1 decrementし、同じ数（元々のべき乗と同じ数）で掛けるという操作が微分に見受けられるようです。例えば最後の例では$x^4$が$x^3$に減少しています。これを一度見た後は、他の項もどのように進むか予想できるでしょう。$x^5$を微分すると$5x^4$になるはずです。または、$x^6$を微分すると$6x^5$になります。疑問がある場合は、これらのうちの一つを試してみてください。

次に$y = x^5$を試みます。
\begin{align*}
y + dy &= (x + dx)^5 \\
       &= x^5 + 5x^4 \cdot dx + 10x^3(dx)^2 + 10x^2(dx)^3 + 5x(dx)^4 + (dx)^5.
\end{align*}

すべての高次の小量を無視すると、次のように残ります：
\begin{align*}
y + dy &= x^5 + 5x^4 \, dx, \\
\lintertext{\rlap{そしてこれを引くと}}
y &= x^5 \text{ から} \\
dy &= 5x^4 \, dx, \\
\lintertext{従って}
\frac{dy}{dx} &= 5x^4, \text{ これは私たちの仮定と一致する。}
\end{align*}

次に、論理的に観察を展開すると、任意の高次のべき乗---それを $n$ と呼ぶ---に対しても同じ方法で扱えることになるでしょう。
\begin{align*}
\lintertext{\indent まず}
y &= x^n, \\
\intertext{とすると、期待される結果は}
\frac{dy}{dx} &= nx^{(n-1)}.
\end{align*}

例えば、$n=8$ の場合、$y=x^8$ であり、微分すると $\dfrac{dy}{dx} = 8x^7$ となります。

実際、$x^n$ を微分すると $nx^{n-1}$ となる規則は、$n$ が正の整数の場合すべてに適用されます。[二項定理を使って $(x + dx)^n$ を展開すればすぐにわかります。]しかし、$n$ が負または分数の値を持つ場合においてこの規則が成り立つかどうかはさらに検討が必要です。

\subsection*{負のべき乗の場合}
$y = x^{-2}$ の場合、次のように進めてみましょう：
\begin{align*}
y + dy &= (x + dx)^{-2} \\
      &= x^{-2} \left(1 + \frac{dx}{x}\right)^{-2}.
\end{align*}

二項定理を使って展開すると、次のようになります：
\begin{align*}
&= x^{-2} \left[1 - \frac{2 \, dx}{x} + \frac{2(2+1)}{1 \times 2} \left(\frac{dx}{x}\right)^2 - \text{等} \right] \\
&= x^{-2} - 2x^{-3} \cdot dx + 3x^{-4}(dx)^2 - 4x^{-5}(dx)^3 + \text{等.} \\
\lintertext{\indent したがって、高次の小量を無視すると}
y + dy &= x^{-2} - 2x^{-3} \cdot dx.
\intertext{元の $y = x^{-2}$ を引くと}
dy &= -2x^{-3}dx, \\
\frac{dy}{dx} &= -2x^{-3}.
\end{align*}
これは先に導いた規則と一致しています。

\subsection{分数指数の場合}
$y = x^{\frac{1}{2}}$ とする。すると、以前と同様に、

\begin{align*}
y + dy &= (x + dx)^{\frac{1}{2}} = x^{\frac{1}{2}}
        \left(1 + \frac{dx}{x} \right)^{\frac{1}{2}} \\
     &= \sqrt{x} + \frac{1}{2} \frac{dx}{\sqrt{x}} - \frac{1}{8}
        \frac{(dx)^2}{x\sqrt{x}} +
        \text{($dx$ の高次の項)}
% blank removed
\end{align*}

元の $y = x^{\frac{1}{2}}$ を引いて、高次の項を省略すると、

\[
dy = \frac{1}{2} \frac{dx}{\sqrt{x}} = \frac{1}{2} x^{-\frac{1}{2}} · dx,
\]

よって $\frac{dy}{dx} = \frac{1}{2} x^{-\frac{1}{2}}$ となり、一般の法則と一致する。

\par\medskip\noindent\textbf{まとめ。} これまでどれだけ進んだか見てみよう。以下のようなルールに到着した：


$x^n$ の微分は、指数を掛けて指数を1減らすことで $nx^{n-1}$ を得る。

{I}（解答は \pageref{AnsEx:I} を参照）

次の関数を微分せよ：

\begin{enumerate}
\item (1) $y = x^{13}$
\item (2) $y = x^{-\frac{3}{2}}$


\item (3) $y = x^{2a}$
\item (4) $u = t^{2.4}$


\item (5) $z = \sqrt[3]{u}$
\item (6) $y = \sqrt[3]{x^{-5}}$


\item (7) $u = \sqrt[5]{\frac{1}{x^8}}$
\item (8) $y = 2x^a$ 


\item (9) $y = \sqrt[q]{x^3}$
\item (10) $y = \sqrt[n]{\frac{1}{x^m}}$

$x$ のべきの微分はいかに簡単かを学んだ。それはどれほど簡単か！

\end{enumerate}



\chapter{次の段階：定数の扱い方}
\noindent\textit{\large Next Stage. What to do with Constants}\\[
2mm]
\noindent\textbf{\small 【章の論理とテーマ】加法定数と乗法定数・定数の微分がゼロになる理由}\\[4mm]
次章では、定数の存在が微分過程に及ぼす影響について探ります。定数とは、$x$ または $y$ の値が変化しても変わらない数です。
\#\#\# 加わる定数
まず、加わる簡単な定数のケースから始めましょう：
\begin{align*}
y &= x^3 + 5.
\end{align*}
$x$ が $x+dx$ に増大すると仮定し、それに従って $y$ も $y+dy$ に増大します。
\begin{align*}
y + dy &= (x + dx)^3 + 5 \\
       &= x^3 + 3x^2 \, dx + 3x(dx)^2 + (dx)^3 + 5.
\end{align*}
高次の小量を無視すると、式は次のように簡略化されます：
\begin{align*}
y + dy &= x^3 + 3x^2 \cdot dx + 5. \\
dy &= 3x^2 \, dx. \\
\frac{dy}{dx} &= 3x^2.
\end{align*}
このように、定数 $5$ は完全に消えました。それは $x$ の増大に寄与せず、微分係数には含まれません。$7$ や $700$ や他の任意の数を代入しても同じ結果になります。
もし加わる定数が負の値であれば（例えば $-5$ または $-b$）、同様に消えます。
\#\#\# 乗じられる定数
次に、乗じられる簡単な定数のケースを考えてみましょう：
Let \( y = 7x^2 \).
Then on proceeding as before, we get:
\begin{align*}
y + dy &= 7(x+dx)^2 \\
       &= 7\{x^2 + 2x·dx + (dx)^2\} \\
       &= 7x^2 + 14x·dx + 7(dx)^2. \\
\intertext{Then, subtracting the original \( y = 7x^2 \), and neglecting
the last term, we have}
dy &= 14x·dx.\\
\frac{dy}{dx} &= 14x.
\end{align*}
Let us illustrate this example by working out the graphs of the equations \( y = 7x^2 \) and \( \frac{dy}{dx} = 14x \), by assigning to \( x \) a set of successive values, \( 0 \), \( 1 \), \( 2 \), \( 3 \), etc., and finding the corresponding values of \( y \) and of \( \frac{dy}{dx} \).
These values we tabulate as follows:
\[
\begin{array}{|c||*{6}{r|}|*{3}{r|}}
\hline
x & 0 & 1 & 2 & 3 & 4 & 5 & -1 & -2 & -3 \\
\hline
y & 0 & 7 & 28 & 63 & 112 & 175 & 7 & 28 & 63 \\
\hline
\dfrac{dy}{dx} & 0 & 14 & 28 & 42 & 56 & 70 & -14 & -28 & -42 \\
\hline
\end{array}
\]

Now plot these values to some convenient scale, and we obtain the two curves, \Fig{6} and \Fig{6a}.

Carefully compare the two figures, and verify by inspection that the height of the ordinate of the derived curve, \Fig{6a}, is proportional to the slope of the original curve, \Fig{6}, at the corresponding value of \( x \). To the left of the origin, where the original curve slopes negatively (that is, downward from left to right), the corresponding ordinates of the derived curve are negative.

Now if we look back at \Pageref{diffexample}, we shall see that simply differentiating \( x^2 \) gives us \( 2x \). So that the differential coefficient of \( 7x^2 \) is just 7 times as big as that of \( x^2 \). If we had taken \( 8x^2 \), the differential coefficient would have come out eight times as great as that of \( x^2 \). If we put \( y = ax^2 \), we shall get
\[
\frac{dy}{dx} = a \times  2x.
\]

もし最初に \( y = ax^n \) から始めたら、導関数は
\[
\dfrac{dy}{dx} = a\times nx^{n-1}
\]
となるべきでした。つまり、定数倍が行われた場合でも、微分するとただの定数倍に戻るということが分かります。そして、乗算に関する真理は除算についても同様に成り立ちます：例えば、先ほどの例で \( 7 \) を代わりに \(\frac{1}{7}\) として計算した場合でも、微分すると同じ \(\frac{1}{7}\) が結果に出力されます。

\par\textbf{【例題 さらにいくつかの例。】}\par
以下は完全に計算された追加の例で、これらを経験することで、普通の代数式に対する導関数の過程を完全にマスターし、この章末尾にある問題を自分で解くことができます。

(1) \( y = \dfrac{x^5}{7} - \dfrac{3}{5} \) を微分せよ。

\(\dfrac{3}{5}\) は加算された定数で消えます（参照：\Pageref{addconst}）。

したがって、
\[
\dfrac{dy}{dx} = \dfrac{1}{7} \times  5 \times  x^{5-1}
\]
または
\[
\dfrac{dy}{dx} = \dfrac{5}{7} x^4.
\]

(2) \( y = a\sqrt{x} - \dfrac{1}{2}\sqrt{a} \) を微分せよ。

項 \(\dfrac{1}{2}\sqrt{a}\) は加算された定数で消えます；そして、\( a\sqrt{x} \) は指数形式で \( ax^{\frac{1}{2}} \) と書けるため、
\[
\dfrac{dy}{dx}
  = a \times  \dfrac{1}{2} \times  x^{\frac{1}{2}-1}
  = \dfrac{a}{2} \times  x^{-\frac{1}{2}}
\]
または
\[
\dfrac{dy}{dx} = \dfrac{a}{2\sqrt{x}}.
\]

(3) \( ay + bx = by - ax + (x+y)\sqrt{a^2 - b^2} \) において、\( y \) を \( x \) についての微分係数を求めよ。

通常、このような表現はこれまでに習得した知識よりも少し深い理解が必要となるでしょう；しかし、常に式をより単純な形に変形できるかどうか試してみる価値があります。

まず最初に、\( y = \) \( x \) だけを含む何かの表現に持っていく必要があります。

この式は以下のようになります：
\[
(a-b)y + (a + b)x = (x+y) \sqrt{a^2 - b^2}.
\]

squaring, 我々は次を得る：
\[
(a-b)^2 y^2 + (a+b)^2 x^2 + 2(a+b)(a-b)xy = (x^2+y^2+2xy)(a^2-b^2),
\]
これは簡略化して
%
\begin{align*}
(a-b)^2 y^2 + (a+b)^2 x^2 &= x^2(a^2 - b^2) + y^2(a^2 - b^2); \\
\lintertext{または}
[(a-b)^2 - (a^2 - b^2)]y^2 &= [(a^2 - b^2) - (a+b)^2]x^2, \\
\lintertext{すなわち}
2b(b-a)y^2 &= -2b(b+a)x^2;
\end{align*}
\begin{gather*}
\lintertext{従って}
y = \sqrt{\frac{a+b}{a-b}} x \quad\text{および}\quad \frac{dy}{dx} = \sqrt{\frac{a+b}{a-b}}.
\end{gather*}

(4) 半径 $r$ と高さ $h$ の円柱の体積は公式で与えられる：
\[
 V = \pi r^2 h. 
\]
$r = 5.5$ インチ、$h=20$ インチの場合に、半径に対する体積の変化率を求めよ。また、$r=h$ の場合、体積が $400$ 立方インチ増えるようにする円柱の寸法を求めよ。

体積 $V$ に対する半径 $r$ の変化率は
\[
\frac{dV}{dr} = 2 \pi r h.
\]

$r = 5.5$ インチ、$h=20$ インチの場合、これは $690.8$ となる。半径が $1$ インチ変化すると体積が $690.8$ 立方インチ変化することを意味する。これについては、$r = 5$ インチと $r = 6$ インチの場合の体積はそれぞれ $1570$ 立方インチと $2260.8$ 立方インチであり、$2260.8 - 1570 = 690.8$ であることが容易に確認できる。

また、
\[
r=h, \quad \frac{dV}{dr} = 2\pi r^2 = 400 \quad \text{および}\quad r = h = \sqrt{\frac{400}{2\pi}} = 7.98 \text{ インチ}.
\]

(5) フェリーの放射温度計の読み取り $\theta$ は、観測体の摂氏温度 $t$ と関連付けられ、その関係は
\[
\frac{\theta}{\theta_1} = \left(\frac{t}{t_1}\right)^4,
\]
ここで $\theta_1$ は既知の温度 $t_1$ に対応する読み取りである。

800$^\circ$C、1000$^\circ$C、1200$^\circ$C の温度での温度計の感受性を比較せよ。温度が 1000$^\circ$C 時に読み取りが 25 であったことから。

感度は温度に対する読み取りの変化率であり、すなわち $\frac{d\theta}{dt}$ である。式は書き換えて
\[
\theta = \frac{\theta_1}{t_1^4} t^4 = \frac{25t^4}{1000^4},
\]
と表せる。したがって
\[
\frac{d\theta}{dt} = \frac{100t^3}{1000^4} = \frac{t^3}{10,000,000,000}.
\]

$t=800$、$1000$、および $1200$ のとき、$\frac{d\theta}{dt}$ はそれぞれ $0.0512$、$0.1$、および $0.1728$ となる。

熱敏度は約$800^\circ$から$1000^\circ$までで2倍になり、さらに$1200^\circ$まで上昇すると3分の4になります。

演習 II（解答は \pageref{AnsEx:II} を参照してください。）

次の関数を微分せよ：

\begin{enumerate}
\item (1) $y = ax^3 + 6$。
\item (2) $y = 13x^{\frac{3}{2}} - c$。


\item (3) $y = 12x^{\frac{1}{2}} + c^{\frac{1}{2}}$。
\item (4) $y = c^{\frac{1}{2}} x^{\frac{1}{2}}$。


\item (5) $u = \dfrac{az^n - 1}{c}$。
\item (6) $y = 1.18t^2 + 22.4$。

自分で他の例題を作成し、それらの微分を試みてみよ。

\item (7) 鉄棒の長さが$t^\circ\mathrm{C}$と$0^\circ\mathrm{C}$でそれぞれ$l_t$と$l_0$であるとき、$l_t = l_0(1 + 0.000012t)$です。$1^\circ\mathrm{C}$あたりの長さの変化を求めよ。

\item (8) 赤外電球の燐光力$c$が電圧$V$に従って$c = aV^b$で与えられ、$a$と$b$は定数である。$c$が$V$に対してどのように変化するかを求めよ。また、$a = 0.5 \times 10^{-10}$、$b=6$の電球において、$80$、$100$、$120$ボルトでの燐光力の変化を計算せよ。

\item (9) 直径$D$、長さ$L$、比重$\sigma$で張られた力$T$で引張られた弦の振動数$n$は
\[
n = \dfrac{1}{DL} \sqrt{\dfrac{gT}{\pi\sigma}}
\]
で与えられる。$D$、$L$、$\sigma$、および$T$がそれぞれ単独に変化したときの$n$の変化率を求めよ。

\item (10) 管が収縮しないで支持できる最大外部圧力$P$は
\[
  P = \left(\dfrac{2E}{1-\sigma^2}\right) \dfrac{t^3}{D^3}
\]
で与えられる。ここで、$E$と$\sigma$は定数であり、$t$は管の厚さ、$D$は直径である。（この式は$4t$が$D$に対して小さい場合に適用されるものとする。）

$P$が薄さの微小な変化と直径の微小な変化によりどのように変化するかを比較せよ。

\item (11) 半径$r$に関する次の量の変化率を求めよ：

\item[(a)] 半径$r$の円の周長；
\item[(b)] 半径$r$の円の面積；
\item[(c)] 斜辺長$l$の円錐の側面積；
\item[(d)] 半径$r$、高さ$h$の円錐の体積。

\item[(e)] 半径 \( r \) の球体の面積；

\item[(f)] 半径 \( r \) の球体の体積．

\item[(12)] 温度 \( T \) における鉄棒の長さ \( L \) は \( L = l_t \bigl[1 + 0.000012(T - t)\bigr] \) で与えられるとき、温度 \( T \) が変化するときに適した収縮用鉄輪の直径 \( D \) の変化率を求めよ．
\end{enumerate}



\chapter{和・差・積および商の微分法}
\noindent\textit{\large Sums, Differences, Products and Quotients}\\[
2mm]
\noindent\textbf{\small 【章の論理とテーマ】複合関数の基本微分則・積の微分法 (u·v) と商の微分法 (u/v)}\\[4mm]
VI 総和、差、積と商
1. 我々は単純な代数関数 $x^2 + c$ または $ax^4$ の微分法を学んできましたが、今度は二つ以上の関数の総和について考えなければなりません。
例えば，
\[
y = (x^2+c) + (ax^4+b);
\]
この関数の $\dfrac{dy}{dx}$ は何か？新しい問題にどのように取り組むべきか？

この質問に対する答えは簡単です：それぞれを順次微分すればよいのです：
\[
\dfrac{dy}{dx} = 2x + 4ax^3.\quad (\textit{答}．)
\]

もし疑問があるなら、基本原理から導くより一般的な例を試してみてください。

$y = u+v$ とします。ここで $u$ は $x$ の任意の関数であり、$v$ は別の $x$ の関数です。$x$ が $dx$ 増加すると、$y$ も $dy$ 増加し、$u$ も $du$ 増加し、$v$ も $dv$ 増加します。

したがって，
\begin{align*}
  y+dy &= u+du + v+dv. \\
\intertext{\indent ここで元の関係式 $y = u+v$ を引くと}
dy &= du+dv, \\
\intertext{さらに両辺を $dx$ で割ると}
\dfrac{dy}{dx} &= \dfrac{du}{dx} + \dfrac{dv}{dx}.
\end{align*}

この式は微分法の正当性を示しています。それぞれの関数を別々に微分し、結果を足すだけです。したがって、先ほどの例に戻り、二つの関数の値を入れると，
\begin{alignat*}{2}
\frac{dy}{dx} & = \frac{d(x^2+c)}{dx} &&+ \frac{d(ax^4+b)}{dx} \\
  & = 2x                  &&+ 4ax^3,
\end{alignat*}
と一致します。

もし三つの $x$ の関数 $u$, $v$, $w$ がある場合、
\begin{align*}
y &= u+v+w; \\
\lintertext{すると}
\frac{dy}{dx} &= \frac{du}{dx} + \frac{dv}{dx} + \frac{dw}{dx}.
\end{align*}

差分については、すぐに分かります。もし関数 $v$ に負の符号がついていた場合、その微分係数も負になります。したがって，
\begin{align*}
y &= u-v, \\
\lintertext{とすると}
\frac{dy}{dx} &= \frac{du}{dx} - \frac{dv}{dx}.
\end{align*}

しかし、積について考えると、事情は少し複雑になります。

まず、式
\[
y = (x^2+c) \times  (ax^4+b)
\]
を微分する場合、どのような手順で進めるべきか？

その結果は決して
\[
2x \times  4ax^3
\]
ではないでしょう。なぜなら、$c \times  ax^4$ や $x^2 \times  b$ がその積に含まれるとは限らないからです。

この問題には二つの方法があります。

まず第一法：最初に乗算を行い、その後微分します。

したがって、$x^2 + c$ と $ax^4 + b$ を乗算します。

これにより
\[
ax^6 + acx^4 + bx^2 + bc
\]
となります。

これを微分すると
\[
\dfrac{dy}{dx} = 6ax^5 + 4acx^3 + 2bx
\]
が得られます。

第二法：最初の原理に戻り、$y = u \times  v$ の式を考えます。ここで $u$ は $x$ の関数であり、$v$ もまた $x$ の関数です。すると、$x$ が $x+dx$ に増加し、$y$ が $y+dy$ に増加し、$u$ が $u+du$ に増加し、$v$ が $v+dv$ に増加するときには
\begin{align*}
 y + dy &= (u + du) \times  (v + dv) \\
        &= u · v + u · dv + v · du + du · dv.
\end{align*}

ここで $du · dv$ は二番目の微小量であり、極限において無視できるため
\[
y + dy = u · v + u · dv + v · du
\]
となります。

元の式 $y = u· v$ を引くと
\[
dy = u · dv + v · du
\]
となり、これを $dx$ で割ると
\[
\dfrac{dy}{dx} = u\, \dfrac{dv}{dx} + v\, \dfrac{du}{dx}
\]
が得られます。

これは次の指示を意味します：
\emph{二つの関数の積を微分するには、それぞれの関数に他方の関数の微分係数を乗じて両者を加える}。

この過程は以下のように解釈できます：$u$ を定数として $v$ を微分し、その後 $v$ を定数として $u$ を微分します。全体の微分係数 $\dfrac{dy}{dx}$ はこれらの二つの結果の合計となります。

このルールを見つけてから、先ほど考慮した具体的な例に適用します。

式
\[
(x^2 + c) \times  (ax^4 + b)
\]
を微分したいです。

$(x^2 + c)$ を $u$ とし、$(ax^4 + b)$ を $v$ とします。

次に、先ほど確立した一般の法則により、以下の式を書くことができます：
\begin{alignat*}{2}
\dfrac{dy}{dx} &= (x^2 + c) \, \frac{d(ax^4 + b)}{dx} &&+ (ax^4 + b) \, \frac{d(x^2 + c)}{dx} \\
              &= (x^2 + c) \, 4ax^3                  &&+ (ax^4 + b) \, 2x \\
              &= 4ax^5 + 4acx^3                     &&+ 2ax^5 + 2bx,  \\
%
\dfrac{dy}{dx} &= 6ax^5 + 4acx^3                     &&+ 2bx,
\end{alignat*}
これと全く同じ結果が得られます。

最後に、商の微分法について考えます。
例として $y = \dfrac{bx^5 + c}{x^2 + a}$ を考えてみましょう。このような場合、$x^2 + a$ で $bx^5 + c$ を割る前に計算を試みることは無意味です。なぜなら、$x^2 + a$ は $bx^5 + c$ で割り切れないし、共通の因数もありません。そこで最初の原理に戻って、法則を見つけるしか方法はありません。

したがって、
\[
y = \frac{u}{v};
\]
とします。ここで $u$ と $v$ は独立変数 $x$ の異なる関数です。すると $x$ が $x + dx$ に変化するとき、$y$ は $y + dy$ に、$u$ は $u + du$ に、$v$ は $v + dv$ にそれぞれ変化します。よって
\[
y + dy = \frac{u + du}{v + dv}.
\]

次に代数的な除算を行います：
\begin{align*}
& \TmpColA \;
  \PadTo[l]{\TmpColB}{u + du}
  \begin{tabular}{|c@{}}
    $\dfrac{u}{v} + \dfrac{du}{v} - \dfrac{u· dv}{v^2}\DStrut$\\\hline
  \end{tabular} \\[
-2.ex]
& \phantom{\TmpColA}
  \TmpColB \\
& \phantom{\TmpColA \;u + {}}
  \begin{tabular}{@{}l@{}}
    $du - \dfrac{u· dv}{v}$ \\
    $du + \dfrac{du· dv}{v}\DStrut$ \\\hline
  \end{tabular}\\
& \phantom{\TmpColA \;u + du}
  \begin{tabular}{@{}l@{}}
    ${} - \dfrac{u· dv}{v} - \dfrac{du· dv}{v}\DStrut$ \\
    ${} - \dfrac{u· dv}{v} - \dfrac{u· dv· dv}{v^2}\DStrut$ \\\hline
  \end{tabular}\\
& \phantom{\TmpColA \;u + du - \dfrac{u· dv}{v}}
  {} - \dfrac{du· dv}{v} + \dfrac{u· dv· dv}{v^2}.
\end{align*}
これらの余項は二階の小量であり、無視することができます。さらに余りが小さくなるため、ここで除算を終了しても問題ありません。
So we have got:
\begin{align*}
y + dy &= \frac{u}{v} + \frac{du}{v} - \frac{u \cdot dv}{v^2}; \\
\intertext{which may be written}
%
&= \frac{u}{v} + \frac{v \cdot du - u \cdot dv}{v^2}. \\
\intertext{\indent Now subtract the original $y = \frac{u}{v}$, and we have left:}
%
dy &= \frac{v \cdot du - u \cdot dv}{v^2}; \\
%
\lintertext{whence}
\frac{dy}{dx} &= \frac{v \, \frac{du}{dx} - u \, \frac{dv}{dx}}{v^2}.
\end{align*}
This gives us our instructions as to \textit{how to differentiate
a quotient \emph{of two functions}. Multiply the
divisor function by the differential coefficient of
the dividend function; then multiply the dividend
function by the differential coefficient of the divisor
function; and subtract. Lastly divide by the square
of the divisor function}.
Going back to our example $y = \frac{bx^5 + c}{x^2 + a}$,
\begin{align*}
\lintertext{write}
bx^5 + c &= u; \\
\lintertext{and}
x^2  + a &= v.
\end{align*}
Then
\begin{align*}
\frac{dy}{dx} &= \frac{(x^2 + a) \, \frac{d(bx^5 + c)}{dx} - (bx^5 + c) \, \frac{d(x^2 + a)}{dx}}{(x^2 + a)^2} \\
&= \frac{(x^2 + a)(5bx^4) - (bx^5 + c)(2x)}{(x^2 + a)^2}, \\
\frac{dy}{dx} &= \frac{3bx^6 + 5abx^4 - 2cx}{(x^2 + a)^2}. \quad (\textit{Answer}.)
\end{align*}
The working out of quotients is often tedious, but
there is nothing difficult about it.
Some further examples fully worked out are given
hereafter. (1) Differentiate $y = \frac{a}{b^2} x^3 - \frac{a^2}{b} x + \frac{a^2}{b^2}$.
Being a constant, $\frac{a^2}{b^2}$ vanishes,
and we have
\[
\frac{dy}{dx} = \frac{a}{b^2} \times 3 \times x^{3-1} - \frac{a^2}{b} \times 1 \times x^{1-1}.
\]

But $x^{1-1} = x^0 = 1$; so we get:
\[
\frac{dy}{dx} = \frac{3a}{b^2} x^2 - \frac{a^2}{b}.
\]

(2) Differentiate $y = 2a\sqrt{bx^3} - \frac{3b \sqrt[3]{a}}{x} - 2\sqrt{ab}$.

Putting $x$ in the index form, we get
\[
y = 2a\sqrt{b} x^{\frac{3}{2}} - 3b \sqrt[3]{a} x^{-1} - 2\sqrt{ab}.
\]

Now
\begin{gather*}
\frac{dy}{dx} = 2a\sqrt{b} \times \tfrac{3}{2} \times x^{\frac{3}{2}-1} - 3b\sqrt[3]{a} \times (-1) \times x^{-1-1}; \\
\lintertext{or,}
\frac{dy}{dx} = 3a\sqrt{bx} + \frac{3b\sqrt[3]{a}}{x^2}.
\end{gather*}

(3) Differentiate $z = 1.8 \sqrt[3]{\frac{1}{\theta^2}} - \frac{4.4}{\sqrt[5]{\theta}} - 27^\circ$.

This may be written: $z= 1.8\, \theta^{-\frac{2}{3}} - 4.4\, \theta^{-\frac{1}{5}} - 27^\circ$.

The 27$^\circ$ vanishes, and we have
\begin{gather*}
\frac{dz}{d\theta}
  = 1.8 \times  -\tfrac{2}{3} \times  \theta^{-\frac{2}{3}-1}
  - 4.4 \times  \left(-\tfrac{1}{5}\right)\theta^{-\frac{1}{5}-1}; \\
\text{or,}
\frac{dz}{d\theta}
  = -1.2\, \theta^{-\frac{5}{3}} + 0.88\, \theta^{-\frac{6}{5}}; \\
\text{or,}
\frac{dz}{d\theta} = \frac{0.88}{\sqrt[5]{\theta^6}}
  - \frac{1.2}{\sqrt[3]{\theta^5}}.
\end{gather*}

(4) $v = (3t^2 - 1.2 t + 1)^3$ の微分を計算します。

直接的な方法は後で説明します（参照：\Pageref{dodge}）；しかし、現在では難なく計算できます。

立方を展開すると
\[
v = 27t^6 - 32.4t^5 + 39.96t^4 - 23.328t^3 + 13.32t^2 - 3.6t + 1;
\]
したがって
\[
\frac{dv}{dt} = 162t^5 - 162t^4 + 159.84t^3 - 69.984t^2 + 26.64t - 3.6.
\]

(5) $y = (2x - 3)(x + 1)^2$ の微分を計算します。
\begin{alignat*}{2}
\frac{dy}{dx}
  &= (2x - 3)\, \frac{d\bigl[(x + 1)(x + 1)\bigr]}{dx}
     &&+ (x + 1)^2\, \frac{d(2x - 3)}{dx} \\
  &= (2x - 3) \left[(x + 1)\, \frac{d(x + 1)}{dx}\right.
     &&+ \left.(x + 1)\, \frac{d(x + 1)}{dx}\right] \\
  &  &&+ (x + 1)^2\, \frac{d(2x - 3)}{dx} \\
% [** TN: No natural second alignment point in next line; use fakery]
  &= \rlap{$2(x + 1)\bigl[(2x - 3) + (x + 1)\bigr] = 2(x + 1)(3x - 2)$;}&&
\end{alignat*}
または、より単純に、展開してから微分します。

(6) $y = 0.5 x^3(x-3)$ の微分を計算します。
\begin{align*}
\frac{dy}{dx}
  &= 0.5\left[x^3 \frac{d(x-3)}{dx} + (x-3) \frac{d(x^3)}{dx}\right] \\
  &= 0.5\left[x^3 + (x-3) \times  3x^2\right] = 2x^3 - 4.5x^2.
\end{align*}

前例と同様の注意が必要です。

(7) $w = \left(\theta + \dfrac{1}{\theta}\right)
   \left(\sqrt{\theta} + \dfrac{1}{\sqrt{\theta}}\right)$ の微分を計算します。

$$w = (\theta + \theta^{-1})(\theta^{\frac{1}{2}} + \theta^{-\frac{1}{2}}).$$

\begin{aligned}
\dfrac{dw}{d\theta}
  &= (\theta + \theta^{-1})
     \dfrac{d(\theta^{\frac{1}{2}} + \theta^{-\frac{1}{2}})}{d\theta}
   + (\theta^{\frac{1}{2}} + \theta^{-\frac{1}{2}})
     \dfrac{d(\theta+\theta^{-1})}{d\theta} \\
%
  &= (\theta + \theta^{-1})(\tfrac{1}{2}\theta^{-\frac{1}{2}}
                          - \tfrac{1}{2}\theta^{-\frac{3}{2}})
   + (\theta^{\frac{1}{2}} + \theta^{-\frac{1}{2}})(1 - \theta^{-2}) \\
%
  &= \tfrac{1}{2}(\theta^{ \frac{1}{2}} + \theta^{-\frac{3}{2}}
                - \theta^{-\frac{1}{2}} - \theta^{-\frac{5}{2}})
   + (\theta^{ \frac{1}{2}} + \theta^{-\frac{1}{2}}
    - \theta^{-\frac{3}{2}} - \theta^{-\frac{5}{2}}) \\%
%
  &= \tfrac{3}{2} \left(\sqrt{\theta} - \frac{1}{\sqrt{\theta^5}}\right)
   + \tfrac{1}{2} \left(\frac{1}{\sqrt{\theta}} - \frac{1}{\sqrt{\theta^3}}\right).
\end{aligned}

また、この結果は最初に二つの因子を乗算し、その後微分することでより簡単に得られます。しかし、これは常に可能ではありません；例えば、例8（\Pageref{example1}）では積の微分のルールが必須となります。

(8) $y = \dfrac{a}{1 + a\sqrt{x} + a^2x}$ の微分を計算します。
\begin{align*}
\dfrac{dy}{dx}
  &= \frac{(1 + ax^{\frac{1}{2}} + a^2x) \times  0 - a\dfrac{d(1 + ax^{\frac{1}{2}} + a^2x)}{dx}}
          {(1 + a\sqrt{x} + a^2x)^2} \\
  &= - \frac{a(\frac{1}{2}ax^{-\frac{1}{2}} + a^2)}
            {(1 + ax^{\frac{1}{2}} + a^2x)^2}.
\end{align*}

(9) $y = \dfrac{x^2}{x^2 + 1}$ の微分を計算します。
\[
\dfrac{dy}{dx} = \dfrac{(x^2 + 1)\, 2x - x^2 \times  2x}{(x^2 + 1)^2} = \dfrac{2x}{(x^2 + 1)^2}.
\]

(10) $y = \dfrac{a + \sqrt{x}}{a - \sqrt{x}}$ の微分を計算します。

インデックス形式では、$y = \dfrac{a + x^{\frac{1}{2}}}{a - x^{\frac{1}{2}}}$ です。
\begin{gather*}
\dfrac{dy}{dx}
  = \frac{(a - x^{\frac{1}{2}})( \tfrac{1}{2} x^{-\frac{1}{2}})
        - (a + x^{\frac{1}{2}})(-\tfrac{1}{2} x^{-\frac{1}{2}})}
         {(a - x^{\frac{1}{2}})^2}
  = \frac{ a - x^{\frac{1}{2}}
         + a + x^{\frac{1}{2}}}
        {2(a - x^{\frac{1}{2}})^2\, x^{\frac{1}{2}}}； \\
\lintertext{\rlap{従って}}
\dfrac{dy}{dx} = \frac{a}{(a - \sqrt{x})^2\, \sqrt{x}}.
\end{gather*}

\begin{align*}
\theta &= \frac{1 - a t^{\frac{2}{3}}}{1 + a t^{\frac{3}{2}}} \\
\frac{d\theta}{dt} 
  &= \frac{(1 + a t^{\frac{3}{2}}) \left(-\frac{2}{3} a t^{-\frac{1}{3}}\right)
         - (1 - a t^{\frac{2}{3}}) \cdot \frac{3}{2} a t^{\frac{1}{2}}}
          {(1 + a t^{\frac{3}{2}})^2} \\
  &= \frac{5a^2 t^{\frac{7}{6}} - \frac{4a}{\sqrt[3]{t}} - 9a \sqrt{t}}
          {6(1 + a \sqrt{t^3})^2}.
\end{align*}

（12）底面が垂直方向に45$^\circ$の傾斜角度を持つ正方形断面の貯水槽がある。底面の辺は200フィートである。水深が1フィート変化するときの流入または流出量を求めよ；その後、水深が24時間で14フィートから10フィートに減少したときに、毎時間どのくらいの量が引き出されたか求めよ。

高さが$H$、底面が$A$、上部が$a$である台形棱柱の体積は$V = \frac{H}{3} (A + a + \sqrt{Aa})$である。傾斜が45$^\circ$である場合、水深が$h$であるとき、水面の正方形の辺の長さは$200 + 2h$フィートであり、したがって水の体積は
\[
V = \frac{h}{3} [200^2 + (200 + 2h)^2 + 200(200 + 2h)]
= 40,000h + 400h^2 + \frac{4h^3}{3}
\]
である。

$\frac{dV}{dh} = 40,000 + 800h + 4h^2$ は深さ1フィートの変化あたりの立方フィート数である。水深が14フィートから10フィートに減少する平均水位は12フィートであり、$h = 12$ のとき $\frac{dV}{dh} = 50,176$ 立方フィート。

24時間で水深が4フィート変化した場合の每時ガロン数は$\frac{4 \times 50,176 \times 6.25}{24} = 52,267$ガロンである。

（13）Dulongによって与えられた、温度$t^\circ\mathrm{C}$で飽和した蒸気の絶対圧力$P$は、$t$が80℃以上である限り $P = \left( \frac{40 + t}{140} \right)^5$ とされる。100℃での圧力の温度に対する変化率を求めよ。

展開の分子は二項定理により計算します（参照：\Pageref{binomtheo}）。
\[
P = \frac{1}{140^5} (40^5 + 5\times 40^4 t + 10 \times  40^3 t^2 + 10 \times  40^2 t^3
                            + 5 \times  40t^4 + t^5);
\]
%[** TN: \BindMath gives bad line spacing after DPalign]
\begin{align*}
\lintertext{したがって} \frac{dP}{dt} = &\frac{1}{537,824 \times  10^5}\\
      &(5 \times  40^4 + 20 \times  40^3 t + 30 \times  40^2 t^2 + 20 \times  40t^3 + 5t^4),
\end{align*}
$t = 100$ のとき、これは $0.036$ アトモスフィアを 1 センチグラードの温度変化に対応する値になります。

第III章の問題（解答は \Pageref{AnsEx:III} を参照）。

\begin{enumerate}
\item[(1)] $u = 1 + x + \frac{x^2}{1 \times  2} + \frac{x^3}{1 \times  2 \times  3} + \dotsb$ の微分を計算します（例題参照：Pagelabel{examples2}）。

\item[(\textit{a})] $y = ax^2 + bx + c$. \hfil (\textit{c})\textasciitilde{}$y = (x + a)^2$.

\item[(\textit{d})] $y = (x + a)^3$.

\item[(2)] $w = at - \frac{1}{2}bt^2$ のとき、$\frac{dw}{dt}$ を求めます。

\item[(3)] $y = (x + \sqrt{-1}) \times  (x - \sqrt{-1})$ の微分係数を求めます。

\item[(4)] $y = (197x - 34x^2) \times  (7 + 22x - 83x^3)$ の微分を求めます。

\item[(5)] $x = (y + 3) \times  (y + 5)$ のとき、$\frac{dx}{dy}$ を求めます。

\item[(6)] $y = 1.3709x \times  (112.6 + 45.202x^2)$ の微分を求めます。

次のように微分係数を求めます：

\item[(7)] $y = \frac{2x + 3}{3x + 2}$。

\item[(8)] $y = \frac{1 + x + 2x^2 + 3x^3}{1 + x + 2x^2}$。

\item[(9)] $y = \frac{ax + b}{cx + d}$。

\item[(10)] $y = \frac{x^n + a}{x^{-n} + b}$。

\item[(11)] 発光電球のフィラメントの温度 $t$ は、その電流と関連付けられており、関係式は
\[
C = a + bt + ct^2.
\]

温度が変化したときに対応する電流の変動を表す式を求めます。

\item[(12)] 次の公式は $0^\circ$ センチグラードでの同一の導線の抵抗 $R_0$ に対する温度 $t^\circ$\;C. での電気抵抗 $R$ の関係を表現するために提案されています。$a$, $b$, $c$ は定数です。
\begin{align*}
R &= R_0(1 + at + bt^2). \\
R &= R_0(1 + at + b\sqrt{t}). \\
R &= R_0(1 + at + bt^2)^{-1}.
\end{align*}

これらの公式により、抵抗が温度にどのように変化するかの率を求めます。

\item[(13)] ある種類の標準電池の電动势$E$は、温度$t$に従って関係式
\[
E = 1.4340 \bigl[1 - 0.000814(t-15)
                 + 0.000007(t-15)^2\bigr] \text{ ボルト}
\]
で変化することが発見されました。

$15^\circ$, $20^\circ$, および $25^\circ$における電动势の変化量を求めよ。

\item[(14)] 長さ$l$、電流強度$i$を持つ電弧を維持するための電动势は、Mrs.\textasciitilde{}Ayrtonによって
\[
E = a + bl + \frac{c + kl}{i}
\]
で求められました。ここで$a$, $b$, $c$, $k$は定数です。

(\textit{a}) 電弧の长さに対する電动势の変化を、(\textit{b}) 電流强度に対する電动势の変化を求めよ。

\end{enumerate}



\chapter{逐次微分法（高階微分）}
\noindent\textit{\large Successive Differentiation}\\[
2mm]
\noindent\textbf{\small 【章の論理とテーマ】2階微分・3階微分の意味と加速度の物理的導出}\\[4mm]
第7章 微分の連続的適用
まず、関数を複数回微分する操作の効果を試してみましょう（参照：定義 \pageref{function}）。具体的な例から始めます。
$y = x^5$ とする。
\begin{alignat*}{3}
&\text{第一次微分, } &&5x^4.              &&         \\
&\text{第二次微分, } &&5 \times  4x^3           &&= 20x^3. \\
&\text{第三次微分, } &&5 \times  4 \times  3x^2       &&= 60x^2. \\
&\text{第四次微分, } &&5 \times  4 \times  3 \times  2x     &&= 120x.  \\
&\text{第五次微分, } &&5 \times  4 \times  3 \times  2 \times  1  &&= 120.   \\
&\text{第六次微分, } &&                   &&= 0.
\end{alignat*}
ある記法があります（参照：定義 \pageref{notation}）で、一部の著者はこの記法を使用しています。これは非常に便利です。それは任意の $x$ の関数を表す一般的な符号 $f(x)$ を用いることです。ここで、符号 $f(~)$ は「$x$ の関数」と読み取り、具体的にどの関数であるかは言及しません。したがって、$y = f(x)$ という声明は $y$ が $x$ の関数であることを示すだけであり、それは $x^2$ や $ax^n$、または $\cos x$ など、他の複雑な $x$ の関数のいずれかである可能性があります。
微分係数に対応する符号は $f'(x)$ で、$\dfrac{dy}{dx}$ よりも書きやすいです。これは「導関数」と呼ばれます。
再度微分すると、「第二次導関数」または「第二微分係数」を得ることができます。これは $f''(x)$ で表されます。同様に、さらに微分すると「第三次導関数」や「第四次導関数」が得られます。
今度は一般化します。
$y = f(x) = x^n$ とする。
\begin{align*}[m]
&\text{第一次微分,} &&f'(x) &= nx^{n-1}. \\
&\text{第二次微分,} &&f''(x) &= n(n-1)x^{n-2}. \\
&\text{第三次微分,} &&f'''(x) &= n(n-1)(n-2)x^{n-3}. \\
&\text{第四次微分,} &&f''''(x) &= n(n-1)(n-2)(n-3)x^{n-4}. \\
&&\llap{\text{等々。}} \text{ 等々。}
\end{align*}
しかし、これは連続した微分を示す唯一の方法ではありません。例えば，
\begin{align*}
\text{元の関数が} \quad y &= f(x);  \\
\text{1回微分すると} \quad \frac{dy}{dx} &= f'(x); \\
\text{2回微分すると} \quad \frac{d\left(\frac{dy}{dx}\right)}{dx} &= f''(x);
\end{align*}
これはより便利に $\frac{d^2y}{(dx)^2}$ または通常 $\frac{d^2y}{dx^2}$ の形で書くことができます。同様に、3回微分すると $\frac{d^3y}{dx^3} = f'''(x)$ となります。
\begin{align*}
y = f(x) = 7x^4 + 3.5x^3 - \frac{1}{2}x^2 + x - 2 の場合、 \\
\frac{dy}{dx} &= f'(x) = 28x^3 + 10.5x^2 - x + 1, \\
\frac{d^2y}{dx^2} &= f''(x) = 84x^2 + 21x - 1, \\
\frac{d^3y}{dx^3} &= f'''(x) = 168x + 21, \\
\frac{d^4y}{dx^4} &= f''''(x) = 168, \\
\frac{d^5y}{dx^5} &= f'''''(x) = 0.
\end{align*}
同様に、$y = \phi(x) = 3x(x^2 - 4)$ の場合、
\begin{align*}
\phi'(x) &= \frac{dy}{dx} = 3[x \times  2x + (x^2 - 4) \times  1] = 3(3x^2 - 4), \\
\phi''(x) &= \frac{d^2y}{dx^2} = 3 \times  6x = 18x, \\
\phi'''(x) &= \frac{d^3y}{dx^3} = 18, \\
\phi''''(x) &= \frac{d^4y}{dx^4} = 0.
\end{align*}
演習 IV（解答は Pageref[page]{AnsEx:III} を参照してください）。
次の式について $\frac{dy}{dx}$ と $\frac{d^2y}{dx^2}$ を求めよ：
\begin{enumerate}
\item (1) $y = 17x + 12x^2$.
\item (2) $y = \frac{x^2 + a}{x + a}$.
さらに、次の式について2番目の導関数と3番目の導関数を求めよ：
\item (3) $y = 1 + \frac{x}{1} + \frac{x^2}{1\times 2} + \frac{x^3}{1\times 2\times 3} + \frac{x^4}{1\times 2\times 3\times 4}$.
\item (4) 演習 III（Pageref{examples2}）の No. 1 から No. 7、および例題（Pageref{examples3}）の No. 1 から No. 7 の2番目の導関数と3番目の導関数を求めよ。
\end{enumerate}
\chapter{時間が変化するとき（時間変化率）}
\noindent\textit{\large When Time Varies}\\[2mm]
\noindent\textbf{\small 【章の論理とテーマ】速度・加速度・物理現象における時間の導関数と連鎖律の直観}\\[4mm]
第八章 時間が変化するとき
重要な微分積分学の問題のいくつかは、時間が独立した変数であり、その他の量が時間と共に変化するときにその値を考える問題である。あるものは時間が経つにつれて大きくなり、別のものは小さくなる。列車が出発地点からどれだけ離れているかという距離は、時間が経つにつれて永遠に増加する。年月を重ねて木は高くなる。12インチの植物が1ヶ月で14インチになるのに対し、12フィートの木が1年で14フィートになるのではどちらの方が速い成長をしているのか？
この章では、「率」という言葉に多く頼ることにする。これは貧困率や水道料金とは関係ない（ただし、ここでも「率」は比を想起させる）。出生率や死亡率も関係ないが、これらの単語は千人の人口あたりの出産数や死者数を思い起こす。自動車が私たちの眼前を駆け抜けるとき、「何と驚異的な速度だ！」と言ったとき、あるいは浪費家が金を使い果たすとき、「その若者は驚異的な速度で生活している」と言ったとき、私たちは何を意味しているのか？
「率」について何を言っているのか。これらの例では、何かが起こる過程とそれがどれだけの時間を要するかという比較をしている。自動車が秒速10ヤードで私達の眼前を駆け抜ける場合、簡単な計算からこれは1分間に600ヤード、または20マイルを超える速度に相当することになる。
では、なぜ秒速10ヤードは一分間に600ヤードと同じ速度であると言えるのか？10ヤードと600ヤードは同じものではなく、一秒と一分も同じものではない。私たちが「率」が同じであると言っている意味は、距離を通過した割合とその時間を通過するのに要した時間の比が両方の場合で同じであるということだ。
別の例を挙げよう。ある男が所持金が僅かな金額であるとしても、年間数千万ポンドを費やすことができる——ただし、そのような費やし方が数分間だけ続く場合に限る。例えば、あなたがカウンターで何らかの商品を支払いのために1シリング渡すと仮定し、その操作が正確に1秒間にわたることを想像する。すると、その短い操作期間においては、あなたは1秒間に1シリングずつ金銭を手放しており、これは1分間に3ポンド、1時間に180ポンド、1日に4320ポンド、1年間で1576800ポンド相当の速度である！あなたのポケットに10ポンドが入っているなら、そのような速度で金銭を費やすことができるのはわずか5分25秒程度である。
「バン！」とシリングが消えたという話では、サンディはロンドンに入りたてだったにもかかわらず、わずか5分未満でその音がしたと言われている。もし彼が一日中そのような速度で金銭を費やすとした場合、つまり12時間間続けていたとすると、彼は1時間に6シリング、または3ポンド12シリングずつ費やし、1日に21ポンド12シリングずつ費やすことになる。土曜日を除く。
これらのアイデアを微分記号で表してみよう。
この場合、$y$ を金銭、$t$ を時間を表すことにしよう。
もし金銭を費やしているとすると、短時間 $dt$ に費やす金額を $dy$ とする。その費やし速度は $\dfrac{dy}{dt}$ であるが、$dy$ は減少量であり増加量ではないため、マイナス符号をつけた形で書き表すべきであり、つまり $-\dfrac{dy}{dt}$ となるべきである。しかし金銭は微積分の例としては良いものではない。それは通常一気に来たり去ったりするからだ——1年間に200ポンドを稼ぐかもしれないが、それがあなたの前で細い流れのように常に流れてくるわけではない。それは週に一度、月に一度、または四半期ごとにまとまってやって来るし、また支出も一気に支払われる。
（訳者注：文中の「sawbbath」は「土曜日」を意味するイギリス英語で、「Sabbath」の誤植と思われる。）
より適切な速度の概念を示す例は、移動物体の速さです。ロンドン（ユーストン駅）からリバプールまで $200$ マイルあります。列車が午後 $7$ 時にロンドンを出発し、午後 $11$ 時にリバプールに到着するとします。この場合、$4$ 時間で $200$ マイルを移動したので、平均速度は $50$ マイル每時であることになります；なぜなら $\frac{200}{4} = \frac{50}{1}$ です。実際には、距離を通過する時間に対する精神的な比較を行っています。一方を他方で割っているのです。もし $y$ を全距離、$t$ を全時間を表すとするならば、平均速度は $\frac{y}{t}$ となります。しかし、この移動の間、速度は常に一定ではありません：出発時と旅程終了時の減速では速度が低くなります。おそらく下り坂で走っている途中には、速度が $60$ マイル每時を超えるかもしれません。ある特定の時間要素 $dt$ において通過した距離要素を $dy$ とするならば、その部分での速度は $\frac{dy}{dt}$ となります。一方量（ここでは「距離」）が他方量（この場合は「時間」）に対してどのように変化しているかの速さは、それらの量の微分係数で適切に表現されます。科学的に正確な速度は、特定の方向における非常に小さな距離がどの程度の速さで通過されているかを表すものであり、従って次のように書くことができます：
\[
v = \frac{dy}{dt}.
\]

しかし、速度 $v$ が一様でない場合、それは増加しているか減少している必要があります。速度がどのように増加するかの速さは「加速度」と呼ばれます。ある物体が特定の瞬間において時間要素 $dt$ の間に追加された速度 $dv$ を得ているとします。その瞬間での加速度 $a$ は次のように書けます：
\[
a = \frac{dv}{dt};
\]
ここで $dv$ は $\frac{dy}{dt}$ の微分であるため、
\[
a = \frac{d\left( \frac{dy}{dt} \right)}{dt};
\]
これは通常 $a = \frac{d^2y}{dt^2}$ と書き換えられます。つまり、加速度は時間に対する距離の二階微分係数です。加速度は単位時間における速度の変化量として表されます；たとえば、フィート每秒毎秒などと表現します。その記号は「フィート ÷ 秒$^2$」で表されます。

当列車が動き始めたとき、その速度$v$は小さいですが、急速に速さを増しています。エンジンの努力により加速されるため、$\dfrac{d^2y}{dt^2}$ の値は大きいです。最高速度に達すると、もう加速されないため、そのとき $\dfrac{d^2y}{dt^2}$ は0になります。しかし、停止地点に近づくと速度が徐々に減速し始め、ブレーキをかけた場合などには急速に減速します。この減速期において、速度の変化率$\dfrac{dv}{dt}$、つまり $\dfrac{d^2y}{dt^2}$ の値は負になります。

質量$m$を加速するためには継続的な力が必要です。必要な力を加速する質量に比例し、その加速度にも比例します。したがって、力$f$は以下の式で表すことができます。
\begin{align*}
f &= ma;\\
\text{または} \quad f &= m \frac{dv}{dt}; \\
\text{または} \quad f &= m \frac{d^2y}{dt^2}.
\end{align*}

質量$m$とその速度$v$の積は、その運動量と呼ばれ、記号で $mv$ と表されます。運動量を時間について微分すると、運動量の変化率$\dfrac{d(mv)}{dt}$が得られます。ここで、$m$は定数であるため、この式は$m \dfrac{dv}{dt}$に書き換えることができ、これは先ほど述べた$f$と同じです。つまり、力は質量$m$と加速度$a$の積、または運動量の変化率として表すことができます。

また、反対方向に等しい抵抗力に対して力を用いて何かを動かす場合、その力は\emph{作業}を行います。そして、行われた作業の量は、その点が移動する距離（その力の作用方向において）と力の積で測定されます。したがって、力$f$が長さ$y$を移動すると、行われた作業（これを$w$と呼びましょう）は
\[
w = f \times y;
\]
となります。ここで、$f$を恒常的な力をとります。もし力が範囲$y$の各部分で異なる場合、その値を求めなければなりません。長さ$dy$の小さな要素における力$f$があるとき、行われた作業量は$f \times dy$です。しかし、$dy$は長さの要素であるため、作業量も要素となります。ここで$w$を作業とし、要素の作業量を$dw$とすると
\begin{align*}
dw &= f \times dy; \\
\intertext{これは次のように書くこともできます}
dw &= ma \cdot dy; \\
\lintertext{または}
dw &= m \frac{d^2y}{dt^2} \cdot dy; \\
\lintertext{または}
dw &= m \frac{dv}{dt} \cdot dy.     \\
\intertext{\indent さらに、式を移項して書き換えると}
\frac{dw}{dy} &= f.
\end{align*}

これはまた力の定義を与えます。何かを任意の方向に移動させるために力を用いている場合、その方向における力はその方向の長さあたりの作業量の比率または割合で行われている作業量の速さと等しいです。最後の文において、「速さ」という言葉は時間の意味ではなく、比や比例の意味で使われています。

アイザック・ニュートン博士は、レービッツと並んで微分積分学の方法を発明した人物ですが、変化するすべての量を「流れ」と見なしていました。そして、現代では導関数と呼ばれる比率は、その量の流速または「フリクション」であると考えていました。ニュートン博士は dy, dx および dt の記号を使用していません（これはレービッツによるものでした）が、独自の記法を持っていました。もし y が変化する量であれば、その変化率（または「フリクション」）を表す彼の記号は \dot{y} でした。もし x が変数であれば、そのフリクションは \dot{x} と呼ばれました。文字の上に点が付いていることは、その量が微分されたことを示していました。しかし、この記法では、どの独立変数に対する微分が行われたかを教えてくれません。dy/dt を見ると y が t に対して微分されていることがわかります。dy/dx を見ると y が x に対して微分されていることがわかります。しかし、単に \dot{y} を見ても、これは dy/dx, dy/dt, または dy/dz のいずれを意味するかを知るためには文脈を見なければなりません。したがって、このフリクション記法は導関数記法よりも情報量が少なく、結果的に使用から離れています。しかし、その簡潔さは、時間が独立変数である場合に限り、優位性があります。その場合、\dot{y} は dy/dt, \dot{u} は du/dt, \ddot{x} は d2x/dt2 を意味します。

このフリクション記法を採用して、上記の段落で考慮された機械学方程式を以下のように書くことができます：

\begin{tabular}{l@{\qquad\qquad}l}
距離 \& $x$, \\
速度 \& $v = \dot{x}$, \\
加速度 \& $a = \dot{v} = \ddot{x}$, \\
力 \& $f = m\dot{v} = m\ddot{x}$, \\
作業 \& $w = x \times  m \ddot{x}$.
\end{tabular}

例題。
（1）ある物体が一定点Oから距離x（フィート）を移動するとき、その関係はx = 0.2t$^2$ + 10.4と与えられます。ここでtは物体が特定の瞬間から経過した秒数です。5秒後における物体の速度と加速度を求めよ。また、距離が100フィートとなったときの値も求めよ。さらに、最初の10秒間の平均速度を求めよ。（右方向に移動した距離を正とする）

\begin{gather*}
x = 0.2t^2 + 10.4 \\
v = \dot{x} = \frac{dx}{dt} = 0.4t;\quad a = \ddot{x} = \frac{d^2x}{dt^2} = 0.4 = \text{定数。}
\end{gather*}

t = 0のとき、x = 10.4, v = 0である。物体は点Oから右に10.4フィート離れた位置から動き始め、時間はその瞬間から計測される。

t = 5のとき、v = 0.4 $\times$ 5 = 2 ft/sec; a = 0.4 ft/sec$^2$。

x = 100のとき、100 = 0.2t$^2$ + 10.4 より t$^2$ = 448, そして t = 21.17 sec.; v = 0.4 $\times$ 21.17 = 8.468 ft/sec。

t = 10のとき、
\begin{gather*}
移動距離 = 0.2 \times  10² + 10.4 - 10.4 = 20 ft \\
平均速度 = \frac{20}{10} = 2 ft/sec
\end{gather*}

（これは区間の中央、t = 5における速度と同じである。なぜなら、加速度が定数であるため、時間t = 0からt = 10までの間に速度は一様に変化しているからである。）

（2）上記の問題において、
\begin{gather*}
x = 0.2t^2 + 3t + 10.4 \\
v = \dot{x} = \frac{dx}{dt} = 0.4t + 3;\quad a = \ddot{x} = \frac{d^2x}{dt^2} = 0.4 = \text{定数。}
\end{gather*}

t = 0のとき、x = 10.4, v = 3 ft/sec。物体は点Oから10.4フィート離れた位置を通過した瞬間から動き始め、そのとき既に速度が3 ft/secであった。物体が動き始めた時間経過を求めるには、v = 0とするとよい。0.4t + 3 = 0 より t = -\frac{3}{0.4} = -7.5 sec。

物体は時間観測が始まる前から7.5秒動き続け、その後5秒経過すると t = -2.5, v = 0.4 $\times$ (-2.5) + 3 = 2 ft/sec。

x = 100 ftのとき、
\[
100 = 0.2t^2 + 3t + 10.4;\quad \text{または } t^2 + 15t - 448 = 0;
\]
よって t = 14.95 sec., v = 0.4 $\times$ 14.95 + 3 = 8.98 ft/sec。

To find the distance travelled during最初的10秒内的运动，必须知道物体从点O出发时与该点的距离。

$t = -7.5$のとき，
\[
x = 0.2 \times (-7.5)^2 - 3 \times 7.5 + 10.4 = -0.85~\text{ft}.,
\]
つまり，点Oの左側に0.85 ft.

$t = 2.5$のとき，
\[
x = 0.2 \times 2.5^2 + 3 \times 2.5 + 10.4 = 19.15.
\]

したがって，10秒間の移動距離は
$19.15 + 0.85 = 20$ ft.,
平均速度は
\[
\text{平均速度} = \frac{20}{10} = 2~\text{ft./sec}.
\]

(3) 同様の問題を考える．距離が$x = 0.2t^2 - 3t + 10.4$で与えられたとき，速度$v = 0.4t - 3$, 加速$a = 0.4 = \text{定数}$である．$t = 0$のとき，$x = 10.4$であり，かつ
\[
v = -3;
\]
つまり，物体は前の問題とは反対方向に動いている．しかし，加速が正であるため，この速度は時間とともに減少し，ゼロになるまで持続する．$v = 0$のとき，$0.4t - 3 = 0$; よって$t = 7.5$秒．その後，速度は正になり，物体が始まった後5秒後の$t = 12.5$のとき，
\[
v = 0.4 \times 12.5 - 3 = 2~\text{ft./sec}.
\]

$x = 100$のとき，
\begin{gather*}
100 = 0.2t^2 - 3t + 10.4,\quad \text{または } t^2 - 15t - 448 = 0, \\
\lintertext{そして}
t = 29.95;\ v = 0.4 \times 29.95 - 3 = 8.98~\text{ft./sec.}
\end{gather*}

$v$がゼロのとき，$x = 0.2 \times 7.5^2 - 3 \times 7.5 + 10.4 = -0.85$，
これは物体は点Oより0.85 ft.の左側に戻る前に停止することを示している．10秒後
\[
t = 17.5 \text{ と } x = 0.2 \times 17.5^2 - 3 \times 17.5 + 10.4 = 19.15.
\]
移動距離は$0.85 + 19.15 = 20.0$であり，平均速度もまた2 ft./sec.

(4) 同じ種類の別の問題を考える．$x = 0.2t^3 - 3t^2 + 10.4$; $v = 0.6t^2 - 6t$; $a = 1.2t - 6$. 加速は定数ではなくなり，$t = 0$のとき，$x = 10.4$, $v = 0$, $a = -6$. 物体は静止しており，しかし，負の加速を獲得し，点Oに向かって速度を得る準備ができている．

(5) $x = 0.2t^3 - 3t + 10.4$のとき，$v = 0.6t^2 - 3$, $a = 1.2t$. $t = 0$のとき，$x = 10.4$; $v = -3$; $a = 0$. 物体は点Oに向かって3 ft./sec.で動いており，その瞬間速度は一様である．

運動の条件は常に時間距離方程式とその最初および二つ目の導関数から即座に確認できます。後者の場合、最初の10秒間の平均速度と開始5秒後の速度は異なるため、速度が一様に増加していないこと、つまり加速度が定数でないことを考慮に入れる必要があります。

(6) 一定の瞬間からの秒数 \( t \) で与えられる円盤が回転する角度 \(\theta\)（ラジアン）は \(\theta = 3 + 2t - 0.1t^3\) で与えられます。角速度 \(\omega\) と角加速度 \(\alpha\) を求めよ。（a）1秒後；（b）1回転を完了したとき。どの時間に停止し、その時点までに何回転を行っているかを求めよ。

加速度は
\[
\omega = \dot{\theta} = \frac{d\theta}{dt} = 2 - 0.3t^2, \quad \alpha = \ddot{\theta} = \frac{d^2\theta}{dt^2} = -0.6t.
\]

\( t = 0 \) のとき、\(\theta = 3\)；\(\omega = 2\) ラジアン/秒；\(\alpha = 0\)。

\( t = 1 \) のとき
\[
\omega = 2 - 0.3 = 1.7 \text{ ラジアン/秒}; \quad \alpha = -0.6 \text{ ラジアン/秒}^2.
\]

これは減速であり、円盤は遅くなる。

1回転後
\[
\theta = 2\pi = 6.28; \quad 6.28 = 3 + 2t - 0.1t^3.
\]

グラフを描くと、\(\theta = 3 + 2t - 0.1t^3\) の値または値の \( t \) を得ることができます；これらは \( 2.11 \) および \( 3.03 \)（負の値も存在します）。

\( t = 2.11 \) のとき
\begin{gather*}
\theta = 6.28; \quad \omega = 2 - 1.34 = 0.66 \text{ ラジアン/秒}; \\
\alpha = -1.27 \text{ ラジアン/秒}^2.
\end{gather*}

\( t = 3.03 \) のとき
\[
\theta = 6.28; \quad \omega = 2 - 2.754 = -0.754 \text{ ラジアン/秒}; \\
\alpha = -1.82 \text{ ラジアン/秒}^2.
\]

速度が反転します。これらの瞬間の間に円盤は停止しており、\(\omega = 0\) のとき、すなわち \( 0 = 2 - 0.3t^3 \)、または \( t = 2.58 \text{ 秒} \) のときに停止しています。その時点までに
\[
\frac{\theta}{2\pi}
  = \frac{3 + 2 \times  2.58 - 0.1 \times  2.58^3}{6.28} = 1.025 \text{ 回}.
\]

問題集 V（解答は Pageref[page]{AnsEx:V} を参照）。

問題
(1) \( y = a + bt^2 + ct^4 \) のとき、\(\frac{dy}{dt}\) と \(\frac{d^2y}{dt^2}\) を求めよ。
\[
\text{\textit{答} } \frac{dy}{dt} = 2bt + 4ct^3; \quad \frac{d^2y}{dt^2} = 2b + 12ct^2.
\]

\begin{enumerate}
\item{(2)} 自由落下する物体が$t$秒後に移動する空間$s$（フィート）は、方程式によって表される：
\[
 s = 16t^2 
\]
を用いて曲線を描き、物体の速度を求めよ。以下から：$t = 2$秒；$t = 4.6$秒；$t = 0.01$秒。

\item{(3)} $x = at - \frac{1}{2}gt^2$；$\dot{x}$と$\ddot{x}$を求めよ。

\item{(4)} 次の法則に従って物体が動くとき、$t = 4$秒における速度を求めよ。$s$はフィートである：
\[
 s = 12 - 4.5t + 6.2t^2 
\]

\item{(5)} 前問で述べた物体の加速度を求めよ。すべての$t$の値に対して加速度が同じかを確認せよ。

\item{(6)} 回転する輪が時間$t$（秒）経過したとき、角度$\theta$（ラジアン）は以下の法則で与えられる：
\[
 \theta = 2.1 - 3.2t + 4.8t^2 
\]

$1\frac{1}{2}$秒後の角速度（ラジアン/秒）を求めよ。また、角加速度を求めよ。

\item{(7)} 滑り棒が動くとき、最初の部分での距離$s$（インチ）は以下の式で与えられる：
\[
 s = 6.8t^3 - 10.8t \quad (t \text{は秒}) 
\]

任意の時間における速度と加速度を求めよ。また、$3$秒後の速度と加速度を求めよ。

\item{(8)} 水中への上昇する balloon の運動は、その瞬間での高さ$h$（マイル）が以下の式で与えられる：
\[
 h = 0.5 + \frac{1}{10}\sqrt[3]{t-125} \quad (t \text{は秒}) 
\]

任意の時間における速度と加速度を求めよ。最初の十分間に高さ、速度、加速度の変化を示す曲線を描け。

\item{(9)} 水面に達した後$t$秒後の深さ$p$（メートル）は以下の式で与えられる：
\[
 p = \frac{4}{4+t^2} + 0.8t - 1 
\]

任意の時間における速度と加速度を求めよ。$10$秒後の速度と加速度を求めよ。

ある物体は、時間$t$から開始して移動した距離$s$が$t^n$で与えられるように動くとする。ここで$n$は定数である。5秒目から10秒目の間に速度が2倍になるときの$n$の値を求めよ；また、10秒末における速度と加速度の絶対値が等しいときの$n$の値も求めよ。
\end{enumerate}



\chapter{便利な巧妙法（合成関数の微分法・連鎖律）}
\noindent\textit{\large Introducing a Useful Dodge}\\[
2mm]
\noindent\textbf{\small 【章の論理とテーマ】置換変数を用いた dy/dx = (dy/du)·(du/dx) の自在な計算技法}\\[4mm]
第九章 有用な技巧の導入
しばしば、微分すべき式が複雑で直接的に扱いにくい場合がある。
例えば、方程式（）
\[
y = (x^2 + a^2)^{\frac{3}{2}}
\]
は初心者にとっては難しい。

この難題を克服する方法は次の通り：$x^2 + a^2$ の表現を $u$ などのシンボルで表す。すると方程式は
\[
y = u^{\frac{3}{2}}
\]
となり、簡単に扱えるようになる。つまり
\[
\frac{dy}{du} = \frac{3}{2} u^{\frac{1}{2}}.
\]
次に $u = x^2 + a^2$ の式を扱い、$x$ に関する微分を行う。
\[
\frac{du}{dx} = 2x.
\]

それでは残りは平らげばよい；
\begin{align*}
\lintertext{つまり}
\frac{dy}{dx} &= \frac{dy}{du} \times  \frac{du}{dx}; \\
\lintertext{すなわち}
\frac{dy}{dx}
  &= \frac{3}{2} u^{\frac{1}{2}} \times  2x \\
  &= \tfrac{3}{2} (x^2 + a^2)^{\frac{1}{2}} \times  2x \\
  &= 3x(x^2 + a^2)^{\frac{1}{2}};
\end{align*}
そしてこの技巧は完了する。

いずれにせよ、正弦、余弦、指数関数などを扱う方法を学んだ後には、この技巧がますます有用になるだろう。

例題。
いくつかの例を使ってこの技巧を練習しよう。

%
(1) $y = \sqrt{a+x}$ を微分せよ。

$u = a + x$ とおく。
\begin{align*}
\frac{du}{dx} &= 1;\quad y=u^{\frac{1}{2}};\quad
  \frac{dy}{du} = \tfrac{1}{2} u^{-\frac{1}{2}}
                = \tfrac{1}{2} (a+x)^{-\frac{1}{2}}.\\
%
\frac{dy}{dx} &= \frac{dy}{du} \times  \frac{du}{dx} = \frac{1}{2\sqrt{a+x}}.
\end{align*}

(2) $y = \dfrac{1}{\sqrt{a+x^2}}$ を微分せよ。

$u = a + x^2$ とおく。
\begin{align*}
\frac{du}{dx} &= 2x;\quad y=u^{-\frac{1}{2}};\quad
  \frac{dy}{du} = -\tfrac{1}{2}u^{-\frac{3}{2}}.\\
%
\frac{dy}{dx} &= \frac{dy}{du}\times \frac{du}{dx} = - \frac{x}{\sqrt{(a+x^2)^3}}.
\end{align*}

(3) $y = \left(m - nx^{\frac{2}{3}} + \dfrac{p}{x^{\frac{4}{3}}}\right)^a$ を微分せよ。

$u = m - nx^{\frac{2}{3}} + px^{-\frac{4}{3}}$ とおく。
\begin{gather*}
\frac{du}{dx} = -\tfrac{2}{3} nx^{-\frac{1}{3}} - \tfrac{4}{3} px^{-\frac{7}{3}};\\
%
y = u^a;\quad \frac{dy}{du} = a u^{a-1}. \\
%
\frac{dy}{dx} = \frac{dy}{du}\times \frac{du}{dx}
  = -a\left(m -nx^{\frac{2}{3}} + \frac{p}{x^{\frac{4}{3}}}\right)^{a-1}
     (\tfrac{2}{3} nx^{-\frac{1}{3}} + \tfrac{4}{3} px^{-\frac{7}{3}}).
\end{gather*}

(4) $y=\dfrac{1}{\sqrt{x^3 - a^2}}$ を微分せよ。

Let \( u = x^3 - a^2 \).
\begin{align*}
\frac{du}{dx} &= 3x^2; \quad y = u^{-\frac{1}{2}}; \quad
  \frac{dy}{du} = -\frac{1}{2}(x^3 - a^2)^{-\frac{3}{2}}. \\
\frac{dy}{dx} &= \frac{dy}{du} \times \frac{du}{dx} = -\frac{3x^2}{2\sqrt{(x^3 - a^2)^3}}.
\end{align*}

(5) \( y = \sqrt{\dfrac{1-x}{1+x}} \).

これを \( y = \dfrac{(1-x)^{\frac{1}{2}}}{(1+x)^{\frac{1}{2}}} \) と書き換える。
\[
\frac{dy}{dx}
  = \frac{(1+x)^{\frac{1}{2}} \, \dfrac{d(1-x)^{\frac{1}{2}}}{dx}
        - (1-x)^{\frac{1}{2}} \, \dfrac{d(1+x)^{\frac{1}{2}}}{dx}}{1+x}.
\]

（または \( y = (1-x)^{\frac{1}{2}} (1+x)^{-\frac{1}{2}} \) として積の形で微分してもよい。）

上記の例（1）と同様に計算すると，
\[
\dfrac{d(1-x)^{\frac{1}{2}}}{dx} = -\frac{1}{2\sqrt{1-x}};
\quad \text{および} \quad
\dfrac{d(1+x)^{\frac{1}{2}}}{dx} = \frac{1}{2\sqrt{1+x}}.
\]

したがって，
\begin{align*}
\frac{dy}{dx}
  &= - \frac{(1 + x)^{\frac{1}{2}}}{2(1 + x)\sqrt{1-x}}
     - \frac{(1 - x)^{\frac{1}{2}}}{2(1 + x)\sqrt{1+x}} \\
  &= - \frac{1}{2\sqrt{1+x}\sqrt{1-x}} - \frac{\sqrt{1-x}}{2 \sqrt{(1+x)^3}};\\
\lintertext{または}
\frac{dy}{dx}
  &= - \frac{1}{(1+x)\sqrt{1-x^2}}.
\end{align*}

(6) \( y = \sqrt{\dfrac{x^3}{1+x^2}} \).

これを
\begin{gather*}
y = x^{\frac{3}{2}} (1+x^2)^{-\frac{1}{2}}; \\
\frac{dy}{dx}
  = \frac{3}{2} x^{\frac{1}{2}} (1 + x^2)^{-\frac{1}{2}}
  + x^{\frac{3}{2}} \times \dfrac{d\bigl[(1+x^2)^{-\frac{1}{2}}\bigr]}{dx}.
\end{gather*}

(2) の例と同様に \( (1+x^2)^{-\frac{1}{2}} \) を微分すると，
\[
\dfrac{d\bigl[(1+x^2)^{-\frac{1}{2}}\bigr]}{dx} = - \frac{x}{\sqrt{(1+x^2)^3}};
\]
したがって，
\[
\frac{dy}{dx}
  = \frac{3\sqrt{x}}{2\sqrt{1+x^2}} - \frac{\sqrt{x^5}}{\sqrt{(1+x^2)^3}}
  = \frac{\sqrt{x}(3+x^2)}{2\sqrt{(1+x^2)^3}}.
\]

(7) \( y = (x + \sqrt{x^2 + x + a})^3 \).

\( x + \sqrt{x^2 + x + a} = u \).
\begin{gather*}
\frac{du}{dx} = 1 + \dfrac{d\bigl[(x^2 + x + a)^{\frac{1}{2}}\bigr]}{dx}. \\
y = u^3; \quad \text{および} \quad \frac{dy}{du} = 3u^2 = 3\left(x + \sqrt{x^2 + x + a}\right)^2.
\end{gather*}

Now let $(x^2+x+a)^{\frac{1}{2}}=v$ and $(x^2+x+a) = w$.
\begin{align*}
\frac{dw}{dx} &= 2x+1;\quad v = w^{\frac{1}{2}};\quad \frac{dv}{dw} = \frac{1}{2}w^{-\frac{1}{2}}. \\
\frac{dv}{dx} &= \frac{dv}{dw} \times \frac{dw}{dx} = \frac{1}{2}(x^2+x+a)^{-\frac{1}{2}}(2x+1). \\
\text{Hence}
\frac{du}{dx} &= 1 + \frac{2x+1}{2\sqrt{x^2+x+a}}, \\
\frac{dy}{dx} &= \frac{dy}{du} \times \frac{du}{dx}\\
&= 3\left(x+\sqrt{x^2+x+a}\right)^2
      \left(1 +\frac{2x+1}{2\sqrt{x^2+x+a}}\right).
\end{align*}

(8) Differentiate $y=\sqrt{\dfrac{a^2+x^2}{a^2-x^2}} \sqrt[3]{\dfrac{a^2-x^2}{a^2+x^2}}$.

We get
\begin{align*}
y &= \frac{(a^2+x^2)^{\frac{1}{2}} (a^2-x^2)^{\frac{1}{3}}}
          {(a^2-x^2)^{\frac{1}{2}} (a^2+x^2)^{\frac{1}{3}}}
  = (a^2+x^2)^{\frac{1}{6}} (a^2-x^2)^{-\frac{1}{6}}. \\
\frac{dy}{dx} &= (a^2+x^2)^{\frac{1}{6}} \frac{d\bigl[(a^2-x^2)^{-\frac{1}{6}}\bigr]}{dx}
   + \frac{d\bigl[(a^2+x^2)^{\frac{1}{6}}\bigr]}{(a^2-x^2)^{\frac{1}{6}}\, dx}.
\end{align*}

Let $u = (a^2-x^2)^{-\frac{1}{6}}$ and $v = (a^2 - x^2)$.
\begin{align*}
u &= v^{-\frac{1}{6}};\quad
  \frac{du}{dv} = -\frac{1}{6}v^{-\frac{7}{6}};\quad
  \frac{dv}{dx} = -2x. \\
%
\frac{du}{dx} &= \frac{du}{dv} \times \frac{dv}{dx} = \frac{1}{3}x(a^2-x^2)^{-\frac{7}{6}}.
\end{align*}

Let $w = (a^2 + x^2)^{\frac{1}{6}}$ and $z = (a^2 + x^2)$.
\begin{align*}
w &= z^{\frac{1}{6}};\quad
  \frac{dw}{dz} = \frac{1}{6}z^{-\frac{5}{6}};\quad
  \frac{dz}{dx} = 2x. \\
%
\frac{dw}{dx} &= \frac{dw}{dz} \times \frac{dz}{dx} = \frac{1}{3} x(a^2 + x^2)^{-\frac{5}{6}}.
\end{align*}

Hence
\begin{align*}
\frac{dy}{dx} &= (a^2+x^2)^{\frac{1}{6}} \frac{x}{3(a^2-x^2)^{\frac{7}{6}}}
   + \frac{x}{3(a^2-x^2)^{\frac{1}{6}} (a^2+x^2)^{\frac{5}{6}}}; \\
%
\text{or}
\frac{dy}{dx} &= \frac{x}{3}
     \left[\sqrt[6]{\frac{a^2+x^2}{(a^2-x^2)^7}}
           + \frac{1}{\sqrt[6]{(a^2-x^2)(a^2+x^2)^5}} \right].
\end{align*}

(9) Differentiate $y^n$ with respect to\textasciitilde{}$y^5$.
\[
\frac{d(y^n)}{d(y^5)} = \frac{ny^{n-1}}{5y^{5-1}} = \frac{n}{5} y^{n-5}.
\]

(10) Find the first and second differential coefficients of $y = \frac{x}{b} \sqrt{(a-x)x}$.
\[
\frac{dy}{dx} = \frac{x}{b}\,
    \frac{d\bigl\{\bigl[(a-x)x\bigr]^{\frac{1}{2}}\bigr\}}{dx}
  + \frac{\sqrt{(a-x)x}}{b}.
\]

Let \((a-x)x = w\) and \(\bigl[(a-x)x\bigr]^{\frac{1}{2}} = u\); then \(u = w^{\frac{1}{2}}\).
\[
\frac{du}{dw}
  = \frac{1}{2} w^{-\frac{1}{2}}
  = \frac{1}{2w^{\frac{1}{2}}} = \frac{1}{2\sqrt{(a-x)x}}.
\]
\[
\frac{dw}{dx} = a-2x.
\]
\[
\frac{du}{dw} \times \frac{dw}{dx} = \frac{du}{dx} = \frac{a-2x}{2\sqrt{(a-x)x}}.
\]

従って
\[
\frac{dy}{dx}
  = \frac{x(a-2x)}{2b\sqrt{(a-x)x}} + \frac{\sqrt{(a-x)x}}{b}
  = \frac{x(3a-4x)}{2b\sqrt{(a-x)x}}.
\]

今
\[
\frac{d^2y}{dx^2}
  = \frac{2b \sqrt{(a-x)x}\, (3a-8x)
           - \frac{(3ax-4x^2)b(a-2x)}{\sqrt{(a-x)x}}}
          {4b^2(a-x)x} \\
  = \frac{3a^2-12ax+8x^2}{4b(a-x)\sqrt{(a-x)x}}.
\]

(これらの最後の2つの微分係数は後で必要となる。参照: Ex. X. No. 11.)

練習問題 VI (解答は Pageref[page]{AnsEx:VI} を参照のこと。)

次の関数を微分せよ:

\begin{enumerate}
\item[(1)] \(y = \sqrt{x^2 + 1}\).
\item[(2)] \(y = \sqrt{x^2+a^2}\).

\item[(3)] \(y = \frac{1}{\sqrt{a+x}}\).
\item[(4)] \(y = \frac{a}{\sqrt{a-x^2}}\).

\item[(5)] \(y = \frac{\sqrt{x^2-a^2}}{x^2}\).
\item[(6)] \(y = \frac{\sqrt[3]{x^4+a}}{\sqrt[2]{x^3+a}}\).

\item[(7)] \(y = \frac{a^2+x^2}{(a+x)^2}\).

\item[(8)] \(y^5\) を \(y^2\) について微分せよ。

\item[(9)] \(y = \frac{\sqrt{1 - \theta^2}}{1 - \theta}\) を微分せよ。

この過程は3つ以上の微分係数にも拡張できるので、\(\frac{dy}{dx} = \frac{dy}{dz} \times \frac{dz}{dv} \times \frac{dv}{dx}\) となる。

例題.
(1) \(z = 3x^4\); \(v = \frac{7}{z^2}\); \(y = \sqrt{1+v}\), \(\frac{dv}{dx}\) を求めよ。

すると
\[
\frac{dy}{dv} = \frac{1}{2\sqrt{1+v}}; \quad
\frac{dv}{dz} = -\frac{14}{z^3}; \quad
\frac{dz}{dx} = 12x^3.
\]
\[
\frac{dy}{dx} = -\frac{168x^3}{(2\sqrt{1+v})z^3}
               = -\frac{28}{3x^5\sqrt{9x^8+7}}.
\]

(2) \( t = \frac{1}{5\sqrt{\theta}} \); \( x = t^3 + \frac{t}{2} \); \( v = \frac{7x^2}{\sqrt[3]{x-1}} \)，求 \(\frac{dv}{d\theta}\)。

\[
\begin{aligned}
\frac{dv}{dx} &= \frac{7x(5x-6)}{3\sqrt[3]{(x-1)^4}}; \\
\frac{dx}{dt} &= 3t^2 + \frac{1}{2}; \\
\frac{dt}{d\theta} &= -\frac{1}{10\sqrt{\theta^3}}. \\
\end{aligned}
\]
したがって
\[
\frac{dv}{d\theta}
  = -\frac{7x(5x-6)(3t^2+\frac{1}{2})}
          {30\sqrt[3]{(x-1)^4} \sqrt{\theta^3}},
\]
この式において \( x \) をその値、および \( t \) を \(\theta\) に関する値に置き換える必要がある。

(3) \( \theta = \frac{3a^2x}{\sqrt{x^3}} \); \( \omega = \frac{\sqrt{1-\theta^2}}{1+\theta} \); \( \phi = \sqrt{3} - \frac{1}{\omega\sqrt{2}} \)，求 \(\frac{d\phi}{dx}\)。

\[
\begin{aligned}
\theta &= 3a^2x^{-\frac{1}{2}}; \\
\omega &= \sqrt{\frac{1-\theta}{1+\theta}}; \\
\phi &= \sqrt{3} - \frac{1}{\sqrt{2}} \omega^{-1}. \\
\end{aligned}
\]
\[
\begin{gathered}
\frac{d\theta}{dx} = -\frac{3a^2}{2\sqrt{x^3}}; \\
\frac{d\omega}{d\theta} = -\frac{1}{(1+\theta)\sqrt{1-\theta^2}}
\end{gathered}
\]
（例5参照、ページ \pageref{examples4}）；そして
\[
\frac{d\phi}{d\omega} = \frac{1}{\sqrt{2}\omega^2}.
\]

したがって \(\frac{d\theta}{dx} = \frac{1}{\sqrt{2} $\times$ \omega2}
  $\times$ \frac{1}{(1+\theta) \sqrt{1-\theta2}}
  $\times$ \frac{3a2}{2\sqrt{x3}}\)。

次にまず \(\omega\)、その後 \(\theta\) をその値で置き換える。

問題集 VII

あなたは次の問題を成功裏に解くことができます。（解答はページ \pageref[page]{AnsEx:VII} を参照してください。）

\item (1) \( u = \frac{1}{2}x^3 \); \( v = 3(u+u^2) \); \( w = \frac{1}{v^2} \)，求 \(\frac{dw}{dx}\)。

\item (2) \( y = 3x^2 + \sqrt{2} \); \( z = \sqrt{1+y} \); \( v = \frac{1}{\sqrt{3}+4z} \)，求 \(\frac{dv}{dx}\)。

\item (3) \( y = \frac{x^3}{\sqrt{3}} \); \( z = (1+y)^2 \); \( u = \frac{1}{\sqrt{1+z}} \)，求 \(\frac{du}{dx}\)。

\end{enumerate}



\chapter{微分の幾何学的意味}
\noindent\textit{\large Geometrical Meaning of Differentiation}\\[
2mm]
\noindent\textbf{\small 【章の論理とテーマ】接線の傾き・dy/dx の図形的可視化・曲線の傾斜の決定}\\[4mm]
微分の幾何学的意味について考えるのは有用です。
まず、$x$ の任意の関数，たとえば $x^2$，または $\sqrt{x}$，または $ax+b$ は曲線としてプロットすることができます；そして現在では、学校の子供たちも座標軸 $OX$ および $OY$ を用いた曲線プロットの過程に馴染みがあります。
図7（参照図7）において、$PQR$ が $OX$ 軸と $OY$ 軸に対する座標系でプロットされた曲線の一部であるとします。この曲線上の任意の点 $Q$ を考慮し、その点の横座標は $x$、縦座標は $y$ とします。今、$x$ の値が変化するとき $y$ がどのように変化するかを観察しましょう。もし $x$ に小さな増分 $dx$ を加えると、右側へ移動すると、この特定の曲線では $y$ も（実際には）小さな増分 $dy$ で増加することが観察されます（この特定の曲線は上昇曲線であるため）。すると、$dy$ に対する $dx$ の比は、点 $Q$ と $T$ 間の曲線がどの程度傾いているかを測る尺度となります。実際には、図から見ると、$Q$ と $T$ 間の曲線は多くの異なる傾きを持っており、そこで「曲線の傾き」という言葉を使うのは難しいです。しかし、もし $Q$ と $T$ が非常に近接して、曲線の小さな部分 $QT$ が実質的に直線である場合、$dy/dx$ の比は $QT$ 間の曲線の傾きを正確に表すことができます。$QT$ を両端から延長した直線は、$QT$ の部分のみで曲線と接しており、この部分が無限に小さければ、直線は実質的に一点のみで曲線と接し、従って $QT$ における接線となります。
この接線は明らかに $QT$ と同じ傾きを持っていますので、$\dfrac{dy}{dx}$ は点 $Q$ における接線の傾きを表します。
観察すると、$dx$は右に短いステップを表し、$dy$はそれに応じて上に短いステップを表します。これらのステップは可能な限り短く考慮されるべきであり——実際には無限に短い——ただし図ではそれらを視認できるように有限に小さいと表現せざるを得ません。
\textit{ここから$\dfrac{dy}{dx}$が任意の点における曲線の傾きを表すという事実を利用することがしばしば必要となる。}
Figure 90a
ある特定の点で45$^\circ$に傾いている曲線があると仮定し、図8を参照してください。この場合、$dy$と$dx$は等しくなりますので、$\dfrac{dy}{dx} = 1$となります。
\textit{より急な傾き（図9）の場合、$\dfrac{dy}{dx}$は1よりも大きくなります。}
Figure 90a, Figure 90b
非常に緩やかな傾斜の曲線があると仮定し、図10を参照してください。この場合、$\dfrac{dy}{dx}$は1より小さな分数となります。
水平な直線の場合、または曲線上の水平部分では$dy=0$であり、従って$\dfrac{dy}{dx}=0$です。
\textit{下向きに傾いている曲線（図11）の場合、$dy$は下にステップを踏むものであり、そのため負の値とみなされるべきです。したがって$\dfrac{dy}{dx}$も負の符号を持つことになります。}
直線である「曲線」がある場合、図12を参照してください。その直線上のどの点でも$\dfrac{dy}{dx}$の値は同じです。言い換えれば、その傾きは一定です。
\textit{右に進むにつれてより上向きに傾いている曲線があると仮定し、図13を参照してください。この場合、$\dfrac{dy}{dx}$の値は傾きが増すとともにますます大きくなります。}
\textit{一方で、進むにつれて徐々に平らになる曲線があると仮定し、図14を参照してください。この場合、$\dfrac{dy}{dx}$の値は平らな部分に達するにつれてますます小さくなります。}
Figure 90c, Figure 91a, Figure 91b, Figure 92a, Figure 92b, Figure 93a, Figure 93b
曲線が最初に下りてから再び上昇する場合、例えば \Fig{15} のように上向きの凹みを呈するとき、明らかに $\dfrac{dy}{dx}$ は最初に負であり、曲線が平坦になるにつれてその値は減少し、曲線の底点に達したときに $0$ になります。この点から先には $\dfrac{dy}{dx}$ の値は正になり、さらに増大します。このような場合、$y$ は \emph{最小} 値を通過すると言います。$y$ の最小値は必ずしも $y$ の最も小さい値とは限らず、それは曲線の底点に対応する $y$ の値です。例えば、\Fig{28}（\Pageref{fig:28}）では、曲線の底点に対応する $y$ の値は $1$ ですが、他の場所ではこの値よりも小さい $y$ の値が存在します。最小値の特徴は、その両側で $y$ が増加することです。
--- 特定の $x$ の値が $y$ を \emph{最小} 值にするとするとき、$\dfrac{dy}{dx}=0$ となります。
曲線が最初に上昇してから再び下落する場合、$\dfrac{dy}{dx}$ の値は最初に正であり、頂点に達したときに $0$ になります。その後、曲線が下向きに傾斜すると負になります。例えば \Fig{16} のようにです。このような場合、$y$ は \emph{最大} 值を通過すると言います。ただし、$y$ の最大値は必ずしも $y$ の最も大きい値とは限らず、\Fig{28} では $y$ の最大値は $2\frac{1}{3}$ ですが、これは他の曲線の点で $y$ が持つことができる最大値ではありません。
\par\textbf{〔第094b図〕}\par{16}{17}
--- 特定の $x$ の値が $y$ を \emph{最大} 值にするとするとき、$\dfrac{dy}{dx}=0$ となります。
曲線が \Fig{17} のような特殊な形状を持つ場合、$\dfrac{dy}{dx}$ の値は常に正です。しかし、その中で最も傾斜の少ない特定の場所がある点があり、この点での $\dfrac{dy}{dx}$ の値は最小であり、つまり他の曲線のどの部分よりも小さいです。
曲線が \Fig{18} のような形状を持つ場合、上部では $\dfrac{dy}{dx}$ の値は負であり、下部では正になります。また、曲線の鼻の部分で実際には垂直になる場所では、$\dfrac{dy}{dx}$ の値は無限大になります。
現在、$\dfrac{dy}{dx}$ が任意の点での曲線の傾斜を測定することを理解したので、既に学んだ微分法のいくつかの式に取り組みましょう。
(1) 最も単純な場合を次のように取る：
\[
y = x + b.
\]

これは \Fig{19} で等しいスケールを使用してプロットされている。$x$ と $y$ のスケールは同じである。$x = 0$ を代入すると、対応する縦軸の値は $y = b$ となる；すなわち、「曲線」 \par\textbf{〔第096b図〕}\par{19}{20} は $y$ 軸上の高さ $b$ 处で交差する。そこから、45$^\circ$ の角度で上昇する；右側に与えるどの値の $x$ も、等しい $y$ 値を上昇させる。直線の傾きは 1:1 である。

次に、既に学んだルール（pp. \pageref{diffrule1} および \pageref{diffrule2} antehere）を使用して $y = x + b$ を微分すると、$\dfrac{dy}{dx} = 1$ となる。

直線の傾きは、右側に小さなステップ $dx$ のたびに等しい小さなステップ $dy$ 上昇するという性質を持っている。そしてこの傾きは一定であり、常に同じ傾きである。

(2) 次の場合を取る：Pagelabel{Case2}
\[
y = ax + b.
\]
この曲線は、前記のものと同様に $y$ 軸上の高さ $b$ 处から始まる。しかし、曲線を描く前に、その傾きを求めることで微分を行う；これにより $\dfrac{dy}{dx} = a$ となる。傾きは一定であり、ここで $\tan$ の値が $a$ と呼ばれる角度である。数値を与えることを考慮し、例えば $a = \frac{1}{3}$ とする。それでは、$y$ 値を 3 分の 1 上昇させるような傾きを与えなければならない；すなわち、$dx$ は $dy$ の 3 倍である必要があり、これは \Fig{21} で拡大表示されている。したがって、\Fig{20} にこの傾きで線を描く。

\par\textbf{〔第21図〕}\par

(3) 略微難しくなった場合を考える。
\begin{align*}
\lintertext{\indent 令}
y = ax^2 + b.
\end{align*}

また、この曲線は $y$ 軸上の原点より高さ $b$ 处から始まる。

次に微分を行う。[忘れていた場合は \Pageref{diffrule2} を参照；または、むしろ、微分の方法を考え出すことを.]

\[
\dfrac{dy}{dx} = 2ax.
\]

\par\textbf{〔第22図〕}\par

これは傾斜が一定でないことを示している：$x$ の増加に従って増大する。原点の始点 $P$ 处、すなわち $x = 0$ の位置では、曲線 (\Fig{22}) は傾斜を持たない---つまり、水平である。原点の左側で $x$ が負の値を持つとき、$\dfrac{dy}{dx}$ も負の値を持つか、すなわち左から右へと下降するようになる（図参照）。

この例でそれを説明しましょう。方程式
\[
y = \frac{1}{4}x^2 + 3,
\]
を考え、これを微分すると
\[
\frac{dy}{dx} = \frac{1}{2}x.
\]
次に、$x$ の値を $0$ から $5$ まで順番にいくつか設定し、最初の方程式を使って $y$ の対応する値を計算し、二つ目の方程式から $\frac{dy}{dx}$ の値を計算します。結果を表にまとめると次のようになります：
\[
\begin{array}{|c||*{6}{c|}}
\hline
x & 0 & 1 & 2 & 3 & 4 & 5 \\
\hline
y & 3 & 3\frac{1}{4} & 4 & 5\frac{1}{4} & 7 & 9\frac{1}{4} \\
\hline
\frac{dy}{dx} & 0 & \frac{1}{2} & 1 & 1\frac{1}{2} & 2 & 2\frac{1}{2} \\
\hline
\end{array}
\]
これらの値を二つの曲線でプロットし、図23と図24に表します。図23では $y$ の値を $x$ の値に対してプロットし、図24では $\frac{dy}{dx}$ の値を $x$ の値に対してプロットします。任意の $x$ の値に対する正規線の高さは、最初の曲線の傾きと比例します。

図2.5in 099a 099b 23 24

曲線が急な尖端を持つ場合、図25のように、その点での傾きは突然上向きから下向きへと変化します。その場合、$\frac{dy}{dx}$ は明確に正の値から負の値へと急激に変化します。

図2.5in 099c 25

次の例では、先ほど説明した原理のさらなる応用を示します。

(4) $x = -1$ の点での曲線
\[
y = \frac{1}{2x} + 3,
\]
に対する接線の傾きを求めます。また、この接線が曲線 $y = 2x^2 + 2$ と作る角度も求めます。

接線の傾きは、曲線と接する点での曲線の傾きであり（参照：\Pageref{slope}）、これは点における曲線の方程式の$\dfrac{dy}{dx}$に等しい。ここで$\dfrac{dy}{dx} = -\dfrac{1}{2x^2}$で、$x = -1$のとき$\dfrac{dy}{dx} = -\dfrac{1}{2}$となり、これが接線と曲線の傾きである。接線は直線であり、方程式は$y = ax + b$となる。その傾きは$\dfrac{dy}{dx} = a$で、よって$a = -\dfrac{1}{2}$である。また$x = -1$のとき$y = \dfrac{1}{2(-1)} + 3 = 2\dfrac{1}{2}$であり、接線を通っているため、この点の座標は接線の方程式を満たすので
\[
y = -\dfrac{1}{2} x + b,
\]
である。よって$2\dfrac{1}{2} = -\dfrac{1}{2} \times (-1) + b$から$b = 2$となり、接線の方程式は$y = -\dfrac{1}{2} x + 2$である。

二つの曲線が交わるとき、その交点は両方の曲線共通の点であり、その座標はそれぞれの曲線の方程式を満たす。つまり、それは二つの方程式からなる連立方程式の解である。ここでは、次の連立方程式の解によって曲線が交わる
\begin{gather*}
\left\{
\begin{aligned}
  y &= 2x^2 + 2, \\
  y &= -\dfrac{1}{2} x + 2 \quad\text{または}\quad 2x^2 + 2 = -\dfrac{1}{2} x + 2;
\end{aligned}
\right.  \\
\lintertext{すなわち}
x(2x + \dfrac{1}{2}) = 0.
\end{gather*}

この方程式の解は$x = 0$と$x = -\dfrac{1}{4}$である。曲線$y = 2x^2 + 2$における任意の点での傾きは
\[
\dfrac{dy}{dx} = 4x.
\]

$x = 0$の点ではこの傾きはゼロであり、曲線は水平である。$x = -\dfrac{1}{4}$の点では$\dfrac{dy}{dx} = -1$となり、その点での曲線は右下に傾いており、水平方向と角度$\theta$で交わるが、$\tan \theta = 1$であり、つまり水平方向に対して$45^\circ$である。

直線の傾きは$-\dfrac{1}{2}$であり、これは右下に傾いており、水平方向との角度$\phi$を形成し、$\tan \phi = \dfrac{1}{2}$である。すなわち、$26^\circ~34'$である。したがって最初の点では曲線は直線と角度$26^\circ~34'$で交わり、二つ目の点では角度$45^\circ - 26^\circ~34' = 18^\circ~26'$で交わる。

直線を点 $(x = 2, y = -1)$ を通るように引く。この直線は曲線 $y = x^2 - 5x + 6$ の接線である。接触点の座標を求めよ。

接線の傾きは曲線の $\dfrac{dy}{dx}$ と等しいので、それは $2x - 5$ である。

直線の方程式は $y = ax + b$ であり、この直線が $x = 2, y = -1$ を満たすためには $-1 = a \times 2 + b$；また、その $\dfrac{dy}{dx} = a = 2x - 5$。

接触点の座標も接線の方程式と曲線の方程式を満たさなければならない。

次のように四つの方程式を得る：
\begin{align*}
  y &= x^2 - 5x + 6, \tag*{(i)}   \\
  y &= ax + b,       \tag*{(ii)}  \\[
-0.6\baselineskip]
% [** TN: Now set the left brace]
\smash{\left\{\begin{aligned} & \\ & \\ & \\ & \\ \end{aligned} \right.} \phantom{-1} \& \\[-1.5ex]
 -1 &= 2a + b,       \tag*{(iii)} \\
  a &= 2x - 5,       \tag*{(iv)}
\end{align*}
これらは $a$, $b$, $x$, $y$ の四つの方程式である。
(i)と(ii)から $x^2 - 5x + 6 = ax + b$ を得る。
$a$ と $b$ の値を代入すると、
\[
  x^2 - 5x + 6 = (2x - 5)x - 1 - 2(2x - 5),
\]
これは $x^2 - 4x + 3 = 0$ に簡約され、解は $x = 3$ と $x = 1$ である。これらを(i)に代入すると、それぞれ $y = 0$ と $y = 2$ を得る；接触点の座標は $(x, y) = (1, 2)$ と $(x, y) = (3, 0)$ である。

注記。曲線に関するすべての問題では、実際に曲線を描画して導き出された結論を確認することにより、学生たちは非常に有益な学習を得ることができる。

練習問題 VIII（解答は \Pageref[page]{AnsEx:VIII} を参照）

\begin{enumerate}
\item (1) $y = \frac{3}{4} x^2 - 5$ の曲線をミリメートルのスケールで描画する。異なる $x$ の値に対応する点での傾き角を測定する。

微分方程式から傾きの式を求め、自然接線表から求めた傾きと一致するか確認する。
\item (2) $x = 2$ の横座標を持つ特定の点における曲線
\[
y = 0.12x^3 - 2,
\]
の傾きを求めよ。
\item (3) $y = (x - a)(x - b)$ の場合、$\dfrac{dy}{dx} = 0$ となる特定の点での $x$ の値は $\tfrac{1}{2}(a + b)$ であることを示せ。

\item(4) 方程式 \( y = x^3 + 3x \) の \(\dfrac{dy}{dx}\) を求めよ；さらに、\( x = 0 \), \( x = \dfrac{1}{2} \), \( x = 1 \), \( x = 2 \) に対応する点での \(\dfrac{dy}{dx}\) の数值を計算せよ。

\item(5) 方程式 \( x^2 + y^2 = 4 \) を表す曲線上で、傾きが \( 1 \) となる点の \( x \) の値を求めよ。

\item(6) 方程式 \(\dfrac{x^2}{3^2} + \dfrac{y^2}{2^2} = 1\) を表す曲線上の任意一点での傾きを求めよ；また、\( x = 0 \), \( x = 1 \) の点での傾きの数值をそれぞれ求めよ。

\item(7) 曲線 \( y = 5 - 2x + 0.5x^3 \) 上で接線が存在し、その方程式が \( y = mx + n \)（\( m \), \( n \) は定数）の形をとるとき、接点の \( x \) 座標が \( 2 \) の場合、\( m \) および \( n \) の値を求めよ。

\item(8) 曲線 \( y = 3.5x^2 + 2 \) および \( y = x^2 - 5x + 9.5 \) が交わる角度を求めよ。

\item(9) 方程式 \( y = ± \sqrt{25-x^2} \) 上で、\( x = 3 \), \( x = 4 \) の点における接線を引いたとき、これらの接線の交点の座標と、その相互な角度を求めよ。

\item(10) 直線 \( y = 2x - b \) が曲線 \( y = 3x^2 + 2 \) 上の一点で接しているとき、接点の座標と \( b \) の値を求めよ。
\end{enumerate}



\chapter{極大値と極小値（最適化問題）}
\noindent\textit{\large Maxima and Minima}\\[
2mm]
\noindent\textbf{\small 【章の論理とテーマ】導関数が0となる点・2階微分の符号による極大極小判定と実用問題への応用}\\[4mm]
第十一章 最大値と最小値
最初の例として、微分法の主な用途の一つは、微分されたものの値が最大または最小になる条件を見つけることです。これはエンジニアリング問題においてしばしば非常に重要であり、作業コストを最小にするか、効率を最大化するための条件を知りたいことがほとんどです。
まず具体的な例から始めましょう。方程式
\[
y = x^2 - 4x + 7.
\]
を考えます。（図105a参照）

$x$に連続的な値を代入し、対応する$y$の値を求めることで、この方程式は最小値を持つ曲線を表していることがわかります。
\[
\begin{array}{|c||c|c|c|c|c|c|}
\hline
x & 0 & 1 & 2 & 3 & 4 & 5 \\
\hline
y & 7 & 4 & 3 & 4 & 7 & 12 \\
\hline
\end{array}
\]

これらの値は図26にプロットされており、$x = 2$のとき$y$が約3の最小値を取ると推測できます。しかし、最小値が$2$でなく$2 \frac{1}{4}$や$1 \frac{3}{4}$であることは確実でしょうか？

もちろん、どんな代数式に対しても多くの値を計算し、その中から最大または最小の特定の値を見つけることができます。

もう一つの例です：

\[
y = 3x - x^2.
\]

いくつかの値を計算してみましょう：
\[
\begin{array}{|c||c|c|c|c|c|c|c|}
\hline
x & -1 & 0 & 1 & 2 & 3 & 4 & 5 \\
\hline
y & -4 & 0 & 2 & 2 & 0 & -4 & -10 \\
\hline
\end{array}
\]
（図106a参照）

これらの値を図27のようにプロットします。

$x = 1$と$x = 2$の間で最大値があることが明らかであり、$y$の実際の最大値は約$2 \frac{1}{4}$であるように見えます。中間のいくつかの値を試してみましょう。$x = 1 \frac{1}{4}$のとき$y = 2.187$、$x = 1 \frac{1}{2}$のとき$y = 2.25$、$x = 1.6$のとき$y = 2.24$です。$2.25$が実際の最大値であり、それが$x = 1 \frac{1}{2}$で発生するかをどのように確認できるでしょうか？

現在、最大値（または最小値）に直行して到達できる方法があることを保証されても、それは多くの前駆的な試行錯誤や推測を必要としないように聞こえるかもしれません。その方法は微分を利用するものです。前のページ（\pageref{curve}参照）の図14と15についての説明に戻ってみてください。曲線が最大または最小の高さに達する点では、その点での \(\frac{dy}{dx} = 0\) です。

この事実は求めたい方法を示唆します。方程式が与えられたとき、その \(y\) の最小値（または最大値）を求めたい場合、まず微分を行います。その後、\(\frac{dy}{dx}\) を0に等しいとします。そして \(x\) について解きます。この特定の \(x\) の値を元の方程式に戻すと、必要な \(y\) の値が得られます。このプロセスは一般的に「0に等しくする」と呼ばれます。

どのように簡単に機能するかを見てみましょう。この章で最初に扱った例を考えます：
\[
y = x^2 - 4x + 7.
\]

微分すると次のようになります：
\[
\frac{dy}{dx} = 2x - 4.
\]

これを0に等しくします：
\[
2x - 4 = 0.
\]

この方程式を \(x\) について解くと、
\begin{align*}
2x &= 4, \\
x &= 2.
\end{align*}

\(x = 2\) のとき、最大値（または最小値）が発生することがわかります。

\(x = 2\) を元の方程式に戻すと、
\begin{align*}
y &= 2^2 - (4 \times 2) + 7 \\
  &= 4 - 8 + 7       \\
  &= 3.
\end{align*}

次に、図26を参照してください。最小値は \(x = 2\) のときに発生し、その最小値 \(y = 3\) です。

次の例（図24）も試してみましょう：
\[
y = 3x - x^2.
\]

微分すると次のようになります：
\[
\frac{dy}{dx} = 3 - 2x.
\]

これを0に等しくします：
\[
3 - 2x = 0,
\]
従って
\[
x = 1 \frac{1}{2}.
\]

この \(x\) の値を元の方程式に戻すと、
\begin{align*}
y &= 4 \frac{1}{2} - (1 \frac{1}{2} \times 1 \frac{1}{2}) \\
y &= 2 \frac{1}{4}.
\end{align*}

これは多くの値を試す方法で不確定だった情報を正確に提供します。

今、さらなる例に進む前に、二つの注意点があります。$\dfrac{dy}{dx}$がゼロと等しいという指示を受けるとき、最初は（あなた自身に少しでも知恵があるなら）一種の不満を感じます。なぜなら、$\dfrac{dy}{dx}$は曲線の異なる部分でさまざまな値を持つからです。上り下りによってその値が異なりますので、突然
\[
\dfrac{dy}{dx} = 0
\]
と書かれることが告げられると、不満を感じ、それが正しいとは思えないと感じるかもしれません。現在のところ、あなたは自己条件を満たす方程式について扱っていることを理解する必要がありますが、そのような方程式は必ずしも真実であるわけではなく、特定の条件が満たされる場合にのみ真実です。そして、それらの方程式を書き下して解くことで、その条件を見つけることができます。

現在の目標は、曲線が上り斜めでも下り斜めでもないとき、つまり$\dfrac{dy}{dx} = 0$となる特定の場所で$x$の値を見つけ出すことです。したがって、$\dfrac{dy}{dx} = 0$と書くことは常に$=0$であることを意味するわけではありません。しかし、$\dfrac{dy}{dx}$がゼロになる条件としてそれを書き下すことで、$x$がどの程度になるかを見つけることができます。

二番目の注意点は、あなた自身に少しでも知恵があるならおそらく既に気づいていることですが、この大いに称賛される$\dfrac{dy}{dx}$をゼロと等しいという過程は、その結果得られる$x$が$y$の最大値か最小値を与えるかどうかを教えてくれません。もちろん、それが最大値か最小値であるかは自身で判断する必要があります。たとえば、曲線を描いていれば既にそれを知っているかもしれません。

例えば、方程式
\[
y = 4x + \frac{1}{x}
\]
を考えます。

微分せずにその曲線に対応するものとは考えず、微分し、0 に等しい値を求めます：
\begin{align*}
\frac{dy}{dx} &= 4 - x^{-2} = 4 - \frac{1}{x^2} = 0; \\
\therefore x &= \frac{1}{2}; \\
\text{これを代入すると} y &= 4
\end{align*}
これが極大か極小であるかどうかは、後で二階微分を利用する方法が教えられるので、現在のところ、別の $x$ の値を少し変えてその対応する $y$ の値が既に求めたものより大きいか小さいか試すだけでも十分です。

最大と最小に関するもう一つの単純な問題を試してみましょう。例えば、ある数を二つの部分に分けて、その積が最大になるように分割したいとします。この Trick は知らずに、どのようにアプローチすればよいでしょうか？おそらく、繰り返しのルールを使って考え出すことができるでしょう。60 を例として考えてみましょう。2 つの部分に分けると、それらを掛け算して積を求めることができます。例えば、$50 \times 10 = 500, 52 \times 8 = 416, 40 \times 20 = 800, 45 \times 15 = 675, 30 \times 30 = 900$。これは最大の積に見えますので、それを変えてみましょう。$31 \times 29 = 899$, これは少し悪く、$32 \times 28 = 896$ もさらに悪いです。したがって、最大の積は二等分で得られることがわかります。

微分法を使ってみましょう。分割する数を $n$ と呼びます。一つの部分を $x$ とするとき、もう一つの部分は $n-x$ であり、その積は $x(n-x)$ または $nx - x^2$ です。よって $y = nx - x^2$ と表します。微分し、0 に等しい値を求めます：
\begin{align*}
\dfrac{dy}{dx} &= n - 2x = 0 \\
\therefore \dfrac{n}{2} &= x.
\end{align*}
これで何らかの数 $n$ に関わらず、部分積が最大になるためには二等分する必要があることがわかります。その最大値は常に $\frac{1}{4} n^2$ となります。

これは非常に有用なルールであり、任意の因子数にも適用されます。つまり、$m + n + p = $ 恒定数であれば、$m \times n \times p$ が最大になるのは $m = n = p$ のときです。

\subsection{テストケース。}
我々の知識を即座に試すために、次のケースを用いましょう。
\begin{align*}
\lintertext{\indent 令} y &= x^2 - x;
\end{align*}
この関数が最大値または最小値を持っているかどうかを見つけるとともに、それが最大値か最小値かをテストします。

微分すると，
\begin{align*}
\frac{dy}{dx} &= 2x - 1. \\
\intertext{\indent 0に等しいとすると，}
2x - 1 &= 0, \\
\lintertext{従って}
2x &= 1, \\
\lintertext{または}
x &= \tfrac{1}{2}.
\end{align*}

つまり、$x = \tfrac{1}{2}$ のとき、対応する $y$ の値は最大値または最小値のいずれかになります。従って、元の方程式に $x = \tfrac{1}{2}$ を代入すると，
\begin{align*}
y &= (\tfrac{1}{2})^2 - \tfrac{1}{2}, \\
\lintertext{または}
y &= -\tfrac{1}{4}.
\end{align*}

これは最大値か最小値でしょうか？これをテストするために、$x$ を $\tfrac{1}{2}$ より少し大きい値にします。例えば $x = 0.6$ とします。すると，
\[
y = (0.6)^2 - 0.6 = 0.36 - 0.6 = -0.24,
\]
これは $-0.25$ より大きいので、$y = -0.25$ は最小値であることが示されます。

自己で曲線を描いて計算結果を確認してください。

\par\textbf{【例題 さらに興味深い例。】}\par
最大値と最小値を持つ曲線の例は非常に興味深いです。その方程式は次の通りです：
\begin{align*}
y &= \tfrac{1}{3} x^3 - 2x^2 + 3x + 1. \\
\lintertext{\indent 微分すると，}
\dfrac{dy}{dx} &= x^2 - 4x + 3.
\end{align*}

\begin{figure}[h]
    \centering
    \includegraphics[width=2.5in]{113a.png}
    \caption{曲線の例}
    \label{fig:28}
\end{figure}

0に等しいとすると，二次方程式が得られます：
\[
x^2 - 4x + 3 = 0;
\]
この二次方程式を解くと、以下の2つの根を得ます：
\[
\left\{
\begin{aligned}
x &= 3 \\
x &= 1.
\end{aligned}
\right.
\]

$x = 3$ のとき $y = 1$; $x = 1$ のとき $y = 2 \tfrac{1}{3}$。これらの中では、最初のものは最小値であり、次のものは最大値です。

この曲線は（図 \ref{fig:28} 参照）計算された値から描くことができます。
\[
\begin{array}{|c||*{8}{c|}}
\hline
\Td[c]{x} & \Td{-1\Z} & \Td[c]{0} & \Td[c]{1\Z} & \Td[c]{2\Z} & \Td[c]{3} & \Td{4\Z} & \Td{5\Z} & \Td{6} \\
\hline
\Td[c]{y} & \Td{-4\tfrac{1}{3}} & \Td[c]{1} &  \Td[c]{2\tfrac{1}{3}} &  \Td[c]{1\tfrac{2}{3}} & \Td[c]{1} & \Td{2\tfrac{1}{3}} & \Td{7\tfrac{2}{3}} & \Td{19} \\
\hline
\end{array}
\]

最大値と最小値のさらなる練習は、次の例から得られます。

円の半径が \( r \)，中心 \( C \) の座標が \( x = a, y = b \) である場合、その円の方程式は以下のようになります：
\[
(y - b)^2 + (x - a)^2 = r^2.
\]

この式を変形すると：
\[
y = \sqrt{r^2 - (x - a)^2} + b.
\]

図から明らかに、\( x = a \) のとき、\( y \) は最大値 \( b + r \) や最小値 \( b - r \) のいずれかになります。しかし、この事実を利用せず、導関数を用いて \( y \) の極値を求めましょう。
\[
\frac{dy}{dx} = \frac{1}{2} \cdot \frac{1}{\sqrt{r^2 - (x - a)^2}} \times (2a - 2x),
\]
これは簡略化して：
\[
\frac{dy}{dx} = \frac{a - x}{\sqrt{r^2 - (x - a)^2}}.
\]

\( y \) の極値となる条件は以下のようになります：
\[
\frac{a - x}{\sqrt{r^2 - (x - a)^2}} = 0.
\]

分母が無限大になることはないため、この式を満たす \( x \) の値は：
\[
x = a.
\]
元の円の方程式に \( x = a \) を代入すると：
\[
y = \sqrt{r^2} + b;
\]
ここで \( r^2 \) の平方根は \( +r \) または \( -r \) であるため、\( y \) の値は以下のようになります：
\[
\left\{
\begin{aligned}
y &= b + r, \\
y &= b - r.
\end{aligned}
\right.
\]
これらのうち、最大値は上部にあり、最小値は下部にあります。

曲線が極値を持たない場合、導関数を0に等しいと仮定した結果は不可能になります。例えば：
\[
y = ax^3 + bx + c.
\]
このとき：
\[
\frac{dy}{dx} = 3ax^2 + b.
\]

これを0に等しくすると：
\[
3ax^2 + b = 0,
\]
\[
x^2 = \frac{-b}{3a}, \quad x = \sqrt{\frac{-b}{3a}},
\]
これは不可能です。したがって、\( y \) の最大値も最小値も存在しません。

さらにいくつかの計算例を解くことで、この微分法による応用を完全にマスターすることができます。

1. 半径 \( R \) の円内に内接する長方形の面積が最大となるときの辺の長さは何ですか？

もし一方の辺を\textasciitilde{}$x$と呼ぶなら、
\[
\text{他方の辺} = \sqrt{(\text{対角線})^2 - x^2};
\]
そして、長方形の対角線は必ず直径であるから、他方の辺\textasciitilde{}$= \sqrt{4R^2 - x^2}$。

すると、長方形の面積 $S = x\sqrt{4R^2 - x^2}$ であり，
\[
\dfrac{dS}{dx} = x \times  \dfrac{d\left(\sqrt{4R^2 - x^2}\,\right)}{dx} + \sqrt{4R^2 - x^2} \times  \dfrac{d(x)}{dx}.
\]

$\sqrt{4R^2-x^2}$ の微分を忘れた場合、ヒントは次の通りです：$4R^2-x^2=w$ と $y=\sqrt{w}$ として、$\dfrac{dy}{dw}$ および $\dfrac{dw}{dx}$ を求め、それらを使って計算しましょう。それでも進まなければ \Pageref[page]{dodge} を参照してください。

すると
\[
\dfrac{dS}{dx}
  = x \times  -\dfrac{x}{\sqrt{4R^2 - x^2}} + \sqrt{4R^2 - x^2}
  = \dfrac{4R^2 - 2x^2}{\sqrt{4R^2 - x^2}}.
\]

最大値または最小値を得るためには
\[
\dfrac{4R^2 - 2x^2}{\sqrt{4R^2 - x^2}} = 0;
\]
すなわち、$4R^2 - 2x^2 = 0$ から $x = R\sqrt{2}$。

他方の辺 ${} = \sqrt{4R^2 - 2R^2} = R\sqrt{2}$；両辺は等しく、これは半径を底とする正方形の対角線に等しい。この場合、明らかに最大値である。

(2) 極大容量を持つ円錐容器の斜辺が長さ\textasciitilde{}$l$ のとき、その開口部の半径は？

$R$ を半径、$H$ を対応する高さとすると、$H = \sqrt{l^2 - R^2}$。
\[
\text{体積 } V = \pi R^2 \times  \dfrac{H}{3} = \pi R^2 \times  \dfrac{\sqrt{l^2 - R^2}}{3}.
\]

前問と同様の手順で計算すると
\begin{align*}
\dfrac{dV}{dR}
  &= \pi R^2 \times  -\dfrac{R}{3\sqrt{l^2 - R^2}} + \dfrac{2\pi R}{3} \sqrt{l^2 - R^2} \\
  &= \dfrac{2\pi R(l^2 - R^2) - \pi R^3}{3\sqrt{l^2 - R^2}} = 0
\end{align*}
が最大値または最小値を得る条件となる。

すなわち、$2\pi R(l^2 - R^2) - \pi R^2 = 0$ から $R = l\sqrt{\tfrac{2}{3}}$ であり、これは極大であることが明らかである。

(3) 関数
\[
y = \dfrac{x}{4-x} + \dfrac{4-x}{x}
\]
の最大値と最小値を求めよ。

すると
\[
\dfrac{dy}{dx} = \dfrac{(4-x)-(-x)}{(4-x)^2} + \dfrac{-x - (4-x)}{x^2} = 0
\]
が最大値または最小値を得る条件となり、すなわち
\[
\dfrac{4}{(4-x)^2} - \dfrac{4}{x^2} = 0 \quad\text{および}\quad x = 2.
\]

There is only one value, hence only one maximum or minimum.
\begin{align*}
x &= 2 \quad y = 2, \\
x &= 1.5 \quad y = 2.27,   \\
x &= 2.5 \quad y = 2.27;
\end{align*}
it is therefore a minimum. (グラフをプロットすることで理解が深まる.)

(4) 関数 $y = \sqrt{1+x} + \sqrt{1-x}$ の最大値と最小値を求めよ。グラフをプロットすることも有益である。

微分するとすぐに得られる（例題 No.\textasciitilde{}1 参照、\pageref{ExNo1}）
\[
\dfrac{dy}{dx} = \dfrac{1}{2\sqrt{1+x}} - \dfrac{1}{2\sqrt{1-x}} = 0
\]
が最大値または最小値の条件。

したがって $\sqrt{1+x} = \sqrt{1-x}$ で $x = 0$ が唯一の解

$x=0$ のとき、$y=2$。

$x=\pm0.5$ のとき、$y=1.932$ なのでこれは最大値である。

(5) 関数
\[
y = \dfrac{x^2-5}{2x-4}.
\]
の最大値と最小値を求めよ。

微分すると
\[
\dfrac{dy}{dx} = \dfrac{(2x-4) \times  2x - (x^2-5)2}{(2x-4)^2} = 0
\]
が最大値または最小値の条件；または
\[
\dfrac{2x^2 - 8x + 10}{(2x - 4)^2} = 0;
\]
または $x^2 - 4x + 5 = 0$; この方程式の解は
\[
x = \tfrac{5}{2} ± \sqrt{-1}.
\]

これらが虚数であるため、$\dfrac{dy}{dx} = 0$ を満たす実数の $x$ の値はない；したがって最大値も最小値もない。

(6) 関数
\[
(y-x^2)^2 = x^5.
\]
の最大値と最小値を求めよ。

この式は $y = x^2 ± x^{\frac{5}{2}}$ と書き換えることができる。
\[
\dfrac{dy}{dx} = 2x ± \tfrac{5}{2} x^{\frac{3}{2}} = 0 \quad\text{が最大値または最小値の条件};
\]
すなわち $x(2 ± \tfrac{5}{2} x^{\frac{1}{2}}) = 0$ で、これは $x = 0$ のときと $2 ± \tfrac{5}{2} x^{\frac{1}{2}} = 0$ （すなわち $x=\tfrac{16}{25}$）のときを満たす。したがって二つの解がある。

まず $x = 0$ の場合を考える。$x = -0.5$ のとき、$y = 0.25 ± \sqrt[2]{-(0.5)^5}$, そして $x = +0.5$ のとき、$y = 0.25 ± \sqrt[2]{(0.5)^5}$. 一方が虚数であるため、$y$ の値はグラフ上に表示できない；したがってこのグラフは $y$ 軸の右側全体に存在する（図 Fig{30} 参照）。

グラフをプロットすると、

曲線は原点に向かって進み、最小値があるかのように見えるが、実際には最小値を超えて進むべきところを戻る（これは「尖点」と呼ばれる）。したがって最小値は存在せず、$\dfrac{dy}{dx} = 0$ の条件は満たされているにもかかわらずである。常に一方に一つの値を取り、確認する必要がある。

今、$x = \frac{16}{25} = 0.64$とします。このとき、$x = 0.64$のとき、$y = 0.7373$で、$y = 0.0819$；$x = 0.6$のとき、$y$は$0.6389$と$0.0811$に、$x = 0.7$のとき、$y$は$0.8996$と$0.0804$になります。

これは曲線が2つの枝を持つことを示しています；上部の枝は最大値を通過しないのに対し、下部の枝は最大値を通過します。

（7）底面の半径が高さの2倍である円柱が体積を増加させている場合、その各部分は常に同じ比で変化しており、すなわち、任意の瞬間において、円柱は元の円柱に相似です。底面の半径が$r$フィートのとき、表面積は毎秒20平方インチで増加しています。そのとき、体積は何倍每秒で増加していますか？

\begin{align*}
\text{面積}   &= S = 2(\pi r^2) + 2 \pi r \times  2r = 6 \pi r^2.\\
\text{体積}    &= V = \pi r^2 \times  2r = 2 \pi r^3.\\
\frac{dS}{dr}   &= 12\pi r, \quad \frac{dV}{dr}=6 \pi r^2,\\
dS             &= 12\pi r\, dr=20, \quad dr=\frac{20}{12 \pi r},\\
dV             &=  6\pi r^2\, dr = 6 \pi r^2 \times  \frac{20}{12 \pi r} = 10r.
\end{align*}

体積は$10r$立方インチで変化します。

他にも自分で例を考えてみてください。興味深い例が豊富な科目は少ないです。

演習問題 IX（解答は Pageref[page]{AnsEx:IX} を参照してください）

\begin{enumerate}
\item[(1)] $y = \frac{x^2}{x+1}$のとき、$y$が最大値と最小値となる$x$の値は何ですか？

\item[(2)] $y = \frac{x}{a^2+x^2}$のとき、$y$が最大値となる$x$の値は何ですか？

\item[(3)] 長さ$p$の線を4つの部分に切り分けて長方形を作る場合、各辺が$\frac{1}{4}p$であるとき、長方形の面積は最大になります。

\item[Ex9No4]{(4)} 30インチのひもを3つの釘で張り、三角形を作ります。このひもの内接できる最大の三角形の面積は何ですか？

\item[(5)] 方程式
\[
y = \frac{10}{x} + \frac{10}{8-x}
\]
に対応する曲線をプロットし、$\frac{dy}{dx}$を求め、$y$が最小値となる$x$の値を求め、その最小値を求めます。

\item[(6)] $y = x^5 - 5x$のとき、$y$が最大値または最小値となる$x$の値は何ですか？

\item[(7)] 与えられた正方形に内接できる最も小さな正方形は何か？

\item[(8)] 与えられた高さが底面の半径と等しい円錐体内に、体積が最大となる（a）円柱、（b）側面積が最大となる（c）円柱を内接させる。

\item[(9)] 与えられた球体内に、体積が最大となる（a）円柱、（b）側面積が最大となる（c）円柱を内接させる。

\item[(10)] 気球の体積は増大している。半径が \( r \) フィートであるとき、その体積が毎秒 4 立方フィートで増大している場合、その表面積はどのくらい増大しているか。

\item[(11)] 与えられた球体内に、体積が最大となる円錐を内接させる。

\item[(12)] \( N \) 個の相似なボルタ電池からなるバッテリーによって与えられる電流 \( C \) は \( C = \dfrac{n\times E}{R+\dfrac{rn^2}{N}} \) で与えられ、\( E \), \( R \), \( r \) は定数であり、\( n \) は系列に接続されたセルの数である。電流が最大となる \( n \) に対する比例を求める。
\end{enumerate}



\chapter{曲線の曲率}
\noindent\textit{\large Curvature of Curves}\\[
2mm]
\noindent\textbf{\small 【章の論理とテーマ】曲線の曲がり具合・曲率半径の微分幾何学的定式化}\\[4mm]
第十二章 曲線の曲率
連続的な微分法について再び考えると、なぜ二次微分が必要なのか疑問が湧くかもしれません。変数が空間と時間であるとき、二次微分すると物体の加速度が得られますが、幾何学的解釈では \(\frac{dy}{dx}\) は曲線の傾きを意味します。しかし、\(\frac{d^2 y}{dx^2}\) は何を表していますか？明らかにそれは、傾きが単位長さ \(x\) 毎でどの程度変化しているか（すなわち、傾きの曲率を測る指標）を示します。
図124a, 124b
まず、図31のような定数の傾きを考えます。
この場合、\(\frac{dy}{dx}\) は一定の値を持っています。
しかし、図32のように、傾き自体が上向きに増加していると仮定すると、\(\frac{d\left(\frac{dy}{dx}\right)}{dx}\)，すなわち \(\frac{d^2y}{dx^2}\) は正の値になります。
また、図14（参照：第14ページ）や図33のように、右に移動するにつれて傾きが減少している場合、曲線が上向きに進んでいるとしても、その \(\frac{d^2y}{dx^2}\) は負の値になります。
現在、最大値または最小値を得た結果が最大か最小であるかどうかを判別する方法について秘密を教えます。手順は以下の通りです：最初に微分して得られた式を0に等しいと仮定した後、さらに一度微分し、その結果が正の値か負の値かを見るのです。もし \(\frac{d^2y}{dx^2}\) が正であれば、得られた \(y\) の値は最小値であることが分かります；しかし、それが負であれば、得られた \(y\) の値は最大値です。これがルールです。
曲線が最小点を持つ場合（図15や34参照）、または最小点 \( y \) の位置 \( M \) があり、曲線は上向きに凸である場合を考えましょう。\( M \) の左側では傾きは下向きで、つまり負であり、さらに負の値になることがありません。一方、\( M \) の右側では傾きが上向きになり始め、さらに上向きになります。明らかに曲線が \( M \) を通るときの傾きの変化は、\(\frac{d^2y}{dx^2}\) が正であることを示しています。これは \( x \) が右側へ増加するにつれて、下向きの傾きを上向きの傾きに変換することを意味します。
同様に、最大点を持つ曲線（図16や35参照）を考えましょう。この場合、曲線は凸であり、最大点 \( M \) がマークされています。この場合、\( M \) を通って \( x \) が右側へ移動するとき、上向きの傾きは下向きまたは負の傾きに変換されます。したがって、この場合、「傾きの傾き」\(\frac{d^2y}{dx^2}\) は負であることがわかります。
最後の章の例に戻り、このような結論を確認してみてください。特定の場合で最大値または最小値があるかどうかを検証してください。以下にいくつかの計算例があります。
1. \( y = 4x^2 - 9x - 6 \) および \( y = 6 + 9x - 4x^2 \) の最大値や最小値を求めましょう。それぞれの場合で最大値か最小値かを確認してください。
   (a) \(\frac{dy}{dx} = 8x - 9 = 0\) より \( x = 1\frac{1}{8} \), そして \( y = -11.065 \).
   \(\frac{d^2y}{dx^2} = 8\); 它は正であるので、これは最小値です。
   (b) \(\frac{dy}{dx} = 9 - 8x = 0\) より \( x = 1\frac{1}{8} \), そして \( y = +11.065 \).
   \(\frac{d^2y}{dx^2} = -8\); 它は負であるので、これは最大値です。
2. 関数 \( y = x^3 - 3x + 16 \) の最大値と最小値を求めましょう。
   \(\frac{dy}{dx} = 3x^2 - 3 = 0\) より \( x^2 = 1 \), そして \( x = ±1 \).
   \(\frac{d^2y}{dx^2} = 6x\); \( x = 1 \) のとき、それは正であるので、\( x = 1 \) は最小値 \( y = 14 \) を示します。
   \( x = -1 \) のとき、それは負であるので、\( x = -1 \) は最大値 \( y = +18 \) を示します。
(3) y=$\frac{x-1}{x^2+2}$の極値を求めよ。
\[
\frac{dy}{dx} = \frac{(x^2+2) \times  1 - (x-1) \times  2x}{(x^2+2)^2}
  = \frac{2x - x^2 + 2}{(x^2 + 2)^2} = 0;
\]
または $x^2 - 2x - 2 = 0$ の解は $x = +2.73$ および
$x = -0.73$。
\begin{align*}
\dfrac{d^2y}{dx^2}
  &= - \frac{(x^2 + 2)^2 \times  (2x-2) - (x^2 - 2x - 2)(4x^3 + 8x)}{(x^2 + 2)^4} \\
  &= - \frac{2x^5 - 6x^4 - 8x^3 - 8x^2 - 24x + 8}{(x^2 + 2)^4}.
\end{align*}

分母は常に正であるため、 numerator の符号を確認すればよい。

$x = 2.73$ を代入すると numerator は負であり、極大値 $y = 0.183$。

$x = -0.73$ を代入すると numerator は正であり、極小値 $y = -0.683$。

(4) 定数 $a$, $b$, $c$, $d$ が正であるとき、ある工場の週間出荷量 $P$ に応じて費用 $C$ は関係式 $C = aP + \frac{b}{c+P} + d$ で与えられる。最小の費用となる出荷量は何か？

\[
\dfrac{dC}{dP} = a - \frac{b}{(c+P)^2} = 0\quad \text{極大または極小の場合;}
\]
したがって $a = \frac{b}{(c+P)^2}$ で、$P = ±\sqrt{\frac{b}{a}} - c$。

出荷量は負にならないため、$P = +\sqrt{\frac{b}{a}} - c$。

\begin{align*}
\lintertext{\indent 現在}
\frac{d^2C}{dP^2} &= + \frac{b(2c + 2P)}{(c + P)^4},
\end{align*}
これはすべての $P$ の値に対して正であるため、$P = +\sqrt{\frac{b}{a}} - c$ は最小点を表す。

(5) $N$ 個の特定種類のランプで建物を明るく照らす総コスト每時間 $C$ は
\[
C = N\left(\frac{C_l}{t} + \frac{EPC_e}{1000}\right),
\]
ここで $E$ は商業効率（ワット/キャンドル）、
\begin{align*}
&\text{$P$ は各ランプのキャンドルパワー、} \\
&\text{$t$ は各ランプの平均寿命（時間）である、} \\
&\text{$C_l =$ 使用毎時間あたりの更新コスト（ Pence）、} \\
&\text{$C_e =$ 毎時間あたりのエネルギーコスト（1000ワットあたり）。}
\end{align*}

さらに、ランプの平均寿命とそのランプで運転される商業効率との関係は $t = mE^n$ という近似式で与えられる、ここで $m$ および $n$ はランプの種類により定数である。

商業的な効率を求め、照明の総コストが最小となるものを求めよ。
\begin{align*}
C &= N\left(\frac{C_l}{m} E^{-n} + \frac{PC_e}{1000} E\right), \\
\frac{dC}{dE} &= \frac{PC_e}{1000} - \frac{nC_l}{m} E^{-(n+1)} = 0
\end{align*}
これにより最大値または最小値が得られる。
\[
E^{n+1} = \frac{1000 \times  nC_l}{mPC_e}\quad \text{および}\quad
E = \sqrt[n+1]{\frac{1000 \times  nC_l}{mPC_e}}.
\]
これは明らかに最小値であり、
\[
\frac{d^2C}{dE^2} = (n + 1) \frac{nC_l}{m} E^{-(n+2)},
\]
が正であるため。

特定の種類の16燭光のランプにおいて、$C_l=17$ペンス、$C_e=5$ペンス；そして $m=10$ および $n=3.6$ と見つかった。
\[
E = \sqrt[4.6]{\frac{1000 \times  3.6 \times  17}{10 \times  16 \times  5}}
  = 2.6\text{ワット/燭光}.
\]

演習問題 X（数値例のグラフをプロットすることをお勧めします。）（解答は \Pageref{AnsEx:X} を参照してください。）

\begin{enumerate}
\item (1) $y = x^3 + x^2 - 10x + 8$ の最大値と最小値を求めよ。
\item (2) $y = \frac{b}{a}x - cx^2$ について、$\frac{dy}{dx}$ および $\frac{d^2y}{dx^2}$ を求めよ。また、$y$ の最大値または最小値をもたらす $x$ の値を求めよ。さらに、それが最大値か最小値であることを示せ。
\item (3) 曲線 $y = 1 - \frac{x^2}{2} + \frac{x^4}{24}$ における最大値と最小値の数を求める。また、曲線 $y = 1 - \frac{x^2}{2} + \frac{x^4}{24} - \frac{x^6}{720}$ の場合も同様に求めよ。
\item (4) $y=2x+1+\frac{5}{x^2}$ の最大値と最小値を求めよ。
\item (5) $y=\frac{3}{x^2+x+1}$ の最大値と最小値を求めよ。
\item (6) $y=\frac{5x}{2+x^2}$ の最大値と最小値を求めよ。
\item (7) $y=\frac{3x}{x^2-3} + \frac{x}{2} + 5$ の最大値と最小値を求めよ。
\item (8) 数字 $N$ を二つの部分に分割し、一つの部分の三倍の平方ともう一つの部分の二倍の平方の和が最小となるように求めよ。
\item (9) 電力発電機の効率 $u$ は出力 $x$ の異なる値で一般化された式で表される：
\[
u=\frac{x}{a+bx+cx^2};
\]
ここで、$a$ は主に鉄のエネルギー損失に依存し、$c$ は主に銅部品の抵抗に依存する定数である。出力が最大となる効率を表す式を求めよ。

\item[(10)] 某蒸気船の燃料消費量が式 $y = 0.3 + 0.001v^3$ で表されることが分かっているとする；ここで $y$ は時間ごとのトン数、$v$ は海里毎時という単位での速度である。賃金費、資本の利子、そしてその船の損耗費用は、時間ごとに1トンの燃料のコストに等しい。1000 海里的航海の総コストが最小になる速度は何であるか？また、燃料の価格が1トンにつき10シリングの場合、その最小コストはいくらとなるか？

\item[(11)] 次の関数の極値を求めよ：
\[
y = ±\frac{x}{6}\sqrt{x(10-x)}.
\]

\item[(12)] 次の関数の極値を求めよ：
\[
y= 4x^3 - x^2 - 2x + 1.
\]

\end{enumerate}



\chapter{その他の便利な巧妙法（部分分数展開・逆関数）}
\noindent\textit{\large Other Useful Dodges}\\[
2mm]
\noindent\textbf{\small 【章の論理とテーマ】複雑な有理関数の部分分数分解・逆関数微分法 dx/dy = 1/(dy/dx)}\\[4mm]
章 XIII その他の有用な技巧
節 部分分数.
まず最初に、部分分数をどのように得るか見てみましょう。分数を複数のより単純な分数に分割し、それらの和が元の分数と等しい場合、それぞれの単純な式について微分を行うことができます。そして、結果はいくつか（またはそれ以上の）単純な微分の合計となり、最終的な表現は当然同じですが、この技巧を使用せずに得られるものよりも労力を節約し、簡素化された形で得られます。
ではどのようにしてこのような結果を得るか見てみましょう。まず、二つの分数を足し合わせて結果の分数を作る作業から始めます。例えば、$\dfrac{1}{x+1}$ と $\dfrac{2}{x-1}$ の二つの分数を考えます。中学生でもこれらの分数を足し合わせて、その和が $\dfrac{3x+1}{x^2-1}$ であることを計算できます。同様に、三つ以上の分数も足し合わせることができます。
このプロセスは確かに逆方向にも適用可能です：つまり、最後の表現が与えられた場合、それは元の成分や部分分数にどのように分割できるかは確実です。しかし、私たちにはすべての場合でそのような分割方法を知る保証はありません。これを明らかにするために最初に簡単なケースを考慮します。ただし、以下の説明は「適切」な代数的分数（分子の次数が分母の次数より少ない分数）に対してのみ適用されることを忘れないでください。例えば、$\dfrac{1}{x+1}$ と $\dfrac{2}{x-1}$ のような分数です。最高次数 $x$ の指数が分子で分母よりも小さい場合に限ります。
もし $\dfrac{x^2+2}{x^2-1}$ のような分数を扱う必要がある場合、除算により簡素化できます。これは $\dfrac{x^2+2}{x^2-1} = 1 + \dfrac{3}{x^2-1}$ と等しいためで、$\dfrac{3}{x^2-1}$ は部分分数に分割できる適切な代数的分数です。この操作の詳細は後述します。
Case I. 二つ以上の分数を多次数の加算を行う場合、その分母に $x$ の項のみが含まれていて $x^2$, $x^3$，または $x$ の他の任意のべき乗は含まないとき、最終的な結果となる分数の分母は、それらの分数の分母の積と一致することが常に確認される。したがって、最終的な分数の分母を素因数分解することで、探している部分分数のすべての分母を見つけることができる。
$\dfrac{3x+1}{x^2-1}$ から $\dfrac{1}{x+1}$ と $\dfrac{2}{x-1}$ という成分に戻りたい場合を考えてみよう。これらの成分が分からなくても、次のように準備することができる：
\[
\dfrac{3x+1}{x^2-1} = \dfrac{3x+1}{(x+1)(x-1)} = \dfrac{}{x+1} + \dfrac{}{x-1},
\]
分子の場所はまだ空白であるが、それらを埋めることが可能である。部分分数間の符号は常に「プラス」で仮定できる。なぜなら、「マイナス」であれば、対応する分子は負の値となるだけである。また、これらの部分分数は「正規」な分数であり、分子は $x$ の項を含まない単なる数であるため、任意の文字 $A$, $B$, $C \dots$ で表すことができる。したがって、この場合、次の式を得る：
\[
\dfrac{3x+1}{x^2-1} = \dfrac{A}{x+1} + \dfrac{B}{x-1}.
\]

これらの二つの部分分数を加算すると、$\dfrac{A(x-1)+B(x+1)}{(x+1)(x-1)}$ となる。そしてこれは $\dfrac{3x+1}{(x+1)(x-1)}$ と等しいことが必要である。両者の分母が一致しているため、分子も等しくなり、次のように得られる：
\[
3x + 1 = A(x-1) + B(x + 1).
\]

今、これは2つの未知数を持つ方程式です。$A$ と $B$ を求めるためにもう1つの方程式が必要のように思われます。

しかし、この困難を解く別の方法があります。その方程式はすべての $x$ の値に対して成り立つべきであり、したがって、$x-1$ および $x+1$ がそれぞれ0になるような $x$ の値（すなわち $x=1$ と $x=-1$）に対して成り立つ必要があります。$x=1$ に代入すると、$4 = (A \times 0) + (B \times 2)$ となり、$B=2$ となります。また、$x=-1$ に代入すると、$-2 = (A \times -2) + (B \times 0)$ となり、$A=1$ となります。

部分分数の $A$ と $B$ をこれらの新しい値で置き換えると、それぞれ $\frac{1}{x+1}$ および $\frac{2}{x-1}$ になります。これで問題は解決しました。

さらに別の例として、$\frac{4x^2 + 2x - 14}{x^3 + 3x^2 - x - 3}$ の分数を考えます。分母は $x=1$ のとき0になるため、$x-1$ が因子であり、他の因子は $x^2 + 4x + 3$ であることが明らかです。これも $(x+1)(x+3)$ に分解できます。

したがって、分数を次のように書き換えることができます：
\[
\frac{4x^2 + 2x - 14}{x^3 + 3x^2 - x - 3} = \frac{A}{x+1} + \frac{B}{x-1} + \frac{C}{x+3},
\]
3つの部分分数を用意します。

先に述べた方法で計算すると、
\[
4x^2 + 2x - 14 = A(x-1)(x+3) + B(x+1)(x+3) + C(x+1)(x-1).
\]

$x=1$ に代入すると：
\[
-8 = (A \times 0) + B(2 \times 4) + (C \times 0); \quad \text{すなわち } B = -1.
\]

また、$x=-1$ に代入すると：
\[
-12 = A(-2 \times 2) + (B \times 0) + (C \times 0); \quad \text{従って } A = 3.
\]

さらに、$x=-3$ に代入すると：
\[
16 = (A \times 0) + (B \times 0) + C(-2 \times -4); \quad \text{従って } C = 2.
\]

したがって、部分分数は次のようになります：
\[
\frac{3}{x+1} - \frac{1}{x-1} + \frac{2}{x+3},
\]
これは $x$ に関する微分よりも元の複雑な表現から導き出されたものよりもずっと簡単です。

\par\medskip\noindent\textbf{Case II.} 分母の一部の因子が $x^2$ の項を含み、適切な因数分解が難しい場合、対応する分子には $x$ の項と単純な数が含まれることがあります。そのため、この未知の分子は $A$ ではなく $Ax + B$ で表す必要があります。その後の計算は以前と同じように行われます。
\begin{gather*}
\lintertext{\rlap{\indent 例えば試してみる：}}
\frac{-x^2 - 3}{(x^2+1)(x+1)}. \\
\frac{-x^2 - 3}{(x^2+1)(x+1)} = \frac{Ax+B}{x^2+1} + \frac{C}{x+1};\\
-x^2 - 3 = (Ax + B)(x+1) + C(x^2+1).
\end{gather*}

$x = -1$ を代入すると $-4 = C \times  2$; そして $C = -2$;
\begin{align*}
\lintertext{従って}
-x^2 - 3 &= (Ax + B)(x + 1) - 2x^2 - 2; \\
\lintertext{そして}
x^2 - 1 &= Ax(x+1) + B(x+1).
\end{align*}

$x = 0$ を代入すると $-1 = B$;
従って
\begin{gather*}
x^2 - 1 = Ax(x + 1) - x - 1;\quad \text{または } x^2 + x = Ax(x+1); \\
\lintertext{そして}
x+1 = A(x+1),
\end{gather*}

したがって $A=1$ であり、部分分数は
\[
\frac{x-1}{x^2+1} - \frac{2}{x+1}.
\]

別の例として、分数
\[
\frac{x^3-2}{(x^2+1)(x^2+2)}
\]
を取る。

得られるのは
\begin{align*}
\frac{x^3-2}{(x^2+1)(x^2+2)}
  &= \frac{Ax+B}{x^2+1} + \frac{Cx+D}{x^2+2}\\
  &= \frac{(Ax+B)(x^2+2)+(Cx+D)(x^2+1)}{(x^2+1)(x^2+2)}.
\end{align*}

この場合、$A$, $B$, $C$, $D$ の決定は簡単ではありません。次のように進めるのが簡単です：与えられた分数と部分分数を加えた分数が等しいこと、そして分母が同一であることを考慮すると、分子も同一でなければなりません。このような場合、およびここでのような代数的表現に対して、\emph{同じ次数の $x$ の係数は等しく且つ符号も同じ}。

従って
\begin{align*}
x^3-2
  &= (Ax+B)(x^2+2) + (Cx+D)(x^2+1) \\
  &= (A+C)x^3 + (B+D)x^2 + (2A+C)x + 2B+D,
\end{align*}
$1=A+C$; $0=B+D$（左辺の $x^2$ の係数がゼロである）; $0=2A+C$; そして $-2=2B+D$。これらは4つの方程式であり、容易に $A=-1$; $B=-2$; $C=2$; $D=0$ を得ます；したがって部分分数は $\dfrac{2(x+1)}{x^2+2} - \dfrac{x+2}{x^2+1}$ です。

この方法は常に使用できますが、$x$ の因子のみの場合には最初に示した方法の方が速いことが見受けられます。

\par\medskip\noindent\textbf{Case\textasciitilde{}III\@.} 分母の因子の中にいくつかがべき乗になっている場合、その因子の各べき乗に対する部分分数の可能性があることを考慮する必要があります。例えば、分数 $\dfrac{3x^2-2x+1}{(x+1)^2(x-2)}$ を分解するときには、$x+1$ の一次因子も含む可能性があるため、$(x+1)$, $(x+1)^2$, および $(x-2)$ が分母となる部分分数を考慮しなければなりません。

しかし、分母が $(x+1)^2$ である分数の分子に $x$ の項があると仮定して、その分子を $Ax+B$ で表す必要があると考えるかもしれません。したがって、

\[
\frac{3x^2 - 2x + 1}{(x+1)^2(x-2)}
  = \frac{Ax+B}{(x+1)^2} + \frac{C}{x+1} + \frac{D}{x-2}
\]

と書く必要があります。しかし、$A$, $B$, $C$, および $D$ を求めようとすると、4つの未知数に対して3つの関係式しかないので失敗します。実際、

\[
\frac{3x^2 - 2x + 1}{(x+1)^2(x-2)}
  = \frac{x-1}{(x+1)^2} + \frac{1}{x+1} + \frac{1}{x-2}
\]

が成り立ちます。

しかし、

\[
\frac{3x^2 - 2x + 1}{(x+1)^2(x-2)}
  = \frac{A}{(x+1)^2} + \frac{B}{x+1} + \frac{C}{x-2}
\]

と書くと、

\[
3x^2 - 2x + 1 = A(x-2) + B(x+1)(x-2) + C(x+1)^2
\]

となり、$x=2$ のとき $C=1$ が得られます。$C$ の値を代入し、移項して類似項をまとめ、$x-2$ で割ると、

\[
-2x = A + B(x+1)
\]

となり、$x=-1$ のとき $A=-2$ が得られます。$A$ の値を代入すると、

\[
2x = -2 + B(x+1)
\]

となり、$B=2$ が得られます。したがって、部分分数は

\[
\frac{2}{x+1} - \frac{2}{(x+1)^2} + \frac{1}{x-2}
\]

となります。これは先に述べた $\dfrac{1}{x+1} + \dfrac{x-1}{(x+1)^2} + \dfrac{1}{x-2}$ とは異なります。$\dfrac{x-1}{(x+1)^2}$ はさらに分解でき、$\dfrac{1}{x+1} - \dfrac{2}{(x+1)^2}$ に分割されるため、

\[
\frac{1}{x+1} + \frac{1}{x+1} - \frac{2}{(x+1)^2} + \frac{1}{x-2}
  = \frac{2}{x+1} - \frac{2}{(x+1)^2} + \frac{1}{x-2}
\]

となり、これが得られた部分分数と一致します。

このように、各分子に一つの数値項を考慮すれば十分であり、最終的な部分分数が常に得られます。

When thereは有理分数を展開する際に、分母に$x^2$の因子のべき乗がある場合、対応する分子は$Ax+B$の形でなければなりません。例えば、

\[
\frac{3x-1}{(2x^2-1)^2(x+1)} = \frac{Ax+B}{(2x^2-1)^2} + \frac{Cx+D}{2x^2-1} + \frac{E}{x+1},
\]

これは、

\[
3x - 1 = (Ax + B)(x + 1) + (Cx + D)(x + 1)(2x^2 - 1) + E(2x^2 - 1)^2.
\]

$x = -1$のとき、$E = -4$となります。置き換え、移項、類似項をまとめ、$x+1$で割ると、

\[
16x^3 - 16x^2 + 3 = 2Cx^3 + 2Dx^2 + x(A - C) + (B - D).
\]

したがって、$2C = 16$から$C = 8$；$2D = -16$から$D = -8$；$A - C = 0$から$A - 8 = 0$で$A = 8$；最後に、$B - D = 3$から$B = -5$となります。したがって、有理分数は、

\[
\frac{8x-5}{(2x^2-1)^2} + \frac{8(x-1)}{2x^2-1} - \frac{4}{x+1}
\]

となります。

結果を確認するためには、与えられた式と得られた有理分数に$x$の単一値、例えば$+1$を代入するのが最も簡単な方法です。

分母が单一因子のべき乗のみを含む場合、非常に迅速な方法は次の通りです。例えば、

\[
\frac{4x + 1}{(x + 1)^3}
\]

を考えます。ここで$x + 1 = z$と置きます；すると$x = z - 1$となります。

代入すると、

\[
\frac{4(z - 1) + 1}{z^3} = \frac{4z - 3}{z^3} = \frac{4}{z^2} - \frac{3}{z^3}.
\]

したがって、有理分数は、

\[
\frac{4}{(x+1)^2} - \frac{3}{(x+1)^3}
\]

となります。

微分への応用。$y = \dfrac{5-4x}{6x^2 + 7x - 3}$を微分する必要があるとします；すると、

\[
\frac{dy}{dx} = -\frac{(6x^2+7x-3) \times  4 + (5 - 4x)(12x + 7)}{(6x^2 + 7x - 3)^2}
\]

となり、

\[
\frac{dy}{dx} = \frac{24x^2 - 60x - 23}{(6x^2 + 7x - 3)^2}.
\]

与えられた式を

\[
\frac{1}{3x-1} - \frac{2}{2x+3}
\]

に分割すると、

\[
\frac{dy}{dx} = -\frac{3}{(3x-1)^2} + \frac{4}{(2x+3)^2}
\]

となり、これは上記の結果を有理分数に分割したものです。しかし、微分後に分割するとより複雑になります。このような式の積分を考えるとき、有理分数への分割は貴重な補助手段となります（参照：\pageref{partfracs})。

問題集 XI (解答は \pageref[page]{AnsEx:XI} を参照してください)。

次の式を有理分数に分解せよ：

1. $\dfrac{3x + 5}{(x - 3)(x + 4)}$。
2. $\dfrac{3x - 4}{(x - 1)(x - 2)}$。

\begin{enumerate}
\item (3) \(\dfrac{3x + 5}{x^2 + x - 12}\).
\item (4) \(\dfrac{x + 1}{x^2 - 7x + 12}\).

\item (5) \(\dfrac{x - 8}{(2x + 3)(3x - 2)}\).
\item (6) \(\dfrac{x^2 - 13x + 26}{(x - 2)(x - 3)(x - 4)}\).

\item (7) \(\dfrac{x^2 - 3x + 1}{(x - 1)(x + 2)(x - 3)}\).

\item (8) \(\dfrac{5x^2 + 7x + 1}{(2x + 1)(3x - 2)(3x + 1)}\).

\item (9) \(\dfrac{x^2}{x^3 - 1}\).
\item (10) \(\dfrac{x^4 + 1}{x^3 + 1}\).

\item (11) \(\dfrac{5x^2 + 6x + 4}{(x + 1)(x^2 + x + 1)}\).
\item (12) \(\dfrac{x}{(x - 1)(x - 2)^2}\).

\item (13) \(\dfrac{x}{(x^2 - 1)(x + 1)}\).
\item (14) \(\dfrac{x + 3}{(x + 2)^2(x - 1)}\).

\item (15) \(\dfrac{3x^2 + 2x + 1}{(x + 2)(x^2 + x + 1)^2}\).
\item (16) \(\dfrac{5x^2 + 8x - 12}{(x + 4)^3}\).

\item (17) \(\dfrac{7x^2 + 9x - 1}{(3x - 2)^4}\).
\item (18) \(\dfrac{x^2}{(x^3 - 8)(x - 2)}\).

\section*{逆関数の微分。}

関数 \(y = 3x\) を考慮します（参照：\Pageref{function}）。この関数は形 \(x = \dfrac{y}{3}\) で表すことができます。後者の形式は元々与えられた関数の逆関数と呼ばれます。

もし \(y = 3x\) のとき、\(\dfrac{dy}{dx} = 3\)；\(x = \dfrac{y}{3}\) のとき、\(\dfrac{dx}{dy} = \dfrac{1}{3}\) であり、これらから
\[
\dfrac{dy}{dx} = \dfrac{1}{\ \dfrac{dx}{dy}\ }\quad \text{または}\quad
\dfrac{dy}{dx} \times  \dfrac{dx}{dy} = 1.
\]

次に \(y = 4x^2\) を考慮します。このとき、\(\dfrac{dy}{dx} = 8x\)；逆関数は
\[
x = \dfrac{y^{\frac{1}{2}}}{2},\quad \text{および}\quad
\dfrac{dx}{dy} = \dfrac{1}{4\sqrt{y}} = \dfrac{1}{4 \times  2x} = \dfrac{1}{8x}.
\]

ここでまた
\[
\dfrac{dy}{dx}\times \dfrac{dx}{dy} = 1.
\]

すべての関数について、逆関数に変形できるものに対して、常に書き表すことができます
\[
\dfrac{dy}{dx} \times  \dfrac{dx}{dy} = 1\quad \text{または}\quad
\dfrac{dy}{dx} = \dfrac{1}{\ \dfrac{dx}{dy}\ }.
\]

したがって、与えられた関数について、逆関数の微分が計算しやすい場合、その逆関数を微分して、逆関数の微分係数の逆数は与えられた関数の微分係数となります。

例えば、$y=\sqrt{\dfrac{3}{x}-1}$ の微分を考えます。この問題は、$u=\dfrac{3}{x}-1$ とおくことで $\dfrac{dy}{du}$ と $\dfrac{du}{dx}$ を求めることで解くことができます。これにより、

\[
\frac{dy}{dx} = -\frac{3}{2x^2\sqrt{\dfrac{3}{x}-1}}
\]

を得ます。

もしこの方法を忘れてしまったり、別の方法で導関数を求めたい場合や、通常の方法を使用できない理由がある場合は、以下の手順で進めることが可能です：逆関数は $x=\dfrac{3}{1+y^2}$ です。
\[
\frac{dx}{dy} = -\frac{6y}{(1+y^2)^2}
\]
従って
\[
\frac{dy}{dx} = \frac{1}{\frac{dx}{dy}} = -\frac{(1+y^2)^2}{6y} = -\frac{\left(1+\dfrac{3}{x}-1\right)^2}{6\sqrt{\dfrac{3}{x}-1}} = -\frac{3}{2x^2\sqrt{\dfrac{3}{x}-1}}
\]

別の例として、$y=\dfrac{1}{\sqrt[3]{\theta+5}}$ を考えます。

逆関数は $\theta=\dfrac{1}{y^3}-5$ または $\theta=y^{-3}-5$ であり、

\[
\frac{d\theta}{dy} = -3y^{-4} = -3(\theta + 5)^{\frac{4}{3}}
\]

従って、$\dfrac{dy}{dx} = -\dfrac{1}{3(\theta+5)^{\frac{4}{3}}}$ となります。

この方法は後々非常に有用であることがわかります。そのため、Exercise I.（ページ参照：\pageref{Ex:I}）の問題 5, 6, 7；Example （ページ参照：\pageref{ExNo1}）の問題 1, 2, 4；そして Exercise VI.（ページ参照：\pageref{Ex:VI}）の問題 1, 2, 3 および 4 の結果を確認するために、この方法を使って結果を検証することをお勧めします。

この章と前章から、微分積分学は多くの面で芸術よりも科学であることを理解するでしょう。これは他のすべての芸術と同じく練習によって習得されるものです。したがって、多くの例題を解き、自分で新たな問題を作成して解けるか試すことが重要です。これにより、様々な技巧が使用に慣れることになります。

あなたはこの章と前章から、微分積分学は多くの面で芸術よりも科学であることを理解するでしょう。これは他のすべての芸術と同じく練習によって習得されるものです。したがって、多くの例題を解き、自分で新たな問題を作成して解けるか試すことが重要です。これにより、様々な技巧が使用に慣れることになります。
\end{enumerate}



\chapter{真の複利計算と有機的成長の法則}
\noindent\textit{\large On true Compound Interest and the Law of Organic Growth}\\[
2mm]
\noindent\textbf{\small 【章の論理とテーマ】自然対数の底 e・指数関数 e^x の微分・放射性崩壊と人口増加モデル}\\[4mm]
第XIV章 真の複利と有機的な成長法則
まず、ある量が一定時間における増加量がその量自身の大きさに比例するように増加することを考えましょう。これは固定利率で金利を計算する過程に似ています；つまり、資本金が大きいほど、その期間における利息も大きくなります。
しかし、計算において二つのケースを明確に区別しなければなりません。それは算術書で「単利」または「複利」と呼ばれる方法による計算の違いです。前者では資本金は固定ですが、後者では利息が資本金に加算され、資本金は逐次的な増加分によって増大します。
1. 単利の場合。具体的な例を考えてみましょう。最初の資本金が$p$100で、年利率が10%であるとします。すると、資本金の所有者には毎年$p$10の増加が生じます。彼が毎年利息を受け取り、それを靴下に保管したり、金庫に入れて保管したりする場合、10年間続けば、$p$10の10回分の増加を受けることになります。つまり$p$100で、元々の$p$100と合わせて合計$p$200となります。彼の財産は10年で倍になりました。もし年利率が5%であれば、財産を倍にするには20年間待つ必要があります。2%であれば50年です。年利が資本金の$\frac{1}{n}$である場合、彼は財産を倍にするために$n$年間保管する必要があります。
または、元の資本金を$y$とし、年利が$\frac{y}{n}$である場合、$n$年後には彼の財産は以下のようになります。
\[
y + n \cdot \frac{y}{n} = 2y.
\]

以前と同じように、所有者は$p$100の元金から始め、年率10%の利子を得ます。しかし、利息を貯蓄せず、毎年利息を元金に加えることにします。そうすると、元金は年々増えていきます。

最初の1年後には、元金は$p$110になります；2年目（依然として10%）では、この$p$11が$p$11の利子を生みます。彼は3年目に$p$121から始め、その$p$121に対する利子は$p$12.2sとなります；したがって、4年目に$p$133.2sから始めます。これを計算してみると、10年後には総元金が$p$259.7s6dになります。

実際、毎年終了時にそれぞれのポンドは$p$0.1を獲得しており、したがって、常にこの額が加算されると、各年は元金を$\frac{11}{10}$倍にします。10年間続けば、これは元金を2.59374倍にします。

これを記号で表すと、$y_0$を元金、$\frac{1}{n}$を各$n$回の操作ごとに加算される分数、$y_n$を$n$回目の操作終了時の元金の値とします。すると、

\[
y_n = y_0 \left(1 + \frac{1}{n}\right)^n.
\]

しかし、年1回の複利計算は実際には公平ではありません。最初の1年間でも$p$100は成長すべきです。半年後には少なくとも$p$105になっているべきで、2半期目の利息を$p$105に基づいて計算すればより公正です。これは年率5%とみなすことができ、20回の操作があり、各回元金が$\frac{21}{20}$倍になります。

この方法で計算すると、10年後には元金は$p$265.6s7dになります。なぜなら、

\[
(1 + \frac{1}{20})^{20} = 2.653
\]

だからです。

しかし、それでもなお過程は完全には公正ではありません。最初の月が終了するまでに利息が発生しますし、半年ごとの計算では資本が6ヶ月間固定されていると仮定されます。一年を10個の部分に分割し、各部分で1%の利子を計算するとしましょう。これにより100回の操作が10年間にわたって行われます；または
\[
y_n = $p$100 \left( 1 + \frac{1}{100} \right)^{100};
\]
これは
\DPtypo{$p$$270$.\textasciitilde{}$8$\textit{s}.}
       {$p$$270$.\textasciitilde{}$9$\textit{s}.\textasciitilde{}$7\frac{1}{2}$\textit{d}.}

この計算もまだ完了ではありません。10年を1000個の部分に分割し、各部分が1/100年となるようにします。各部分で0.1%の利子を計算すると
\[
y_n = $p$100 \left( 1 + \frac{1}{1000} \right)^{1000};
\]
これは
\DPtypo{$p$$271$.\textasciitilde{}$14$\textit{s}.\textasciitilde{}$2\frac{1}{2}$\textit{d}.}
       {$p$$271$.\textasciitilde{}$13$\textit{s}.\textasciitilde{}$10$\textit{d}.}

さらに細かく分割し、10年を10,000個の部分に分割します。各部分が1/1000年で、利子は1%の1/10となります。すると
\[
y_n = $p$100 \left( 1 + \frac{1}{10,000} \right)^{10,000};
\]
これは
\DPtypo{$p$$271$.\textasciitilde{}$16$\textit{s}.\textasciitilde{}$4$\textit{d}.}
       {$p$$271$.\textasciitilde{}$16$\textit{s}.\textasciitilde{}$3\frac{1}{2}$\textit{d}.}

最後に、我々が探しているものは実際には式 $\left(1 + \frac{1}{n}\right)^n$ の極限値であり、これは2より大きいことがわかります。そして $n$ を大きくするにつれて、この値は特定の極限値に近づきます。どの程度でも $n$ を大きくしても、この式の値は徐々に
\[
2.71828\ldots
\]
という数字に近づきます。これは決して忘れることのない数です。

これらのことを幾何学的な例で説明しましょう。
図36において、$OP$は元の値を表します。$OT$は値が成長する全体の時間です。これは10つの期間に分割され、各期間には等しいステップアップがあります。ここで$\frac{dy}{dx}$は定数であり、もし各ステップアップが元の$OP$の$\frac{1}{10}$であれば、10回のそのようなステップアップにより高さは2倍になります。もし20のステップを取れば、それぞれが示される高さの半分である場合でも、最終的には高さは2倍になります。または$n$回のそのようなステップがあり、各ステップが元の高さ$OP$の$\frac{1}{n}$であれば、高さは2倍になります。これは単利の場合です。ここでは$1$が$2$になるまで成長します。

図37において、対応する幾何級数の例を示しています。各連続的な縦線は前段の$\left(1 + \frac{1}{n}\right)$倍、つまり$n+1$分の$n$倍の高さになります。ステップアップは等しくありませんが、これは各ステップアップが曲線のその部分での縦線の$\frac{1}{n}$であるからです。もし実際には10のステップを取って、$\left(1 + \frac{1}{10}\right)$を乗数として用いれば、最終的な合計は$(1 + \frac{1}{10})^{10}$または$2.594$倍になります。しかし、十分に大きな$n$（そして十分に小さな$\frac{1}{n}$）を選べば、最終的な値$\left(1 + \frac{1}{n}\right)^n$は$2.71828$になります。

図37：縦線の高さが$(1 + \frac{1}{n})^n$倍になる

ε．この神秘的な数値$2.7182818$等、数学家たちはギリシャ文字の$\epsilon$（読み方は「イプシロン」）というシンボルで表しています。学校の生徒は$\pi$（読み方は「パイ」）が$3.141592$等を表すことを知っていますが、どれだけの人が$\epsilon$が$2.71828$であることを知っているでしょうか？しかし、それは$\pi$よりももっと重要な数値です！

では、$\epsilon$とは何なのでしょうか？

もし$1$が単利で$2$になるまで成長させたとします。同じ名目金利で同じ期間を考慮に入れると、真の複利ではなく単利で$1$が成長した場合、その最終的な値は$\epsilon$になります。

この瞬間ごとにその瞬間の大きさに比例して成長する過程を、一部の人々は「対数率での成長」と呼ぶ。単位対数率の成長は、単位時間で1が2.718281になるような成長速度である。また、これは「有機的な成長率」とも呼ばれるかもしれない：なぜなら、特定の状況下では生物体の増加量がその生物体自身の大きさに比例することから、有機的成長の特性として認識されるからだ。

もし100パーセントを率の単位とし、固定された期間を時間の単位とするなら、単位率で単位時間にarithmetically（算術的に）成長させた結果は2となり、同じ時間にlogarithmically（対数的に）成長させた結果は2.71828…となる。

εについて少し補足する。我々は、式$\left(1 + \frac{1}{n}\right)^n$が$n$が無限大になるときの値を知る必要があることを見た。算術的に、以下に多くの値（誰でも通常の対数表を使って計算できる）を列挙する。これらは$n = 2$；$n = 5$；$n = 10$；そしてそれ以降、$n = 10,000$まで順に求めたものである。
\begin{alignat*}{2}
&(1 + \frac{1}{2})^2             &&= 2.25.    \\
&(1 + \frac{1}{5})^5             &&= 2.488.   \\
&(1 + \frac{1}{10})^{10}         &&= 2.594.   \\
&(1 + \frac{1}{20})^{20}         &&= 2.653.   \\
&(1 + \frac{1}{100})^{100}       &&= 2.705.   \\
&(1 + \frac{1}{1000})^{1000}     &&= 2.7169.  \\
&(1 + \frac{1}{10,000})^{10,000} &&= 2.7181.
\end{alignat*}

しかし、この非常に重要な数値を別の方法で計算する価値がある。

したがって、二項定理を利用して、式$\left(1 + \frac{1}{n}\right)^n$をその有名な方法で展開することにしよう。

二項定理（）は、次の規則を提供します：
\begin{align*}
(a + b)^n &= a^n + n \dfrac{a^{n-1} b}{1!} + n(n - 1) \dfrac{a^{n-2} b^2}{2!} \\
  & = a^n + n \dfrac{a^{n-1} b}{1!} + n(n - 1) \dfrac{a^{n-2} b^2}{2!} \\
  & \phantom{= a^n\ } + n(n -1)(n - 2) \dfrac{a^{n-3} b^3}{3!} + \text{等}.
\end{align*}
$a = 1$ および $b = \dfrac{1}{n}$ を代入すると、
\[
\left(1 + \dfrac{1}{n}\right)^n
  = 1 + 1 + \dfrac{1}{2!} \left(\dfrac{n - 1}{n}\right) + \dfrac{1}{3!} \dfrac{(n - 1)(n - 2)}{n^2} \\
  = 1 + 1 + \dfrac{1}{2!} \dfrac{n - 1}{n} + \dfrac{1}{3!} \dfrac{(n - 1)(n - 2)}{n^2} \\
  &\phantom{= 1 + 1\ } + \dfrac{1}{4!} \dfrac{(n - 1)(n - 2)(n - 3)}{n^3} + \text{等}.
\]

今、$n$ を無限大に仮定し、たとえば十億、または十億の十億とするとき、$n - 1$, $n - 2$, $n - 3$ 等はすべて $n$ に近似して等しくなります。すると級数は次のようになります：
\[
\epsilon = 1 + 1 + \dfrac{1}{2!} + \dfrac{1}{3!} + \dfrac{1}{4!} + \text{等}.\ldots
\]

この収束が速い級数を任意の項まで計算することで、必要な精度に達することができます。十項までの計算は次の通りです：

\begin{tabular}{@{}r<{\qquad}@{}l@{}}
                \& $1.000000$ \\
除算により $1$ \& $1.000000$ \\
除算により $2$ \& $0.500000$ \\
除算により $3$ \& $0.166667$ \\
除算により $4$ \& $0.041667$ \\
除算により $5$ \& $0.008333$ \\
除算により $6$ \& $0.001389$ \\
除算により $7$ \& $0.000198$ \\
除算により $8$ \& $0.000025$ \\
除算により $9$ \& $0.000002$ \\
\cline{2-2}
\multicolumn{2}{r@{}}{合計\quad $2.718281$} \\
\cline{2-2}
\end{tabular}

$\epsilon$ は $1$ と無比し、$\pi$ に似て非周期的な不尽小数です。

指数級数について。さらに別の級数が必要になります。

Let us, 再び二項定理を利用し，式 $\left(1 + \dfrac{1}{n}\right)^{nx}$ を展開します。$n$ を無限大にするとこれは $\epsilon^x$ と等しくなります。
\begin{align*}
\epsilon^x
  &= 1^{nx} + nx \frac{1^{nx-1} \left(\dfrac{1}{n}\right)}{1!}
            + nx(nx - 1) \frac{1^{nx - 2} \left(\dfrac{1}{n}\right)^2}{2!} \\
  & = 1 + x + \frac{1}{2!} · \frac{n^2x^2 - nx}{n^2}
    + \frac{1}{3!} · \frac{n^3x^3 - 3n^2x^2 + 2nx}{n^3} + \text{等}.
  &= 1 + x + \frac{x^2 -\dfrac{x}{n}}{2!}
    + \frac{x^3 - \dfrac{3x^2}{n} + \dfrac{2x}{n^2}}{3!} + \text{等}.
\end{align*}

しかし，$n$ を無限大にすると，これは次のように単純化されます：
\[
\epsilon^x
  = 1 + x + \frac{x^2}{2!} + \frac{x^3}{3!} + \frac{x^4}{4!} + \text{等}\dots
\]

この級数は，指数級数と呼ばれます。

$\epsilon$ が重要視される大いなる理由は，$\epsilon^x$ が他の $x$ の任意の関数とは異なる性質を持つことからです。すなわち，その微分すると値が変わらない，または言い換えると，その導函数が自身と等しいという性質を持っています。これは次のように $x$ について微分することで瞬時に確認できます：
\begin{align*}
\frac{d(\epsilon^x)}{dx}
  &= 0 + 1 + \frac{2x}{1 · 2} + \frac{3x^2}{1 · 2 · 3} + \frac{4x^3}{1 · 2 · 3 · 4} \\
& = 1 + x + \frac{x^2}{1 · 2} + \frac{x^3}{1 · 2 · 3} + \frac{x^4}{1 · 2 · 3 · 4} + \text{等}.
\end{align*}
これは元の級数と完全に同じです。

今、別の方法で進むことも可能でした。つまり、$x$ の関数を求めることで、その微分係数が自分自身と一致するものを見つけ出すことから始めることができます。あるいは、$x$ のべき乗のみを含む表現において、微分によって変化しないものは存在しますか？従って、一般に以下の形を仮定します：
\begin{align*}
y &= A + Bx + Cx^2 + Dx^3 + Ex^4 + \cdots, \\
\intertext{(ここで係数 $A$, $B$, $C$ などは決定しなければなりません)そしてそれを微分します。}
\dfrac{dy}{dx} &= B + 2Cx + 3Dx^2 + 4Ex^3 + \cdots
\end{align*}

この新しい表現が本当に元のものと一致するためには、明らかに $A$ は $B$ に等しいことが必要で、$C = \dfrac{B}{2} = \dfrac{A}{1·2}$, $D = \dfrac{C}{3} = \dfrac{A}{1·2·3}$, $E = \dfrac{D}{4} = \dfrac{A}{1·2·3·4}$ などと続きます。

従って、変化の法則は以下のようになります：
\[
y = A\left(1 + \dfrac{x}{1} + \dfrac{x^2}{1·2} + \dfrac{x^3}{1·2·3} + \dfrac{x^4}{1·2·3·4} + \cdots\right).
\]

$A = 1$ と仮定してさらに単純化すると、
\[
y = 1 + \dfrac{x}{1} + \dfrac{x^2}{1·2} + \dfrac{x^3}{1·2·3} + \dfrac{x^4}{1·2·3·4} + \cdots
\]

これを何回微分しても同じ級数が得られます。

$A = 1$ の特別な場合を評価すると、
\begin{align*}
x &= 1 \quad & y &= 2.718281 \text{ etc.};    & \text{すなわち, } y &= e;   \\
x &= 2 \quad & y &=(2.718281 \text{ etc.})^2; & \text{すなわち, } y &= e^2; \\
x &= 3 \quad & y &=(2.718281 \text{ etc.})^3; & \text{すなわち, } y &= e^3;
\end{align*}
従って
\[
x = x \quad y = (2.718281 \text{ etc}.)^x;\quad \text{すなわち, } y=e^x,
\]
これで最終的に
\[
e^x = 1 + \dfrac{x}{1} + \dfrac{x^2}{1·2} + \dfrac{x^3}{1· 2· 3} + \dfrac{x^4}{1· 2· 3· 4} + \cdots
\]

[\textsc{注}.---\textit{指数の読み方}. 手元に教師がいない人にとって、$\epsilon^x$ を「エプシロンのエクス乗」と読むことが有用であるかもしれません。または、一部の人々はそれを「指数エックス」と読むこともあります。したがって $\epsilon^{pt}$ は「エプシロンのピーティース乗」または「指数ピー ティー」で読みます。類似する表現をいくつか挙げると：$\epsilon^{-2}$ は「エプシロンのマイナス二乗」または「指数マイナス二」で読み、$\epsilon^{-ax}$ は「エプシロンのマイナスエックスティース乗」または「指数マイナスエックスティー」で読みます。]

もちろん $\epsilon^y$ は $y$ に関する微分では変化しないことがわかります。また $\epsilon^{ax}$ は $(\epsilon^a)^x$ と等しいので、$x$ に関する微分では $a\epsilon^{ax}$ になります。ここで $a$ は定数です。

\subsection{自然対数やネイピアの対数}
$\epsilon$ が重要であるもう一つの理由は、ネイピア（対数の発明者）によって対数システムの基礎とされたことです。もし $y = \epsilon^x$ のとき、$x$ は $y$ の $\epsilon$ を底とする対数です。つまり
\begin{align*}
                  y &= \epsilon^x, \\
\intertext{ならば} x &= \log_\epsilon y.
\end{align*}

図 \ref{fig:38} と \ref{fig:39} に示す曲線はこれらの式を表しています。

計算された点は以下の通りです：
\begin{align*}
\text{図 \ref{fig:38}} &:
\begin{array}{|c||*{5}{c|}}
\hline
x & 0 & 0.5 & 1 & 1.5 & 2 \\
\hline
y & 1 & 1.65 & 2.71 & 4.50 & 7.39 \\
\hline
\end{array}
\\
\text{図 \ref{fig:39}} &:
\begin{array}{|c||*{5}{c|}}
\hline
y & 1 & 2 & 3 & 4 & 8 \\
\hline
x & 0 & 0.69 & 1.10 & 1.39 & 2.08 \\
\hline
\end{array}
\end{align*}

%[** TN: The graphs below are reversed in the original; exchanged.]
図 \ref{fig:158a} と \ref{fig:158b} は元の順番が逆になっていますので、交換しました。

計算結果から明らかに、描画する点は異なるかもしれませんが、結果自体は同じです。これらの2つの式は実質的に同じ意味を持っています。

多くの人が通常の対数を使用し、これは基底が10であるのに対し、自然対数は基底が$\epsilon$であるため、自然対数に不慣れな場合があります。それらについて少し言及することも価値があるかもしれません。通常の対数の規則として、対数を加えると積の対数になるという法則はまだ成り立ちます；すなわち
\[
\log_\epsilon a + \log_\epsilon b = \log_\epsilon ab.
\]
また、指数法則も依然として有効です；
\[
n \times  \log_\epsilon a = \log_\epsilon a^n.
\]
しかし、10が基底でなくなったため、100や1000を単に指標に2または3を加えることでかけることはできません。自然対数を通常の対数に変換するには、それを0.4343で掛ければよいです；すなわち
\begin{align*}
\log_{10} x &= 0.4343 \times  \log_{\epsilon} x, \\
\text{逆に} & \\
\log_{\epsilon} x &= 2.3026 \times  \log_{10} x.
\end{align*}

自然対数（または双曲線対数）の有用な表は次の通りです：

\begin{table}[hp]
\centering
\textsc{A Useful Table of ``Naperian Logarithms''} \\
(Also called Natural Logarithms or Hyperbolic Logarithms)
\[
\begin{array}{c<{\quad}|>{\ }c>{\qquad\qquad}cr<{\quad}|>{\ }c}
\multicolumn{1}{c}{\text{\footnotesize 数}} &
\multicolumn{1}{c}{\footnotesize\log_{\epsilon}} &&
\multicolumn{1}{c}{\text{\footnotesize 数}} &
\multicolumn{1}{c}{\footnotesize\log_{\epsilon}} \\
\PadTo[l]{1.1}{1}
    & 0.0000 &&      6 & 1.7918 \\
1.1 & 0.0953 &&      7 & 1.9459 \\
1.2 & 0.1823 &&      8 & 2.0794 \\
1.5 & 0.4055 &&      9 & 2.1972 \\
1.7 & 0.5306 &&     10 & 2.3026 \\
2.0 & 0.6931 &&     20 & 2.9957 \\
2.2 & 0.7885 &&     50 & 3.9120 \\
2.5 & 0.9163 &&    100 & 4.6052 \\
2.7 & 0.9933 &&    200 & 5.2983 \\
2.8 & 1.0296 &&    500 & 6.2146 \\
3.0 & 1.0986 &&  1,000 & 6.9078 \\
3.5 & 1.2528 &&  2,000 & 7.6009 \\
4.0 & 1.3863 &&  5,000 & 8.5172 \\
4.5 & 1.5041 && 10,000 & 9.2103 \\
5.0 & 1.6094 && 20,000 & 9.9035 \\
\end{array}
\]
\end{table}

次に、対数や指数を含むいくつかの式について微分してみましょう。

方程式を以下のように変形します：
\[
y = \log_\epsilon x.
\]
まず、これを
\[
\epsilon^y = x,
\]
と書き換えます。これにより、$y$ に関する $\epsilon^y$ の微分は元の関数が変わらない（参照：\Pageref{unchanged}）ので、
\[
\frac{dx}{dy} = \epsilon^y,
\]
となります。逆関数に戻ると、
\[
\frac{dy}{dx}
  = \frac{1}{\frac{dx}{dy}}
  = \frac{1}{\epsilon^y}
  = \frac{1}{x}.
\]

これは非常に興味深い結果です。これは次のように書けます：
\[
\frac{d(\log_\epsilon x)}{dx} = x^{-1}.
\]

$x^{-1}$ は、べき関数の微分ルールでは得られません。そのルール（参照：\Pageref[page]{multipow}）は、べき乗を掛けて指数を1減らすことです。したがって、$x^3$ の微分は $3x^2$ となり、$x^2$ の微分は $2x$ となります。しかし、$x^0$ の微分は $x^{-1}$ や $0 \times x^{-1}$ にはなりません。なぜなら $x^0 = 1$ であり、定数だからです。$\log_\epsilon x$ の微分が $\frac{1}{x}$ となる奇妙な事実については、積分の章で再び触れることになります。

次に、
\begin{align*}
y &= \log_\epsilon(x+a),\\
\epsilon^y &= x+a;
\end{align*}
を微分します。$\epsilon^y$ の微分は $\epsilon^y$ なので、
\[
\frac{d(x+a)}{dy} = \epsilon^y,
\]
となります。

\begin{align*}
\frac{dx}{dy} &= \epsilon^y = x+a; \\
\intertext{逆関数に戻ると（参照：\Pageref{section:3}）}
\frac{dy}{dx} &= \frac{1}{\frac{dx}{dy}} = \frac{1}{x+a}.
\end{align*}
}

次に、
\begin{align*}
y &= \log_{10} x.
\end{align*}

自然対数に変換するために $0.4343$ で掛けてみます。これにより、
\begin{align*}
y &= 0.4343 \log_\epsilon x; \\
\frac{dy}{dx} &= \frac{0.4343}{x}.
\end{align*}

次は少し複雑です。試してみましょう：
\[
y = a^x.
\]

両辺を対数取ると、
\begin{align*}
\log_\epsilon y &= x \log_\epsilon a, \\
x  &= \frac{\log_\epsilon y}{\log_\epsilon a} 
    = \frac{1}{\log_\epsilon a} \times  \log_\epsilon y.
\end{align*}

Since $\dfrac{1}{\log_\epsilon a}$ は定数であるから、次のようになります：
\[
\dfrac{dx}{dy}
  = \dfrac{1}{\log_\epsilon a} \times  \dfrac{1}{y}
  = \dfrac{1}{a^x \times  \log_\epsilon a};
\]
従って元の関数に戻ると、
\[
\dfrac{dy}{dx} = \dfrac{1}{\;\dfrac{dx}{dy}\;} = a^x \times  \log_\epsilon a.
\]

また、
\[
\dfrac{dx}{dy} \times  \dfrac{dy}{dx} = 1\quad\text{および}\quad
\dfrac{dx}{dy} = \dfrac{1}{y} \times  \dfrac{1}{\log_\epsilon a},\quad
\dfrac{1}{y} \times  \dfrac{dy}{dx} = \log_\epsilon a.
\]

したがって、$\log_\epsilon y$ が $ の関数である場合、常に $\dfrac{1}{y}\, \dfrac{dy}{dx} =$ ($ の関数の微分係数) であり、
従って $\log_\epsilon y = x \log_\epsilon a$ からすぐに
\[
\dfrac{1}{y}\, \dfrac{dy}{dx}
  = \log_\epsilon a\quad\text{および}\quad
\dfrac{dy}{dx} = a^x \log_\epsilon a.
\]

次に、さらに例を試みましょう。

例題．
(1) $y=\epsilon^{-ax}$．$-ax=z$ とすると $y=\epsilon^z$．
\[
\dfrac{dy}{dx} = \epsilon^z;\quad
\dfrac{dz}{dx} = -a;\quad\text{従って}\quad
\dfrac{dy}{dx} = -a\epsilon^{-ax}.
\]

または：
\[
\log_\epsilon y = -ax;\quad
\dfrac{1}{y}\, \dfrac{dy}{dx} = -a;\quad
\dfrac{dy}{dx} = -ay = -a\epsilon^{-ax}.
\]

(2) $y=\epsilon^{\frac{x^2}{3}}$．$\dfrac{x^2}{3}=z$ とすると $y=\epsilon^z$．
\[
\dfrac{dy}{dz} = \epsilon^z;\quad
\dfrac{dz}{dx} = \dfrac{2x}{3};\quad
\dfrac{dy}{dx} = \dfrac{2x}{3}\, \epsilon^{\frac{x^2}{3}}.
\]

または：
\[
\log_\epsilon y = \dfrac{x^2}{3};\quad
\dfrac{1}{y}\, \dfrac{dy}{dx} = \dfrac{2x}{3};\quad
\dfrac{dy}{dx} = \dfrac{2x}{3}\, \epsilon^{\frac{x^2}{3}}.
\]

(3) $y = \epsilon^{\frac{2x}{x+1}}$．
\begin{align*}
\log_\epsilon y &= \dfrac{2x}{x+1},\quad
\dfrac{1}{y}\, \dfrac{dy}{dx} = \dfrac{2(x+1)-2x}{(x+1)^2}; \\
従って
\dfrac{dy}{dx} &= \dfrac{2}{(x+1)^2} \epsilon^{\frac{2x}{x+1}}.
\end{align*}

確認のために $\dfrac{2x}{x+1}=z$ と置きます。

(4) \( y = \epsilon^{\sqrt{x^2+a}} \). \quad \(\log_\epsilon y = (x^2+a)^{\frac{1}{2}}\).
\[
\frac{1}{y} \, \frac{dy}{dx} = \frac{x}{(x^2+a)^{\frac{1}{2}}} \quad \text{および} \quad
\frac{dy}{dx} = \frac{x \times \epsilon^{\sqrt{x^2+a}}}{(x^2+a)^{\frac{1}{2}}}.
\]
（もし \( (x^2+a)^{\frac{1}{2}} = u \) および \( x^2+a = v \) のとき、\( u = v^{\frac{1}{2}} \),
\[
\frac{du}{dv} = \frac{1}{2v^{\frac{1}{2}}} ; \quad
\frac{dv}{dx} = 2x ; \quad
\frac{du}{dx} = \frac{x}{(x^2+a)^{\frac{1}{2}}}.
\])

\(\sqrt{x^2+a} = z\) を代入して確認せよ。

(5) \( y = \log(a+x^3) \). 令 \( a+x^3 = z \); すると \( y = \log_\epsilon z \).
\[
\frac{dy}{dz} = \frac{1}{z}; \quad
\frac{dz}{dx} = 3x^2; \quad \text{従って}\quad
\frac{dy}{dx} = \frac{3x^2}{a+x^3}.
\]

(6) \( y = \log_\epsilon \{3x^2 + \sqrt{a+x^2}\} \). 令 \( 3x^2 + \sqrt{a+x^2} = z \);
すると \( y = \log_\epsilon z \).
\begin{align*}
\frac{dy}{dz} &= \frac{1}{z}; \quad \frac{dz}{dx} = 6x + \frac{x}{\sqrt{x^2+a}}; \\
\frac{dy}{dx} &= \frac{6x + \dfrac{x}{\sqrt{x^2+a}}}{3x^2 + \sqrt{a+x^2}}
             = \frac{x(1 + 6\sqrt{x^2+a})}{(3x^2 + \sqrt{a+x^2}) \sqrt{x^2+a}}.
\end{align*}

(7) \( y = (x+3)^2 \sqrt{x-2} \).
\begin{align*}
\log_\epsilon y &= 2 \log_\epsilon(x+3) + \tfrac{1}{2} \log_\epsilon(x-2). \\
\frac{1}{y} \, \frac{dy}{dx} &= \frac{2}{x+3} + \frac{1}{2(x-2)}; \\
\frac{dy}{dx} &= (x+3)^2 \sqrt{x-2} \left\{\frac{2}{x+3} + \frac{1}{2(x-2)}\right\}.
\end{align*}

(8) \( y = (x^2+3)^3 (x^3-2)^{\frac{2}{3}} \).
\begin{align*}
\log_\epsilon y &= 3 \log_\epsilon(x^2+3) + \tfrac{2}{3} \log_\epsilon(x^3-2); \\
\frac{1}{y} \, \frac{dy}{dx} &= 3 \frac{2x}{x^2+3} + \frac{2}{3} \frac{3x^2}{x^3-2}
                            = \frac{6x}{x^2+3} + \frac{2x^2}{x^3-2}.
\end{align*}
（もし \( y = \log_\epsilon(x^2+3) \) のとき、\( x^2+3 = z \) および \( u = \log_\epsilon z \).
\[
\frac{du}{dz} = \frac{1}{z}; \quad
\frac{dz}{dx} = 2x; \quad
\frac{du}{dx} = \frac{2x}{x^2+3}.
\]
同様に、\( v = \log_\epsilon(x^3-2) \) のとき、\(\dfrac{dv}{dx} = \dfrac{3x^2}{x^3-2}\) であり
\[
\frac{dy}{dx}
= (x^2+3)^3 (x^3-2)^{\frac{2}{3}}
  \left\{ \frac{6x}{x^2+3} + \frac{2x^2}{x^3-2} \right\}.
\]

(9) \( y = \frac{\sqrt[2]{x^2 + a}}{\sqrt[3]{x^3 - a}} \).
\begin{align*}
\log_\epsilon y &= \frac{1}{2} \log_\epsilon (x^2 + a) - \frac{1}{3} \log_\epsilon (x^3 - a). \\
\frac{1}{y} \, \frac{dy}{dx} &= \frac{1}{2} \cdot \frac{2x}{x^2 + a} - \frac{1}{3} \cdot \frac{3x^2}{x^3 - a}
   = \frac{x}{x^2 + a} - \frac{x^2}{x^3 - a}. \\
\frac{dy}{dx} &= \frac{\sqrt[2]{x^2 + a}}{\sqrt[3]{x^3 - a}}
     \left\{ \frac{x}{x^2 + a} - \frac{x^2}{x^3 - a} \right\}.
\end{align*}
(10) \( y = \frac{1}{\log_\epsilon x} \)
\[
\frac{dy}{dx} = \frac{\log_\epsilon x \cdot 0 - 1 \cdot \frac{1}{x}} {\log_\epsilon^2 x}
   = -\frac{1}{x \log_\epsilon^2 x}.
\]

(11) \( y = \sqrt[3]{\log_\epsilon x} = (\log_\epsilon x)^{\frac{1}{3}} \). 令 \( z = \log_\epsilon x \); \( y = z^{\frac{1}{3}} \).
\[
\frac{dy}{dz} = \frac{1}{3} z^{-\frac{2}{3}}; \quad
\frac{dz}{dx} = \frac{1}{x}; \quad
\frac{dy}{dx} = \frac{1}{3x \sqrt[3]{\log_\epsilon^2 x}}.
\]

(12) \( y = \left(\dfrac{1}{a^x}\right)^{ax} \).
\begin{align*}
\log_\epsilon y &= ax (\log_\epsilon 1 - \log_\epsilon a^x) = -ax \log_\epsilon a^x. \\
\frac{1}{y} \, \frac{dy}{dx} &= -ax \cdot a^x \log_\epsilon a - a \log_\epsilon a^x.  \\
\frac{dy}{dx} &= -\left(\frac{1}{a^x}\right)^{ax}
      (x \cdot a^{x+1} \log_\epsilon a + a \log_\epsilon a^x).
\end{align*}

次に、以下の練習問題を試みよ。

\par\textbf{【演習 XII】}\par （解答は \Pageref[page]{AnsEx:XII} を参照のこと。）

\begin{enumerate}
\item[(1)] \( y = b (\epsilon^{ax} - \epsilon^{-ax}) \) の微分を求めよ。

\item[(2)] \( u = at^2 + 2 \log_\epsilon t \) に対する \( t \) に関する導関数を求めよ。

\item[(3)] \( y = n^t \) のとき、\( \dfrac{d(\log_\epsilon y)}{dt} \) を求めよ。

\item[(4)] \( y = \dfrac{1}{b} \cdot \dfrac{a^{bx}}{\log_\epsilon a} \) のとき、\( \dfrac{dy}{dx} = a^{bx} \) を示せ。

\item[(5)] \( w = pn^v \) のとき、\( \dfrac{dw}{dv} \) を求めよ。

次に微分せよ。

\item[(6)] \( y = \log_\epsilon x^n \).

\item[(7)] \( y = 3 \epsilon^{-\frac{x}{x-1}} \).

\item[(8)] \( y = (3x^2 + 1) \epsilon^{-5x} \).
\item[(9)] \( y = \log_\epsilon (x^a + a) \).

\item[(10)] \( y = (3x^2 - 1)(\sqrt{x} + 1) \).

\item[(11)] \( y = \dfrac{\log_\epsilon (x+3)}{x+3} \).
\item[(12)] \( y = a^x \times x^a \).

\item[(13)] レオノル・ケルビンによって示されたように、海底ケーブルを通じた信号伝送速度は、芯の外径と包まれた銅線の直径との比に依存します。この比を$y$と呼ぶとき、1分間に送れる信号数$s$は以下の公式で表されます：
\[
s=ay^2 \log_\epsilon \frac{1}{y};
\]
ここで$a$は長さや材料の質によって定数です。これらの値が与えられた場合、$s$は$y=1 ÷ \sqrt{\epsilon}$のときに最大となることを示してください。

\item[(14)] 次の式の極大または極小を求めよ：
\[
y=x^3-\log_\epsilon x.
\]

\item[(15)] $y=\log_\epsilon(ax\epsilon^x)$を微分せよ。

\item[(16)] $y=(\log_\epsilon ax)^3$を微分せよ。

\tb

\section*{対数曲線。}

次に、連続した各 ordinates が幾何級数で関連している曲線について考えましょう。例えば、方程式 $y=bp^x$ によって表されるような曲線です。

$x=0$ のとき、$b$ は $y$ の初期値となります。

次に
\[
x=1,\quad y=bp;\qquad
x=2,\quad y=bp^2;\qquad
x=3,\quad y=bp^3,\quad \text{etc.}
\]
とします。また、$p$ は任意のordinate の高さとその前のordinate の高さとの比の数値を表します。図40では $p=\frac{6}{5}$ としており、各ordinate は前ordinate の $\frac{6}{5}$ 倍の高さです。

\par\textbf{〔第167b図〕}\par{40}{41}

もし連続した2つのordinate が一定の比で関連している場合、その対数値は一定の差を持つことから、$\log_\epsilon y$ をordinate として新しい曲線を描くと、図41のように直線的に上に傾斜する曲線になります。実際には、方程式より
\begin{align*}
\log_\epsilon y &= \log_\epsilon b + x · \log_\epsilon p, \\
\lintertext{従って}
\log_\epsilon y &- \log_\epsilon b = x · \log_\epsilon p.
\end{align*}

ここで $\log_\epsilon p$ はただの数であり、$\log_\epsilon p=a$ と書き換えることができます。したがって
\[
\log_\epsilon \frac{y}{b}=ax,
\]
となり、方程式は新しい形で表されます：
\[
y = b\epsilon^{ax}.
\]

\section*{衰減曲線。}

もし $p$ を分数（1未満）として取ると、曲線は明らかに下に傾くことになります。図42では各ordinate が前ordinate の $\frac{3}{4}$ 倍の高さであることが示されています。

方程式はまだ
\[
y = bp^x;
\]
の形を保っています。
ただし、$p$ が 1 より小さいため $\log_\epsilon p$ は負の量となり、これを $-a$ と書き表すことができます。よって $p = \epsilon^{-a}$ であり、現在の曲線の方程式は
\[
y = b\epsilon^{-ax}
\]
となります。

この表現の重要性は、独立変数が時間である場合、方程式が多くの物理過程を表すことを示していることです。これらの過程では何かが徐々に衰えています。例えば、ニュートンの有名な「冷却法則」では熱体の冷却を
\[
\theta_t = \theta_0 \epsilon^{-at}
\]
で表現します。
ここで $\theta_0$ は熱体の初期温度差、$\theta_t$ は時間 $t$ 経過後の温度差、$a$ は定数であり、これは露出面積や導電率・放散率などの要素による減衰定数です。

同様の式
\[
Q_t = Q_0 \epsilon^{-at}
\]
は、初期に電荷 $Q_0$ を持つ体が一定の減衰定数 $a$ で漏れていったときの電荷を表します。この定数はこの場合、体の容量と漏れ道の抵抗によって決まります。

柔らかいスプリングに与えられた振動は時間とともに衰えていきます。そして、運動の振幅がどのように衰えるかは似た方法で表現できます。

実際、$\epsilon^{-at}$ は減少率が減少する量に対して比例することを示す「衰減因子」であり、または通常の記号では $\dfrac{dy}{dt}$ がその瞬間の $y$ の値に比例しているすべての現象に対応します。つまり、曲線 \Fig{42} を見れば、どの部分でも傾き $\dfrac{dy}{dx}$ は $y$ の高さに対して比例しており、$y$ が小さくなるにつれて曲線は平坦になります。記号で表すと
\begin{gather*}
y = b\epsilon^{-ax}\\
\log_\epsilon y = \log_\epsilon b - ax\\
\frac{1}{y} \frac{dy}{dx} = -a;\\
\therefore \frac{dy}{dx} = -ay;
\end{gather*}
すなわち、曲線の傾きは下向きで $y$ と定数 $a$ の比例関係にあります。

我々は、方程式を以下のように表す形で用いた場合にも同じ結果を得られたはずです：
\begin{align*}
y &= bp^x; \\
\text{すると} \quad \frac{dy}{dx} &= bp^x \times  \log_\epsilon p. \\
\text{しかし} \quad \log_\epsilon p &= -a; \\
\text{よって} \quad \frac{dy}{dx} &= y \times  (-a) = -ay,
\end{align*}
これは以前と同じです。

\par\medskip\noindent\textbf{時定数。}「死滅因子」\textasciitilde{}$\epsilon^{-at}$ の表記において、量\textasciitilde{}$a$ は「\emph{時定数}」と呼ばれる別の量の逆数であり、これを符号\textasciitilde{}$T$ で表すことができます。すると、「死滅因子」は\textasciitilde{}$\epsilon^{-\frac{t}{T}}$ として書き换えることができ、$t = T$ のとき、$T$（または $\dfrac{1}{a}$）の意味が元の量（前文では $\theta_0$ や $Q_0$ と呼ばれていたもの）がその原始値の $\dfrac{1}{\epsilon}$ 分、すなわち $0.3678$ 分だけ減少する時間であることがわかります。

物理学の様々な分野で $\epsilon^x$ および\textasciitilde{}$\epsilon^{-x}$ の値は頻繁に必要とされますが、数学的表にはそれらが非常に少ない数しか含まれていないため、便利のためにいくつかの値を以下に示します。

\begin{table}[p]
%
\[
\setlength{\arraycolsep}{1.5em}
\begin{array}{| .{2,2} | .{5,4} | .{1,6} | .{1,6} |}
\hline
\multicolumn{1}{|c|}{x} &
  \multicolumn{1}{c|}{ε^x} &
  \multicolumn{1}{c|}{ε^{-x}} &
  \multicolumn{1}{c|}{1-ε^{-x}} \\
\hline
0.00   &     1.0000 & 1.0000   & 0.0000   \\
0.10   &     1.1052 & 0.9048   & 0.0952   \\
0.20   &     1.2214 & 0.8187   & 0.1813   \\
0.50   &     1.6487 & 0.6065   & 0.3935   \\
0.75   &     2.1170 & 0.4724   & 0.5276   \\
0.90   &     2.4596 & 0.4066   & 0.5934   \\
1.00   &     2.7183 & 0.3679   & 0.6321   \\
1.10   &     3.0042 & 0.3329   & 0.6671   \\
1.20   &     3.3201 & 0.3012   & 0.6988   \\
1.25   &     3.4903 & 0.2865   & 0.7135   \\
1.50   &     4.4817 & 0.2231   & 0.7769   \\
1.75   &     5.755  & 0.1738   & 0.8262   \\
2.00   &     7.389  & 0.1353   & 0.8647   \\
2.50   &    12.182  & 0.0821   & 0.9179   \\
3.00   &    20.086  & 0.0498   & 0.9502   \\
3.50   &    33.115  & 0.0302   & 0.9698   \\
4.00   &    54.598  & 0.0183   & 0.9817   \\
4.50   &    90.017  & 0.0111   & 0.9889   \\
5.00   &   148.41   & 0.0067   & 0.9933   \\
5.50   &   244.69   & 0.0041   & 0.9959   \\
6.00   &   403.43   & 0.00248  & 0.99752  \\
7.50   &  1808.04   & 0.00055  & 0.99947  \\
10.00  & 22026.5    & 0.000045 & 0.999955 \\
\hline
\end{array}
\]
\end{table}

この表の使用例として、ある熱い物体が冷却しているとします。実験開始時（すなわち \( t = 0 \)）に、その物体は周囲の物体よりも72$^\circ$F hotterであり、その冷却の時間定数は20分間です（つまり、その温度余剰が72$^\circ$Fの\(\frac{1}{ε}\)となるまで20分かかる）。このとき、任意の時間 \( t \) におけるその温度余剰を計算することができます。例えば、\( t = 60 \) 分の場合を考えます。このとき \(\frac{t}{T} = 60 ÷ 20 = 3\) であり、\(\epsilon^{-3}\) の値を求めなければなりません。その後、元の72$^\circ$Fにこの値を乗じます。表から \(\epsilon^{-3}\) の値は0.0498 であることがわかります。したがって、60分後の温度余剰は \( 72^\circ \times  0.0498 = 3.586^\circ \) となります。

(1) 電流の強度は時間 \( t \) 秒後に電動力によって生成される導体において、以下のように与えられます：
\[
 C = \frac{E}{R} \left\{ 1 - \epsilon^{-\frac{Rt}{L}} \right\}. 
\]

時定数は \(\frac{L}{R}\) です。

もし \( E = 10 \), \( R = 1 \), \( L = 0.01 \) のとき、\( t \) が非常に大きいと項 \(\epsilon^{-\frac{Rt}{L}}\) は 1 に近づき、したがって \( C = \frac{E}{R} = 10 \) となります。また
\[
 \frac{L}{R} = T = 0.01. 
\]

任意の時間における値は以下のように表すことができます：
\[
 C = 10 - 10\epsilon^{-\frac{t}{0.01}}, 
\]
時定数が \( 0.01 \) であることを意味します。これは、変動項が元の値 \( 10\epsilon^{-\frac{0}{0.01}} = 10 \) のうち \(\frac{1}{\epsilon} = 0.3678\) 分だけ減少するのに約 \( 0.01 \) 秒かかるという意味です。

\( t = 0.001 \) 秒のときの電流を求めるために、\(\frac{t}{T} = 0.1\) であり、\(\epsilon^{-0.1} = 0.9048\)（表から）となります。したがって \( 0.001 \) 秒後には変動項は \( 0.9048 \times  10 = 9.048 \) となり、実際の電流は \( 10 - 9.048 = 0.952 \) となります。

同様に、\( 0.1 \) 秒後の値は
\[
 \frac{t}{T} = 10; \quad \epsilon^{-10} = 0.000045, 
\]
変動項は \( 10 \times  0.000045 = 0.00045 \) となり、電流は \( 9.9995 \) となります。

(2) 透明な媒質を通過した光の強度 \( I \) は以下のように与えられます：
\[
 I = I_0\epsilon^{-Kl}, 
\]
ここで \( I_0 \) は光の初期強度、\( K \) は「吸収定数」です。

この定数は実験によって求められます。例えば、ある透明な媒質を \( 10 \) cm 通過したときに光の強度が 18\% 减弱すると仮定すると、これは \( 82 = 100 \times  \epsilon^{-K\times 10} \) や \(\epsilon^{-10K} = 0.82\) を意味し、表から \( 10K = 0.20 \) と近似できます。したがって \( K = 0.02 \) となります。

強度を半分に減らすために必要な厚さ \( l \) を求めるには、等式 \( 50 = 100 \times  \epsilon^{-0.02l} \) や \( 0.5 = \epsilon^{-0.02l} \) を解く必要があります。これは対数形式に書き換えることで
\[
 \log 0.5 = -0.02 \times  l \times  \log \epsilon, 
\]
となり、これから
\[
 l = \frac{-0.3010}{-0.02 \times  0.4343} = 34.7 \text{ cm} \approx 34.5 \text{ cm}. 
\]

(3) 放射性物質の未変化量 \( Q \) は、その初期量 \( Q_0 \) と関係し、その関係は \( Q = Q_0 \epsilon^{-\lambda t} \) で与えられることが知られています。ここで、\(\lambda\) は定数であり、\( t \) は変化が始まってから経過した秒数です。

「ラジウム A」の場合、時間が秒単位で表されるとき、実験により \(\lambda = 3.85 \times  10^{-3}\) と示されています。半分の物質が変化するのに必要な時間を求めましょう。（この時間は「平均寿命」と呼ばれます）

\[
0.5 = \epsilon^{-0.00385t}
\]

\begin{align*}
\log 0.5 &= -0.00385t \times \log \epsilon; \\
\therefore t &= 3\text{ 分钟 非常に近似}。
\end{align*}

\par\textbf{【演習 XIII】}\par (答えは \Pageref[page]{AnsEx:XIII} を参照してください。)

\item[(1)] \( y = b \epsilon^{-\frac{t}{T}} \) の曲線を描き、\( b = 12 \), \( T = 8 \)，そして \( t \) が \( 0 \) から \( 20 \) までの値で与えられる場合。

\item[(2)] 熱い物体が冷却し、その温度余分が初期量の半分になるまでに \( 24 \) 分かかるとき、時間定数を求め、元の温度余分の \( 1 \) ％になるまでにどれくらい時間がかかるか求めよ。

\item[(3)] 曲線 \( y = 100(1-\epsilon^{-2t}) \) をプロットせよ。

\item[XIII:4]{(4)} 次の式は非常に類似した曲線を描きます：
\begin{align*}
\text{(i)}\   y &= \frac{ax}{x + b}; \\
\text{(ii)}\  y &= a(1 - \epsilon^{-\frac{x}{b}}); \\
\text{(iii)}\ y &= \frac{a}{90^\circ} \arctan \left(\frac{x}{b}\right).
\end{align*}

\( a = 100 \) ミリメートル、\( b = 30 \) ミリメートルとしてすべての曲線を描きなさい。

\item[(5)] \( y \) に対する \( x \) の微分係数を求めよ：
\[
(\textit{a})~y = x^x;\quad
(\textit{b})~y = (\epsilon^x)^x;\quad
(\textit{c})~y = \epsilon^{x^x}.
\]

\item[(6)] 「トーヒウム A」の場合、\(\lambda\) の値は \( 5 \) であり、半分の初期量 \( Q_0 \) となるまでに変化する時間（平均寿命）を求めよ。式は
\[
Q = Q_0 \epsilon^{-\lambda t};
\]
ここで \( t \) は秒単位です。

\item[(7)] 容量 \( K = 4 \times  10^{-6} \)，電圧 \( V_0 = 20 \) で充電されたコンデンサが \( 10,000 \) オームの抵抗を通じて放電しているとき、(\textit{a}) \( 0.1 \) 秒後；(\textit{b}) \( 0.01 \) 秒後の電圧 \( V \) を求めよ。電圧の減少は以下の規則に従うと仮定：\( V = V_0 \epsilon^{-\frac{t}{KR}} \)。

\item[(8)] 電離された孤立した金属球の電荷 \( Q \) は、10分間に20から16単位まで減少します。電荷 \( Q = Q_0 \times \epsilon^{-\mu t} \) の中で、\( Q_0 \) は初期電荷であり、\( t \) は秒数です。漏れによる係数 \( \mu \) を求めよ。その後、半分の電荷が漏れるのにかかる時間を求める。

\item[(9)] 電話線における減衰は、関係式 \( i = i_0 \epsilon^{-\beta l} \) から確認できる。ここで、\( t \) 秒後に初期強度 \( i_0 \) の電話信号の強度 \( i \) が与えられ、\( l \) は線路の長さ（km）であり、定数 \(\beta\) である。1910年に敷設されたフランスとイギリス間の海底ケーブルの場合、\(\beta = 0.0114\) である。ケーブルの終端での減衰を求める。また、非常に良い聴覚における元の電流の8%に達するまでの線路の長さを求めよ。

\item[(10)] 海抜 \( h \) km における大気圧 \( p \) は \( p = p_0 \epsilon^{-kh} \) で与えられる。ここで、\( p_0 \) は海面での圧力（760mm）である。10km、20km、50kmにおける大気圧がそれぞれ199.2、42.2、0.32である場合、各ケースで \( k \) を求めよ。\( k \) の平均値を使用して、各ケースでの誤差を求める。

\item[(11)] \( y = x^x \) の極値を求めよ。

\item[(12)] \( y = x^{\frac{1}{x}} \) の極値を求めよ。

\item[(13)] \( y = xa^{\frac{1}{x}} \) の極値を求めよ。

\end{enumerate}



\chapter{正弦（サイン）と余弦（コサイン）の微分法}
\noindent\textit{\large How to deal with Sines and Cosines}\\[
2mm]
\noindent\textbf{\small 【章の論理とテーマ】三角関数の幾何学的増分・sin x と cos x の導関数と周期的現象}\\[4mm]
章[SINES AND COSINES]{XV}{正弦と余弦の扱い}
ギリシャ文字が角度を表すために一般的に使用されるので、任意の変動する角度を表す通常の文字として$\theta$（「theta」）を使用します。
関数
\[
y = \sin \theta.
\]

図177a（43）

調査すべきは$\dfrac{d(\sin \theta)}{d \theta}$の値です；つまり、角度$\theta$が変化する場合、正弦の増分と角度の増分との関係を求める必要があります。両方の増分は自身で無限に小さいものとして考えます。図43（178.png）を見てください。円の半径が1であるとき、$y$の高さは正弦であり、$\theta$は角度です。現在、$\theta$に小さな角度$d \theta$を加えて増加させると、$y$の高さ、すなわち正弦は小さな要素$dy$で増加します。新しい高さ$y + dy$は新しい角度$\theta + d \theta$の正弦であり、方程式として表すと
\[
y+dy = \sin(\theta + d \theta);
\]
これを最初の方程式から引くと
\[
dy = \sin(\theta + d \theta)- \sin \theta.
\]

右辺の量は2つの正弦の差であり、三角法の本ではこの計算方法が説明されています。それは、$M$と$N$が異なる角度である場合、
\[
\sin M - \sin N = 2 \cos\frac{M+N}{2}·\sin\frac{M-N}{2}.
\]

したがって、$M= \theta + d \theta$を一つの角度として、$N= \theta$をもう一つの角度として置くと
\begin{align*}
dy &= 2 \cos\frac{\theta + d\theta + \theta}{2}
      · \sin\frac{\theta + d\theta - \theta}{2},\\
\lintertext{または}
dy &= 2\cos(\theta + \tfrac{1}{2}d\theta)
     · \sin\tfrac{1}{2} d\theta.
\end{align*}

しかし、$d \theta$が無限に小さいと見なされる場合、極限において$\frac{1}{2} d \theta$を$\theta$と比較して無視できます。また、$\sin\tfrac{1}{2} d \theta$は$\tfrac{1}{2} d \theta$とみなすことができます。方程式は次のようになります：
\begin{align*}
dy &= 2 \cos \theta \times  \tfrac{1}{2} d \theta; \\
dy &= \cos \theta · d \theta, \\
\lintertext{そして、最後に}
\dfrac{dy}{d \theta} &= \cos \theta.
\end{align*}

添付の曲線、図44および図45は、スケールに従って、$y=\sin \theta$と$\dfrac{dy}{d\theta}=\cos\theta$の値を、対応する$\theta$の値に対してプロットしています。
%[** TN: 原文の図は逆で、コサインは図44にラベル付けされていますなど。]
図44 179a ページラベル：erratum2
図45 179b
---
次にコサインを考慮します。ページラベル：differcos

$y=\cos \theta$とします。

現在、$\cos \theta=\sin\left(\frac{\pi}{2}-\theta\right)$です。

したがって
\begin{align*}
&\begin{aligned}
dy = d\left(\sin\left(\frac{\pi}{2} - \theta\right)\right)
  &= \cos\left(\frac{\pi}{2} - \theta\right) \times  d(-\theta), \\
  &= \cos\left(\frac{\pi}{2} - \theta\right) \times  (-d\theta),
\end{aligned} \\
\&\frac{dy}{d\theta} = -\cos\left(\frac{\pi}{2} - \theta\right).
\intertext{そして}
\&\frac{dy}{d\theta} = -\sin \theta.
\end{align*}

---
最後に、タンジェントを考慮します。
%
\begin{align*}
\lintertext{まず}
y  &= \tan \theta, \\
dy &= \tan(\theta + d\theta) - \tan\theta. \\
\intertext{三角法の教科書で示されるように展開すると}
\tan(\theta + d\theta)
   &= \frac{\tan\theta + \tan d\theta}
           {1 - \tan\theta·\tan d\theta}; \\
\lintertext{従って}
dy &= \frac{\tan\theta + \tan d\theta}
           {1-\tan\theta·\tan d\theta} - \tan\theta \\
   &= \frac{(1 + \tan^2\theta)\tan d\theta}
           {1-\tan\theta·\tan d\theta}.
\end{align*}
---
今、$d\theta$が無限に小さくなると、$\tan d\theta$の値は$d\theta$と同一視され、$\tan\theta · d\theta$は1に対して微小であるため、式は
\begin{align*}
dy &= \frac{(1+\tan^2 \theta)\, d\theta}{1}, \\
\lintertext{よって}
\frac{dy}{d\theta} &= 1 + \tan^2\theta, \\
\lintertext{または}
\frac{dy}{d\theta} &= \sec^2 \theta.
\end{align*}

これらの結果を集めて、以下の表を得ます：
\[
\begin{array}{|*{2}{>{\quad}c<{\quad}|}}
\hline
\DStrut y   & \dfrac{dy}{d\theta} \\
\hline
\sin\theta & \cos\theta \\
\cos\theta & -\sin\theta \\
\tan\theta & \sec^2 \theta\\
\hline
\end{array}
\]

時々、機械や物理学の問題において、例えば単純な調和運動や波動運動のような問題では、角度が時間に比例して増加することがあります。したがって、$T$ が一回の完全な周期（円周を一周する動き）の時間を表すとき、円周全体の角度は $2\pi$ ラジアンまたは $360^\circ$ であるため、時間 $t$ における角度の移動量は以下のようになります：
\[
\theta = 2\pi \frac{t}{T}, \quad \text{(ラジアン)}
\]
または
\[
\theta = 360 \frac{t}{T}, \quad \text{(度)}
\]

周期数（一秒あたりの周期数）を $n$ で表すと、$n = \dfrac{1}{T}$ であり、次のように書くことができます：
\[
\theta = 2\pi nt
\]
すると
\[
y = \sin 2\pi nt
\]

今、$\sin$ の時間に対する変化を知りたい場合、$\theta$ ではなく $t$ に対して微分する必要があります。そのためには第 IX 章で説明した方法を使用し、以下の式を使うことができます：
\[
\frac{dy}{dt} = \frac{dy}{d\theta} · \frac{d\theta}{dt}
\]

ここで $\dfrac{d\theta}{dt}$ は明らかに $2\pi n$ であるため、
\begin{align*}
\frac{dy}{dt} &= \cos \theta \times  2\pi n \\
              &= 2\pi n · \cos 2\pi nt.
\end{align*}

同様に、以下のことが導けます：
\[
\frac{d(\cos 2\pi nt)}{dt} = -2\pi n · \sin 2\pi nt.
\]

二重微分係数についての節

私たちは $\sin \theta$ を $\theta$ に対して微分すると $\cos \theta$ になることを見た。また、$\cos \theta$ を $\theta$ に対して微分すると $-\sin \theta$ になることも確認した。つまり、シンボルで表すと以下のようになります：
\[
\frac{d^2(\sin \theta)}{d\theta^2} = -\sin \theta.
\]

この結果は興味深いものであり、二度微分すると元の関数に戻るが符号が反転することを示しています。

これはコサインについても同じことが言えます。$\cos \theta$ を $\theta$ に対して微分すると $-\sin \theta$ になり、さらに $-\sin \theta$ を $\theta$ に対して微分すると $-\cos \theta$ になります。つまり：
\[
\frac{d^2(\cos \theta)}{d\theta^2} = -\cos \theta.
\]

$\sin$ と $\cos$ は唯一の二重微分係数が元の関数と符号を反転させて等しい関数です。

(1) $y=\arcsin x$。

もし$y$が$\sin y=x$の弧であるならば、$x=\sin y$。
\[
\frac{dx}{dy}=\cos y.
\]

逆関数から元の関数に移り換えると、
\begin{align*}
\frac{dy}{dx}
  &= \frac{1}{\;\dfrac{dx}{dy}\;} = \frac{1}{\cos y}. \\
\text{ここで}
\cos y
  &= \sqrt{1-\sin^2 y}=\sqrt{1-x^2}; \\
\text{従って}
\frac{dy}{dx}
  &= \frac{1}{\sqrt{1-x^2}},
\end{align*}
これは予想外の結果である。
(2) $y=\cos^3 \theta$。

これは$y=(\cos \theta)^3$と同じである。

$\cos\theta=v$とすると、$y=v^3$であり、$\frac{dy}{dv}=3v^2$。
\begin{align*}
\frac{dv}{d\theta} &= -\sin\theta.\\
\frac{dy}{d\theta} &=  \frac{dy}{dv} \times  \frac{dv}{d\theta}
                    = -3 \cos^2 \theta \sin\theta.
\end{align*}

(3) $y=\sin(x+a)$。

$x+a=v$とすると、$y=\sin v$。
\[
\frac{dy}{dv}=\cos v;\qquad
\frac{dv}{dx}=1 \quad\text{および}\quad
\frac{dy}{dx}=\cos(x+a).
\]

(4) $y=\log_\epsilon \sin \theta$。

$\sin\theta=v$とすると、$y=\log_\epsilon v$。
\begin{align*}
\frac{dy}{dv}      &= \frac{1}{v};\quad \frac{dv}{d\theta}=\cos\theta;\\
\frac{dy}{d\theta} &= \frac{1}{\sin\theta} \times  \cos\theta = \cot\theta.
\end{align*}

(5) $y=\cot\theta=\dfrac{\cos\theta}{\sin\theta}$。
\begin{align*}
\frac{dy}{d\theta}
  &= \frac{-\sin^2\theta - \cos^2 \theta}{\sin^2 \theta}\\
  &= -(1+\cot^2 \theta) = -\cosec^2 \theta.
\end{align*}

(6) $y=\tan 3\theta$。

$3\theta=v$とすると、$y=\tan v$であり、$\frac{dy}{dv}=\sec^2 v$。
\[
\frac{dv}{d\theta}=3;\quad
\frac{dy}{d\theta}=3 \sec^2 3\theta.
\]

(7) $y = \sqrt{1+3\tan^2\theta}$；$y=(1+3 \tan^2 \theta)^{\frac{1}{2}}$。

$3\tan^2\theta=v$とすると。
\begin{align*}
y &= (1+v)^{\frac{1}{2}};\quad
\frac{dy}{dv} = \frac{1}{2\sqrt{1+v}} \text{（参照：\Pageref{ExNo1}）;} \\
%
\frac{dv}{d\theta}
  &= 6\tan\theta \sec^2 \theta \\
%
\text{ここで}
v &= 3u^2;\quad \frac{dv}{du} = 6u;\quad \frac{du}{d\theta} = \sec^2 \theta; \\
%
\text{従って}
\frac{dv}{d\theta}
  &= 6 (\tan \theta \sec^2 \theta) \\
\text{従って}
%
\frac{dy}{d\theta}
  &= \frac{6\tan\theta \sec^2\theta}{2\sqrt{1 + 3\tan^2\theta}}.
\end{align*}

(8) \( y = \sin x \cos x \). 
\begin{align*}
\frac{dy}{dx} 
  &= \sin x(-\sin x) + \cos x \times \cos x \\
  &= \cos^2 x - \sin^2 x.
\end{align*}

\par\textbf{【演習 XIV】}\par (Answers are found on \Pageref[page]{AnsEx:XIV}.)

\begin{enumerate}
\item[(1)] 次の関数を微分せよ：
\begin{align*}
\text{(i)}\quad   y &= A \sin\left(\theta - \frac{\pi}{2}\right).\\
\text{(ii)}\quad  y &= \sin^2 \theta;\quad \text{and } y = \sin 2\theta.\\
\text{(iii)}\quad y &= \sin^3 \theta;\quad \text{and } y = \sin 3\theta.
\end{align*}

\item[(2)] \(\sin\theta \times \cos\theta\) の最大値をもつ \(\theta\) の値を求めよ。

\item[(3)] \(y=\dfrac{1}{2\pi} \cos 2\pi nt\) を微分せよ。

\item[(4)] \(y = \sin a^x\) のとき, \(\frac{dy}{dx}\) を求めよ。

\item[(5)] \(y=\log_\epsilon \cos x\) を微分せよ。

\item[(6)] \(y=18.2 \sin(x+26^\circ)\) を微分せよ。

\item[(7)] 曲線 \(y=100 \sin(\theta-15^\circ)\) を描き, \(\theta = 75^\circ\) における曲線の傾きが最大傾きの半分であることを示せ。

\item[(8)] \(y=\sin \theta·\sin 2\theta\) のとき, \(\frac{dy}{d\theta}\) を求めよ。

\item[(9)] \(y=a·\tan^m(\theta^n)\) のとき, \(\theta\) に関する微分係数を求めよ。

\item[(10)] \(y=\epsilon^x \sin^2 x\) を微分せよ。

\item[(11)] Exercises XIII (No. 4) (\Pageref{XIII:4}) の3つの方程式を微分し, \(x\) の非常に小さい値や非常に大きい値, または \(x=30\) の近傍での微分係数が等しいか, 近似して等しいかどうかを比較せよ。

\item[(12)] 次の関数を微分せよ：
\begin{align*}
\text{(i)}\quad   y &= \sec x.    &
\text{(ii)}\quad  y &= \arccos x. \\
\text{(iii)}\quad y &= \arctan x. &
\text{(iv)}\quad  y &= \arcsec x. \\
\text{(v)}\quad   y &= \tan x \times  \sqrt{3 \sec x}. &&
\end{align*}

\item[(13)] \(y=\sin(2\theta +3)^{2.3}\) を微分せよ。

\item[(14)] \(y=\theta^3+3 \sin(\theta+3)-3^{\sin \theta} - 3^\theta\) を微分せよ。

\item[(15)] \(y=\theta \cos \theta\) の最大値または最小値を求めよ。

\end{enumerate}



\chapter{偏微分法}
\noindent\textit{\large Partial Differentiation}\\[
2mm]
\noindent\textbf{\small 【章の論理とテーマ】多変数関数・一方の変数を固定して微分する偏導関数の直観的理解}\\[4mm]
第XVI章　部分微分
しばしば、一つの独立変数以外にも依存変数がある量に遭遇することがあります。例えば、$y$ が二つの他の変動する量 $u$ と $v$ の関数である場合があります。記号で表すと、
\begin{align*}
y &= f(u, v).
\end{align*}
最も単純な具体的な例を考えてみましょう。
\begin{align*}
y &= u \times v.
\end{align*}
では、どうすればよいでしょうか？$v$ を定数として $u$ について微分すると、
\begin{align*}
dy_v &= v \, du;
\end{align*}
または、$u$ を定数として $v$ について微分すると、
\begin{align*}
dy_u &= u \, dv.
\end{align*}
ここで、下付きの小文字は、どの量が操作中に定数とされたかを示しています。
また、微分が一部のみ行われたことを表すために、通常の $d$ ではなくギリシャ語の $\partial$ を使用します。この方法では、
\begin{align*}
\frac{\partial y}{\partial u} &= v, \\
\frac{\partial y}{\partial v} &= u.
\end{align*}
これらの値をそれぞれ $v$ と $u$ に代入すると、
\[
dy_v = \frac{\partial y}{\partial u} \, du, \quad dy_u = \frac{\partial y}{\partial v} \, dv
\]

となり、これらは「部分微分」です。

しかし、考えてみると、$y$ の総変動はこれらの両方の影響を受けることがわかります。つまり、どちらも変化している場合、実際の $dy$ は、

\[
dy = \frac{\partial y}{\partial u} \, du + \frac{\partial y}{\partial v} \, dv
\]

と書かれるべきであり、これが「総微分」です。一部の本では、

\[
dy = \left( \frac{dy}{dx} \right) \, dx + \left( \frac{dy}{dv} \right) \, dv
\]

と表されます。

例（1）．式 $w = 2ax^2 + 3bxy + 4cy^3$ の部分微分係数を求めます。答えは、

\[
\frac{\partial w}{\partial x} = 4ax + 3by, \quad \frac{\partial w}{\partial y} = 3bx + 12cy^2
\]

となります。

まず $y$ を定数として微分したものが最初の式、次に $x$ を定数として微分したものが第二の式です。すると、

\[
dw = (4ax + 3by) \, dx + (3bx + 12cy^2) \, dy
\]

\par\textbf{【例題 (2).】}\par \( z = x^y \)とする。最初に \( y \) を定数として、次に \( x \) を定数として通常の方法で微分すると、

\[
\left.
\begin{aligned}
\dfrac{\partial z}{\partial x} &= yx^{y-1}, \\
\dfrac{\partial z}{\partial y} &= x^y \log_\epsilon x,
\end{aligned}\right\}
\]

よって \( dz = yx^{y-1} \, dx + x^y \log_\epsilon x \, dy \)。

\par\textbf{【例題 (3).】}\par 高さが \( h \) で、底面の半径が \( r \) の円錐について、体積は \( V = \frac{1}{3} \pi r^2 h \) である。高さが固定され、半径 \( r \) だけが変化するときと、高さが変化し、半径 \( r \) が固定されるときの体積の変化率は異なる。これは、

\[
\left.
\begin{aligned}
\dfrac{\partial V}{\partial r} &= \dfrac{2\pi}{3} rh, \\
\dfrac{\partial V}{\partial h} &= \dfrac{\pi}{3} r^2,
\end{aligned}\right\}
\]

によって示される。両方の変数が変化するときの変動は \( dV = \dfrac{2\pi}{3} rh \, dr + \dfrac{\pi}{3} r^2 \, dh \) である。

例（4）。例 \pageref{Example4}において、$F$ と $f$ はどのような形の任意の関数を表します。例えば、これらは正弦関数や指数関数、あるいは独立変数 $t$ および $x$ の単なる代数関数であるかもしれません。この点が理解された上で、次のような式を考えます：
\begin{align*}
y &= F(x+at) + f(x-at), \\
\text{または} \quad y &= F(w) + f(v); \\
\text{ここで} \quad w &= x+at, \quad v = x-at. \\
\text{すると} \quad \frac{\partial y}{\partial x}
&= \frac{\partial F(w)}{\partial w} \cdot \frac{\partial w}{\partial x}
+ \frac{\partial f(v)}{\partial v} \cdot \frac{\partial v}{\partial x} \\
&= F'(w) \cdot 1 + f'(v) \cdot 1
\quad (ここで、数字 $1$ は $w$ および $v$ の $x$ の係数です); \\
\text{そして} \quad \frac{\partial^2 y}{\partial x^2}
&= F''(w) + f''(v). && \\
\text{また} \quad \frac{\partial y}{\partial t}
&= \frac{\partial F(w)}{\partial w} \cdot \frac{\partial w}{\partial t}
+ \frac{\partial f(v)}{\partial v} \cdot \frac{\partial v}{\partial t} \\
&= F'(w) \cdot a - f'(v) a; \\
\text{そして} \quad \frac{\partial^2 y}{\partial t^2}
&= F''(w)a^2 + f''(v)a^2; \\
\text{従って} \quad \frac{\partial^2 y}{\partial t^2}
&= a^2 \, \frac{\partial^2 y}{\partial x^2}.
\end{align*}

この微分方程式は数学物理学において極めて重要です。

二つの独立変数の関数の最大値と最小値。

例（5）。再度、問題 IX. の No. 4 \pageref{Ex9No4} を取り上げます。

$x$ および $y$ は弦の二つの部分の長さを表します。残りの部分は $30-(x+y)$ であり、三角形の面積は $A = \sqrt{s(s-x)(s-y)(s-30+x+y)}$ で与えられ、$s$ は半周長 $15$ であるので、$A = \sqrt{15P}$ と表されます。ここで
\begin{align*}
P &= (15-x)(15-y)(x+y-15) \\
  &= xy^2 + x^2y - 15x^2 - 15y^2 - 45xy + 450x + 450y - 3375.
\end{align*}

Clearly \( A \)は最大となるとき、\( P \)も最大となります。
\[
dP = \frac{\partial P}{\partial x}\, dx + \frac{\partial P}{\partial y}\, dy.
\]
最大値の場合（最小値にはならない）、同時に
%
\[
\frac{\partial P}{\partial x} = 0 \quad\text{and}\quad
\frac{\partial P}{\partial y} = 0;
\]
\begin{align*}
%[** Can't top-align the equations without enlarging the right brace]
\lintertext{すなわち}
\left.
\begin{aligned}
2xy - 30x + y^2 - 45y + 450 &= 0, \\
2xy - 30y + x^2 - 45x + 450 &= 0.
\end{aligned}
\right\}
\end{align*}}%

直感的な解は \( x = y \) です。

今、この条件を \( P \) の値に導入すると、
\[
P = (15-x)^2 (2x-15) = 2x^3 - 75x^2 + 900x - 3375.
\]
最大または最小の場合、\(\frac{dP}{dx} = 6x^2 - 150x + 900 = 0\) となり、これより \( x = 15 \) や \( x = 10 \)。

明らかに \( x = 15 \) は最小面積をもたらし、\( x = 10 \) は最大面積をもたらす。なぜなら、\(\frac{d^2 P}{dx^2} = 12x - 150\) であり、\( x = 15 \) のとき +30、\( x = 10 \) のとき -30 である。

例題 (6)。与えられた体積 \( V \) に対して、側面と底板の合計面積が最小となる普通の鉄道用煤車の寸法を求めよ。

この車は上部を開けた長方体である。\( x \) を長さ、\( y \) を幅とすると、深さは \( \frac{V}{xy} \) となる。表面積は \( S = xy + \frac{2V}{x} + \frac{2V}{y} \) である。
\[
dS = \frac{\partial S}{\partial x}\, dx
   + \frac{\partial S}{\partial y}\, dy
   = \left(y - \frac{2V}{x^2}\right) dx
   + \left(x - \frac{2V}{y^2}\right) dy.
\]
最小値の場合（最大値にはならない）、
\[
y - \frac{2V}{x^2} = 0,\quad
x - \frac{2V}{y^2} = 0.
\]

また、直感的な解は \( x = y \) であり、このとき
\( S = x^2 + \frac{4V}{x} \), \(\frac{dS}{dx}= 2x - \frac{4V}{x^2} = 0\) となり最小値をもたらし、
\[
x = \sqrt[3]{2V}.
\]

演習問題 XV（解答は Pageref[page]{AnsEx:XV} を参照のこと）

\begin{enumerate}
\item[(1)] \( x^3 - 2x^3y - 2y^2x + \frac{y}{3} \) を \( x \) だけ、また \( y \) だけ微分せよ。

\item[(2)] 表現式
\[
x^2yz + xy^2z + xyz^2 + x^2y^2z^2
\]
を \( x \), \( y \), \( z \) について偏微分せよ。

\item[(3)] \( r^2 = (x-a)^2 + (y-b)^2 + (z-c)^2 \).

$r$に関する$x$、$y$、$z$に対する偏微分の和を求めよ。また、二重偏微分の和も求めよ。
$$\frac{\partial r}{\partial x} + \frac{\partial r}{\partial y} + \frac{\partial r}{\partial z}$$
および
$$\frac{\partial^2r}{\partial x^2} + \frac{\partial^2r}{\partial y^2} + \frac{\partial^2r}{\partial z^2}$$

(4) $y = u^v$の全微分を求めよ。
(DPPageSep{193.png}{181}%)

(5) $y = u^3 \sin v$、$y = (\sin x)^u$、および $y = \frac{\log_\epsilon u}{v}$ の全微分を求めよ。

(6) 乘积为常数$k$的三个量$x$、$y$、$z$之和在这些三个量相等时达到最大值。

(7) 函数
$$u = x + 2xy + y$$
的最大或最小值求出。

(8) 邮局規定，包裹的长度加上周长不得超过6英尺。分别在（a）矩形截面；（b）圆形截面的情况下，可以寄送的最大体积是多少？

(9) 将$\pi$分为三部分，使得它们正弦值的连乘积达到最大或最小。

(10) 函数
$$u = \frac{\epsilon^{x+y}}{xy}$$
的最大或最小值求出。

(11) 函数
$$u = y + 2x - 2 \log_\epsilon y - \log_\epsilon x$$
的最大和最小值求出。

(12) 给定容量的吊桶具有水平等腰三角形棱柱形状，且顶点在下方，対面面开放。为了使建造所需铁皮最少，其尺寸应如何确定？
\end{enumerate}



\chapter{積分法}
\noindent\textit{\large Integration}\\[
2mm]
\noindent\textbf{\small 【章の論理とテーマ】微分の逆演算としての積分・不定積分と積分定数 C の正体}\\[4mm]
第XVII章 インテグラル
大いなる秘密は既に明らかになっています。この謎めいた記号$\displaystyle \int$、結局は長いSの形ですが、「合計」または「そのような量すべての合計」という意味です。従って、他の記号である$\sum$（ギリシャ語のシグマ）も同様に合計を表す記号です。しかし、数学的な男たちがこれらの記号を使用する際の慣習にはこの違いがあります：$\sum$は一般的に有限な量の合計を示すために使用されますが、インテグラル記号$\displaystyle \int$は大量の微小な量（無限に小さな要素）の合計を表すために一般的に使用されます。したがって、$\displaystyle \int dy = y$ および $\displaystyle \int dx = x$。
誰でも何か全体が多くの小さな部分から成り立つという考え方は理解できるでしょう；そして、部分が小さくなるほどその数は増える。例えば、1インチの線は10個の各$\frac{1}{10}$インチの部分で構成されるか、または100個の各部分が$\frac{1}{100}$インチの長さを持つ部分で構成されるか、あるいは1,000,000個の各部分が$\frac{1}{1,000,000}$インチの長さを持つ部分で構成されるかのように考えることができます；または、考えられる限りの極限まで推し進めると、無数の要素から成り立つと見なすことができます。これらの要素はすべて無限に小さく、無限小です。
はい、あなたは言うでしょうが、そのような考え方にはどのような用途がありますか？一度に全体を思い描くのはなぜなりませんか？単純な理由は、何かの大きさを直接計算することができない多くの場合があるからです。「積分」のプロセスは、それらの小さな部分の合計を計算することで、その全体の大きさを直接計算できないものを計算するのに役立ちます。
まず、この概念に慣れるためにいくつかの単純な例を考えてみましょう。
次のような数列を考えます：
\[
1 + \frac{1}{2} + \frac{1}{4} + \frac{1}{8}
  + \frac{1}{16} + \frac{1}{32} + \frac{1}{64} + \text{等}
\]

各項は前項の半分の値を持つように形成されます。無限に多くの項を考えた場合、その総和は何になりますか？学校で習う通り、答えは2です。もし好きなら、それを線として考えてみてください。最初に1インチから始め、半分の1インチを加え、四分の一インチを加え、八分の一インチを加えていきます。操作のどの時点でも中断したとしても、まだ全体の2インチを作るための欠けた部分がありますし、その欠けた部分は最後に追加された部分と同じ大きさです。例えば、1, $\frac{1}{2}$, そして $\frac{1}{4}$ を合計した後で中断すると、$\frac{1}{4}$ の欠けた部分があります。$\frac{1}{64}$ を加えた後でもまだ $\frac{1}{64}$ の欠けた部分があります。必要な余分は常に最後に追加された項と同じです。無限の操作によってのみ実際の2インチに達することができます。実際には、描画できないほど小さな部分まで到達した時点でほぼ達成されます——これには約10項程度が必要で、第11項は $\frac{1}{1024}$ です。 Whitworth の測定機械でも検出できないほど小さい部分まで行きたければ、約20項程度まで進める必要があります。マイクロスコープではすら18番目の項が見えません！したがって、無限の操作はそんなに恐ろしいものではありません。

積分は単にその全部を表します。しかし、後で見るように、積分微分法は指定された傾斜（または各点での傾斜）から全体の曲線を再構築するための逆過程を実行することができます。曲線の傾きについて少し前もって調査しましょう。なぜなら、曲線を微分することはその傾斜（または各点での傾斜）を見つけることと等しいからです。指定された傾斜（または各点での傾斜）が与えられた場合に全体の曲線を再構築することができるでしょうか？

Case (2) に戻り、\Pageref{Case2} を参照してください。ここでは最も単純な曲線、傾斜した直線があり、方程式は
\[
y = ax + b
\]
です。\par\textbf{〔第47図〕}\par

我们知道$b$代表当$x=0$时$y$的初始高度，而$a=\dfrac{dy}{dx}$是该直线的“斜率”。这条线具有恒定的斜率。整条线上，基本三角形（）的高度与底边的比例相同。

假設有有限大小的$dx$和$dy$，使得10个$dx$组成一英寸，則会有十个類似的小三角形（

  

现在假設我们仅从$\dfrac{dy}{dx}=a$的信息出发被要求重建“曲线”。我们该怎么办？仍然将这些小$d$视为有限大小，我们可以画出10个这样的三角形，并且它们具有相同的斜率，然后像这样连接起来：
\par\textbf{〔第48図〕}\par

由于所有斜率相同，它们会组合成一条如图\Fig{48}所示的倾斜线，其斜率为$\dfrac{dy}{dx}=a$。无论我们是将$dy$和$dx$视为有限大小还是无穷小，因为它们都一样，显然有$\dfrac{y}{x}=a$，如果我们将$y$视为所有$dy$之和，而$x$视为所有$dx$之和。但是，这条倾斜线应放置在何处？是从原点$O$开始，还是更高处？由于我们仅有关于斜率的信息，没有关于相対于$O$的具体高度的指示；事实上，初始高度是未确定的。无论初始高度如何，斜率都是相同的。

因此，让我们尝试可能需要的情况，并从高于$O$的高度$C$开始这条倾斜线。
即我们有方程
\[
y = ax + C.
\]

现在显而易见的是，在这种情况下，附加常数表示当$x=0$时$y$的具体值。

接下来考虑一个更难的问题，即斜率不恒定的直线，其斜率随$x$的增长而逐渐增大。假設上坡斜率随着$x$的增长而变得越来越大。用符号表示为：
\[
\dfrac{dy}{dx} = ax.
\]

或者，举个具体的例子，設$a=\dfrac{1}{5}$，則有
\[
\dfrac{dy}{dx} = \tfrac{1}{5} x.
\]

次に、異なる$x$の値における傾きのいくつかの値を計算し、それらを小さな図形で表すことが一番良いでしょう。

\begin{align*}
x &= 0 \quad \frac{dy}{dx} = 0, &  \\
x &= 1 \quad \frac{dy}{dx} = 0.2, &  \\
x &= 2 \quad \frac{dy}{dx} = 0.4, &  \\
x &= 3 \quad \frac{dy}{dx} = 0.6, &  \\
x &= 4 \quad \frac{dy}{dx} = 0.8, &  \\
x &= 5 \quad \frac{dy}{dx} = 1.0. & 
\end{align*}

これらの部分を組み合わせて、それぞれの基底の中央が適切な距離にあるように並べ、角で接続します（図49参照）。結果は当然、滑らかな曲線ではありませんが、その近似です。各部分が半分に短く、2倍多くなると（図50参照）、より良い近似になります。しかし、完全な曲線を得るには、それぞれの$dx$と対応する$dy$を無限に小さくし、無数にすべきです。

\par\textbf{〔第49図〕}\par

どの$y$の値がどれほど必要かを考えましょう。明らかに、曲線上の任意の点$P$における$y$の値は、$0$からそのレベルまでのすべての小さな$dy$の和であり、つまり$\displaystyle \int dy = y$です。そして各$dy$が$\frac{1}{5}x \cdot dx$に等しいので、全体の$y$は$\frac{1}{5}x \cdot dx$のような各部分の和であり、または書くと$\displaystyle \int \tfrac{1}{5}x \cdot dx$です。

もし$x$が一定であれば、$\displaystyle \int \tfrac{1}{5}x \cdot dx$は$\frac{1}{5} x \displaystyle \int dx$や$\frac{1}{5}x^2$と等しいでしょう。しかし$x$は最初に$0$であり、点$P$の特定の値まで増加するので、その間の平均値は$\frac{1}{2}x$です。したがって$\displaystyle \int \tfrac{1}{5} x \, dx = \tfrac{1}{10} x^2$であり、または$y=\tfrac{1}{10}x^2$です。

しかし、前の例と同様に、この式には未定の定数$C$を加える必要があります。なぜなら$x=0$での曲線の開始高度が与えられていないからです。したがって、図51で描かれた曲線の方程式は
\[
y = \tfrac{1}{10}x^2 + C.
\]
となります。

\par\textbf{〔第51図〕}\par

\par\textbf{【演習 XVI】}\par (See \Pageref[page]{AnsEx:XVI} for Answers.)

\begin{enumerate}
\item[(1)] $\frac{2}{3} + \frac{1}{3} + \frac{1}{6} + \frac{1}{12} + \frac{1}{24} + \text{etc}$ の極限和を求めよ。

\item[(2)] 級数 $1 - \frac{1}{2} + \frac{1}{3} - \frac{1}{4} + \frac{1}{5} - \frac{1}{6} + \frac{1}{7}$ など、収束することを示し、8項までの和を求めよ。

\item[(3)] $\log_\epsilon(1+x) = x - \dfrac{x^2}{2} + \dfrac{x^3}{3} - \dfrac{x^4}{4} + \text{etc}$ から $\log_\epsilon 1.3$ を求めよ。

\item[(4)] この章で説明したような論理を用いて、次を求めよ。
\[
\text{（a）$\frac{dy}{dx} = \frac{1}{4} x$ のとき；（b）$\frac{dy}{dx} = \cos x$ のとき}
\]

\item[(5)] $\dfrac{dy}{dx} = 2x + 3$ から $y$ を求めよ。

\end{enumerate}



\chapter{微分の逆演算としての積分計算}
\noindent\textit{\large Integrating as the Reverse of Differentiating}\\[
2mm]
\noindent\textbf{\small 【章の論理とテーマ】基本関数の積分公式・多項式・分数式・べき関数の積分}\\[4mm]
XVIII 適分法とは
微分法は、$y$が$x$の関数として与えられたとき、$\dfrac{dy}{dx}$を求めることです。
逆差分化分法と同様に、数学的な操作全体が逆転する可能性があります。例えば、$y = x^4$を微分すると$\dfrac{dy}{dx} = 4x^3$となりますが、逆に$\dfrac{dy}{dx} = 4x^3$から始めると$y = x^4$となります。しかし、ここで奇妙な点が現れます。もし$y = x^4$, $y = x^4 + a$, $y = x^4 + c$, または任意の定数を加えた$x^4$のいずれかから始めても$\dfrac{dy}{dx} = 4x^3$となります。したがって、$\dfrac{dy}{dx}$から$y$に戻る際に追加された定数があることを考慮する必要があります。その値は他の方法で特定されるまで未確定です。
よって、$x^n$を微分すると$nx^{n-1}$となるので、逆に$\dfrac{dy}{dx} = nx^{n-1}$から始めると$ = x^n$ + C$となります。ここで$C$はまだ未定義の可能性のある定数です。
$x$の冪に関しては、逆方向でのルールは次のようになります：指数を1増やし、その増加した指数で割り、未定義の定数を追加します。
例えば
\[
\dfrac{dy}{dx} = x^n
\]
の場合、逆方向では
\[
y = \dfrac{1}{n + 1} x^{n+1} + C
\]
となります。

また、$y = ax^n$を微分すると$\dfrac{dy}{dx} = anx^{n-1}$となるので、逆に$\dfrac{dy}{dx} = anx^{n-1}$から始めると$y = ax^n$となります。したがって、乗数がある場合、積分の結果にその定数を乗数として加える必要があります。

例えば、$\dfrac{dy}{dx} = 4x^2$の場合、逆方向では
\[
y = \dfrac{4}{3} x^3
\]
となります。しかし、これは不完全です。もし最初から
\[
y = ax^n + C
\]
（$C$は任意の定数）を始めても同様に$\dfrac{dy}{dx} = anx^{n-1}$となることを忘れないでください。

したがって、逆方向での操作を行う際には常に未定義の定数を追加する必要があります。その値はまだ特定されていないとしてもです。

この過程は微分の逆であるため、\emph{積分}と呼ばれる。与えられた \( dy \) または \( \dfrac{dy}{dx} \) の表現から全量 \( y \) の値を求める作業である。これまで可能であれば \( dy \) と \( dx \) を微分係数として一緒に扱ってきましたが、これからはそれらを分離することが多くなる。

単純な場合から始めましょう：
\[
\dfrac{dy}{dx} = x^2.
\]

この式を書き換えると、
\[
dy = x^2 \, dx.
\]

これは「微分方程式」であり、\( y \) の要素が \( x \) の要素に \( x^2 \) を掛けて得られるという情報を提供します。私たちが求めたいのは積分です。よって、適切な記号を使用して両辺を積分する指示を書き下します：
\[
\int dy = \int x^2 \, dx.
\]

積分の読み方についての注記：上記は以下のようになります：
```
「Integral dee-wy equals integral eks-squared dee-eks」
```

まだ積分していません。ただ、積分するための指示を書き下しただけです---もし可能であれば。試してみましょう。他の多くの愚かな者はできるので、なぜ私たちはできないのでしょうか？左辺は非常に単純です。すべての \( y \) の要素の合計は \( y \) そのものと等しいからです。よって、すぐに以下のように書けます：
\[
y = \int x^2 \, dx.
\]

しかし、方程式の右辺に到達すると、求積すべき項はすべての$dx$ではなく、$x^2\, dx$のような項であることに注意しなければなりません。そしてこれは$x^2 \displaystyle \int dx$とは異なるものになります。なぜなら$x^2$は定数ではないからです。いくつかの$dx$は大きな$x^2$の値で乗算され、他のいくつかは小さな$x^2$の値で乗算されます。$x$がどのような値を取るかによります。そこで積分の逆過程である微分の逆過程について考え直す必要があります。つまり、$x^n$の場合の規則（\Pageref{revdif}参照）は「指数を1増加させ、その増加した指数と同じ数で割る」というものです。（）

つまり、$x^2\, dx$は変換されて$\frac{1}{3} x^3$になります（末尾の小さな$dx$がどこにいったのかと疑問を呈するかもしれません。実際には微分係数の一部であり、$x^2\, dx$のように右辺に移動すると$x$が操作を行う独立変数であることを思い出させる役割を果たします。そして積分された結果として、$x$の指数は1つ増加します）。すぐにこの手順に慣れることでしょう。

これを方程式に代入し、積分定数$C$も忘れずに追加します。すると次のようになります：
\[
y = \frac{1}{3} x^3 + C.
\]

あなたは実際には積分を完了しました。それほど簡単ではありませんか？

もう一つの単純な例を見てみましょう。

\begin{align*}
\dfrac{dy}{dx} &= ax^{12},
\end{align*}

ここで$a$は任意の定数倍です。微分するとき（\Pageref{differ}参照）で確認したように、$y$の値に含まれる任意の定数因子は$\frac{dy}{dx}$の値にも変化なく現れます。逆過程である積分においても同じことが起こりますので、前の手順と同じように作業を進めることができます。

\begin{align*}
dy &= ax^{12} \, dx,\\
\int dy &= \int ax^{12} \, dx,\\
\int dy &= a \int x^{12} \, dx,\\
y &= a \times  \frac{1}{13} x^{13} + C.
\end{align*}

これで完了です。それほど簡単ではありませんか？

現在、積分は微分と比較して「戻り道を見つける」プロセスであることを理解するようになります。微分の過程で特定の表現を見つけた場合（例えば $ax^{12}$）から元の $y$ に戻る方法を見つけることができます。二つのプロセスの対比は、有名な教師による以下のような言葉で例示できます。トラファルガー広場に迷路入りした Stranger がエストン駅まで行くのは困難かもしれませんが、エストン駅からトラファルガー広場までの道を実際に案内された場合、彼には比較的容易にその逆の道を見つけることができます。

\section*{二つの関数の和または差の積分}

\begin{align*}
\lintertext{\indent 令}
\frac{dy}{dx} &= x^2 + x^3, \\
\lintertext{則}
dy &= x^2\, dx + x^3\, dx.
\end{align*}

各項を個別に積分することは理由がないわけではありません：二つの独立な関数の和を微分した場合、導関数はその二つの独立な微分の和であることが \Pageref{sumdiffer} で示されていますからです。逆方向に戻るとき、積分を行うと、それは二つの独立な積分の和になります。

その後の指示は以下の通りです：
\begin{align*}
\int dy
  &= \int (x^2 + x^3)\, dx \\
  &= \int x^2\, dx + \int x^3\, dx   \\
y &= \frac{1}{3} x^3 + \frac{1}{4} x^4 + C.
\end{align*}

もしどちらかの項が負の量であれば、積分式における対応する項も負になります。したがって、和と差は同様に扱えます。

\section*{定数項の処理方法}

被積分式の中に定数項がある場合（例えば $ = x^n$ + b$）を考えてみましょう。

これは笑えるほど簡単です。$y = ax$ の微分が $\dfrac{dy}{dx} = a$ であることを思い出すだけです。したがって、逆方向に積分を行うと定数は $x$ で乗算されます。従って、
\begin{align*}
dy &= x^n\, dx + b \cdot dx,  \\
\int dy &= \int x^n\, dx + \int b\, dx, \\
y &= \frac{1}{n+1} x^{n+1} + bx + C.
\end{align*}

あなたの新しい能力を試すための多くの例題があります。

\tb

\par\medskip\noindent\textbf{【例題】}\par
(1) $\dfrac{dy}{dx} = 24x^{11}$ が与えられたとき、$y$ を求めよ。\qquad \textit{答}.\ $y = 2x^{12} + C$.

(2) $$\int (a + b)(x + 1)\, dx$$ を求めよ。これは $(a + b) \int (x + 1)\, dx$ または $$(a + b) \left[\int x\, dx + \int dx\right]$$ または $$(a + b) \left(\frac{x^2}{2} + x\right) + C$$ である。

(3) $$\frac{du}{dt} = gt^{\frac{1}{2}}$$ が与えられたとき、$u$ を求めよ。答えは $u = \frac{2}{3} gt^{\frac{3}{2}} + C$ である。

(4) $$\frac{dy}{dx} = x^3 - x^2 + x$$ から $y$ を求めよ。
\begin{align*}
dy &= (x^3 - x^2 + x)\, dx \quad \text{または} \\
dy &= x^3\, dx - x^2\, dx + x\, dx; \quad
y = \int x^3\, dx - \int x^2\, dx + \int x\, dx;
\end{align*}
したがって
\[
 y = \frac{1}{4} x^4 - \frac{1}{3} x^3 + \frac{1}{2} x^2 + C. 
\]

(5) $$9.75x^{2.25}\, dx$$ を積分せよ。答えは $y = 3x^{3.25} + C$ である。

これらはすべて比較的簡単である。別の例を試してみよう。
最简单的积分

\begin{align*}
\lintertext{令}
\frac{dy}{dx} &= ax^{-1}.
\end{align*}

前と同じように進めると
\[
 dy = a x^{-1} \cdot dx, \quad \int dy = a \int x^{-1}\, dx. 
\]

しかし、$x^{-1}\, dx$ の積分は一体どうなるか？

$x^2$, $x^3$, そして $x^n$ 等の微分結果を振り返ってみれば、どの場合も $\frac{dy}{dx}$ の値として $x^{-1}$ を得ることはなかった。$x^3$ から $3x^2$，$x^2$ から $2x$，そして $x^1$（つまり $x$ 自身）から $1$ を得たが、$x^0$ からは $x^{-1}$ を得ることはなかった。二つの理由がある。まず第一に、$x^0$ は単に $= 1$ であり、定数であるため微分係数を持てない。第二に、もし微分可能であれば、通常のルールを忠実に適用しても微分係数は $0 \times x^{-1}$ となり、ゼロと乗算するとその値もゼロになる！したがって、現在 $x^{-1}\, dx$ の積分を考えると、それは $x$ のべき乗のルールによって与えられるべき式のどこにも含まれないことがわかる：
\[
 \int x^n\, dx = \frac{1}{n+1} x^{n+1}. 
\]
これは例外的なケースである。

しかし、もう一度試してみよう。様々な $x$ の関数から得られた微分をすべて見渡し、その中に $x^{-1}$ が含まれているか探す。十分な調査は、実際には $y = \log_\epsilon x$（参照：\pageref{differlog}）の微分結果として $\frac{dy}{dx} = x^{-1}$ を得たことを示すだろう。

次に、$\log_\epsilon x$ を微分すると $x^{-1}$ となることを知っているので、逆の過程で $dy = x^{-1}\, dx$ を積分すると $y = \log_\epsilon x$ となります。しかし、与えられた定数因子 $a$ を忘れず、不定積分に加える未定の積分定数も忘れないことが重要です。これにより現在の問題の解は
\[
y = a \log_\epsilon x + C.
\]

ここで注意すべき非常に興味深い事実は、上記の場合のように $x^{-1}\, dx$ の積分ができなかったら、対応する微分が分かっていなければならなかったことです。もし $\log_\epsilon x$ を微分すると $x^{-1}$ となることが誰も発見していなければ、$x^{-1}\, dx$ の積分問題は全く手詰まりになっていたでしょう。実際、微分積分学の奇妙な特徴の一つとして認められるべき点は、何かを微分した結果が求める積分式であることを先に知らなければ、何ものも積分できないことです。今日でも、一般形の積分
\[
\frac{dy}{dx} = a^{-x^2}
\]
を見つけ出す人は誰一人いません。なぜなら $a^{-x^2}$ は何かを微分した結果としてまだ見つかっていないからです。

別の簡単な例題。
$(x + 1)(x + 2)\, dx$ を積分せよ。

被積分関数を見ると、これは $x$ の異なる2つの関数の積であることがわかります。それぞれ $(x + 1)\, dx$ や $(x + 2)\, dx$ を独立に積分することはできます。もちろんできます。しかし、積をどうするか？学んだ微分法から得られる導函数はこのような積ではありません。そのようなものは見つからない場合、最も簡単なことは2つの関数を乗算し、その後積分することです。これにより
\[
\int (x^2 + 3x + 2)\, dx.
\]
となり、これは
\[
\int x^2\, dx + \int 3x\, dx + \int 2\, dx.
\]
と等しいことがわかります。各積分を計算すると
\[
\tfrac{1}{3} x^3 + \tfrac{3}{2} x^2 + 2x + C.
\]

現在、積分は微分の逆であることを知っているので、既に知る微分係数を調べて、それらはどの関数から導かれたものかを見ることができます。これにより、次の積分が得られます：
\begin{alignat*}{4}
&x^{-1} &&(\Pageref{differlog}); \qquad &&
  \int x^{-1}\, dx      &&= \log_\epsilon x + C. \\
%
&\frac{1}{x+a} &&(\Pageref{differ2}); &&
  \int \frac{1}{x+a}\, dx &&= \log_\epsilon (x+a) + C. \\
%
&\epsilon^x &&(\Pageref{unchanged}); &&
  \int \epsilon^x\, dx    &&= \epsilon ^x + C. \\
%
&\epsilon^{-x} &&&&
  \int \epsilon^{-x}\, dx &&= -\epsilon^{-x} + C \\
%
\intertext{(for if $y = - \dfrac{1}{\epsilon^x}$,\quad $\dfrac{dy}{dx} = -\dfrac{\epsilon^x \times  0 - 1 \times  \epsilon^x}{\epsilon^{2x}} = \epsilon^{-x}$).}
%
&\sin x &&(\Pageref{differcos}); &&
  \int \sin x\, dx        &&= -\cos x + C. \\
%
&\cos x &&(\Pageref{differsin}); &&
  \int \cos x\, dx        &&= \sin x + C. \\
%
\intertext{\indent Also we may deduce the following:}
%
&\log_\epsilon x; &&&&
  \int\log_\epsilon x\, dx &&= x(\log_\epsilon x - 1) + C \\
%
\intertext{(for if $y = x \log_\epsilon x - x$,\quad $\dfrac{dy}{dx} = \dfrac{x}{x} + \log_\epsilon x - 1 = \log_\epsilon x$).}
%
&\log_{10} x;   &&&&
  \int\log_{10} x\, dx &&= 0.4343x (\log_\epsilon x - 1) + C. \\
%
&a^x &&(\Pageref{diffexp}); &&
  \int a^x\, dx        &&= \dfrac{a^x}{\log_\epsilon a} + C. \\
%
&\cos ax; &&&&
  \int\cos ax\, dx     &&= \frac{1}{a} \sin ax + C \\
\intertext{(for if $y = \sin ax$, $\dfrac{dy}{dx} = a \cos ax$; hence to get $\cos ax$
one must differentiate $y = \dfrac{1}{a} \sin ax$).}
%
&\sin ax; &&&&
  \int\sin ax\, dx     &&= -\frac{1}{a} \cos ax + C. \\
\end{alignat*}

試してみて$\cos^2\theta$; 少しの手口で計算が簡素化されます：
\begin{align*}
\cos 2\theta &= \cos^2\theta - \sin^2\theta = 2\cos^2\theta - 1; \\
\therefore \quad \cos^2\theta &= \frac{1}{2}(\cos 2\theta + 1),
\end{align*}
したがって
\begin{align*}
\int \cos^2 \theta\, d\theta &= \frac{1}{2} \int (\cos 2\theta + 1)\, d\theta \\
&= \frac{1}{2} \int \cos 2 \theta\, d\theta + \frac{1}{2} \int d\theta \\
&= \frac{\sin 2\theta}{4} + \frac{\theta}{2} + C. \quad (\text{参照例：}\pageref{moreexamples})
\end{align*}

表の標準形を参照してください（ページ範囲：\pageref{stdforms1}から\pageref{stdforms2}）。
あなた自身でそのような表を作成し、成功した微分と積分を行う一般的な関数だけを記入することをお勧めします。この表が徐々に成長することを確認してください。

二重積分や三重積分について（複数の変数を含む式の積分）。
しばしば、その中にあるいくつかの関数を2つ以上の変数について積分する必要があります；その場合、積分記号は複数回現れます。例えば，
\[
\iint f(x,y)\, dx\, dy
\]
これは、$x$と$y$のいずれかの変数についてそれぞれ積分することを意味します。どの順序で行うかは関係ありません。例えば、関数$x^2 + y^2$を考えます。この関数を$x$について積分すると：
\[
\int (x^2+y^2)\, dx = \frac{1}{3} x^3 + xy^2.
\]

次に、この結果を$y$について積分します：
\[
\int (\frac{1}{3} x^3 + xy^2)\, dy = \frac{1}{3} x^3y + \frac{1}{3} xy^3,
\]
もちろん、定数も加えます。もし積分の順序を逆にしたとしても、結果は同じになります。

面積や体積を考える際に、しばしば長さと幅について積分する必要があります。そのような形の積分は
\[
\iint u \cdot dx\, dy,
\]
ここで$u$は各点で$x$と$y$に依存する何かの性質です。この積分は「面積積分」と呼ばれます。これは、小さな長方形$dx$$\times$$dy$の各点での$u$の値を全長と全幅にわたって合計することを示します。

相似の原理は、三次元の立体の場合にも適用されます。例えば、任意の体積要素を考え、そのサイズが $dx$\textasciitilde{}$dy$\textasciitilde{}$dz$ の小さな立方体です。立体の形状を関数 $f(x, y, z)$ で表すと、全体の立体の体積は \emph{体積積分} として与えられます：
\[
\text{体積} = \iiint f(x,y,z) \cdot dx \cdot dy \cdot dz.
\]
当然、各次元において適切な境界値間に積分を取る必要があります\footnote
  {境界値間での積分については \Pageref{limits} を参照してください。}
；そして、境界が $x$，$y$，および $z$ のどの値に依存するかを知っていることが積分を行う上で必要です。例えば、$x$ の境界は $x_1$ から $x_2$，$y$ の境界は $y_1$ から $y_2$，そして $z$ の境界は $z_1$ から $z_2$ であるとすると、明らかに
\[
\text{体積} = \int_{z1}^{z2} \int_{y1}^{y2} \int_{x1}^{x2} f(x,y,z) \cdot dx \cdot dy \cdot dz.
\]
当然、複雑な場合もたくさんありますが、一般的には、特定の積分が与えられた表面全体または与えられた立体空間全体で行われることを示すシンボルの意味はすぐに理解できます。

単純な積分の練習問題（解答は \Pageref{AnsEx:XVII} を参照）：

\begin{enumerate}
\item (1) $y^2 = 4ax$ のとき $\displaystyle \int y\, dx$ を求めよ。
\item (2) $\displaystyle \int \frac{3}{x^4}\, dx$ を求めよ。
\item (3) $\displaystyle \int \frac{1}{a} x^3\, dx$ を求めよ。
\item (4) $\displaystyle \int (x^2 + a)\, dx$ を求めよ。
\item (5) $5x^{-\frac{7}{2}}$ の積分を求めよ。
\item (6) $\displaystyle \int (4x^3 + 3x^2 + 2x + 1)\, dx$ を求めよ。
\item (7) $\dfrac{dy}{dx} = \dfrac{ax}{2} + \dfrac{bx^2}{3} + \dfrac{cx^3}{4}$ のとき $y$ を求めよ。
\item (8) $\displaystyle \int \left(\frac{x^2 + a}{x + a}\right) dx$ を求めよ。
\item (9) $\displaystyle \int (x + 3)^3\, dx$ を求めよ。
\item (10) $\displaystyle \int (x + 2)(x - a)\, dx$ を求めよ。
\item (11) $\displaystyle \int (\sqrt x + \sqrt[3] x) 3a^2\, dx$ を求めよ。
\item (12) $\displaystyle \int (\sin \theta - \tfrac{1}{2})\, \frac{d\theta}{3}$ を求めよ。
\item (13) $\displaystyle \int \cos^2 a \theta\, d\theta$ を求めよ。
\item (14) $\displaystyle \int \sin^2 \theta\, d\theta$ を求めよ。
\item (15) $\displaystyle \int \sin^2 a \theta\, d\theta$ を求めよ。
\item (16) $\displaystyle \int \epsilon^{3x}\, dx$ を求めよ。
\item (17) $\displaystyle \int \dfrac{dx}{1 + x}$ を求めよ。
\item (18) $\displaystyle \int \dfrac{dx}{1 - x}$ を求めよ。

\end{enumerate}



\chapter{積分による面積の求値法（定積分）}
\noindent\textit{\large On Finding Areas by Integrating}\\[
2mm]
\noindent\textbf{\small 【章の論理とテーマ】定積分・微小区間の無限和・曲線の下の面積の厳密計算}\\[4mm]
積分計算の一つの用途は、曲線で囲まれた面積の値を求めることが可能です。
まず少しずつこの主題について考えてみましょう。
![2.5インチ](218a){\#fig:52}
曲線 $AB$（図\[52
\]）があり、その方程式が既知です。つまり、この曲線上の $y$ は $x$ のある已知の関数です。点 $P$ から点 $Q$ までの曲線の一部を考えます。

点 $P$ に垂直な線 $PM$ を引き、点 $Q$ に垂直な線 $QN$ も引きます。すると $OM = x_1$ と $ON = x_2$ とし、縦座標 $PM = y_1$ と $QN = y_2$ とします。これで曲線 $PQ$ の下側に位置する面積 $PQNM$ を定義しました。問題は、この面積の値をどのように計算できるかです？

各幅が $dx$ の非常に狭い帯状部分に分割することで、この問題を解決する秘訣があります。$dx$ が小さければ小さいほど、$x_1$ と $x_2$ 間に含まれる帯状部分の数が増えます。明らかに全体の面積はこれらの各帯状部分の面積の合計に等しいです。そこで、一つの狭い帯状部分の面積を求める式を見つけ出し、それら全てを加算する必要があります。どの一つの帯状部分を考えても、その形状は二つの垂直な側面で囲まれたものであり、底辺が $dx$ で、上端がやや曲がった平らなものとします。この帯状部分の平均高さを $y$ として考えると、その幅が $dx$ であるため、その面積は $y \, dx$ となります。そして、幅をどれだけ狭くしても良いので、十分に狭ければ平均高さは中点での高さと等しくなります。全体の面積を未知数として $S$ と呼び、一つの帯状部分の面積を $dS$ と呼ぶことにします。したがって、

\[
1\text{つの帯状部分の面積} = dS = y \, dx
\]

となります。これら全ての帯状部分を加算すると、

\[
全体の面積 S = \int dS = \int y \, dx
\]

となります。

したがって、$S$ を求めることは、特定の場合において $y$ が $x$ の関数として与えられているとき、$y \, dx$ を積分することができるかどうかに依存します。

例えば、曲線 \( y = b + ax^2 \) について与えられた場合、おそらくあなたはその値を式に代入し、次のように言えるでしょう：「それでは、\(\int (b + ax^2) \, dx\) を求めなければなりません。」 %[ ** \displaystyle]

これは十分ですが、少し考えると、もっとすべきことが明らかになります。なぜなら、探している面積は曲線の全体ではなく、左側で \( PM \) に制限され、右側で \( QN \) に制限される面積であるからです。したがって、私たちにはその間の「制限」を定義する必要があります。「制限」 %[ ** F2: 単引用符を使用しています。他は双引用符）

この文脈で新しい概念が導入されます：「境界における積分」というものです。\( x \) が変化すると仮定し、現在の目的では \( OM \) の下にある任意の \( x \) の値（すなわち \( x_1 \)) や \( ON \) の上の任意の \( x \) の値（すなわち \( x_2 \)) は必要としません。積分が二つの境界間で定義されるとき、下の値を「 inferior limit （ Inferior 境限）」、上の値を「 superior limit （ Superior 境限）」と呼びます。これらの制限内で定義された任意の積分は、「 definite integral （ Definite 確定積分）」と呼んで区別します。

積分を行う指示記号において、境界は積分記号の上端と下端にマークされます。例えば、次の指示
\[
\int_{x=x_1}^{x=x_2} y \, dx
\]
は「 inferior limit が \( x_1 \) から superior limit が \( x_2 \) までの \( y \, dx \) の積分を求めよ」と読み解きます。

時には、このことをより単純に書き表すこともできます：
\[
\int_{x_1}^{x_2} y \, dx.
\]
では、これらの指示を得て境界における積分をどのように求めますか？

再び図 Fig{52}（\Pageref{fig:52}）を見てみましょう。もし \( A \) から \( Q \) の大きな曲線部分の下側の面積、すなわち \( x = 0 \) から \( x = x_2 \) の面積 \( AQNO \) を求めることができた場合、そして \( A \) から \( P \) の小さな曲線部分の下側の面積、すなわち \( x = 0 \) から \( x = x_1 \) の面積 \( APMO \) を求めることができた場合、これらを引くと残る余分な面積は \( PQNM \) となり、これが探している面積です。ここに私たちがすべきことを示すヒントがあります：二つの境界間の確定積分は、 superior limit で計算した積分から inferior limit で計算した積分を引いた「差」です。

次に進みましょう。まず、一般解を求めるために次の積分を求めます：
\[
\int y\, dx,
\]
そして、曲線（図52）の方程式 \(y = b + ax^2\) から、
\[
\int (b + ax^2)\, dx
\]
が求め trebuieする一般解です。

この積分をルール（第9節参照）で計算すると、次のようになります：
\[
bx + \frac{a}{3} x^3 + C;
\]
そしてこれは、任意の \(x\) の値まで全域積分となるため、
\(0\) から任意の \(x\) 値までの面積は
\[
bx_2 + \frac{a}{3} x_2^3 + C;
\]
また、下限 \(x_1\) までの面積は
\[
bx_1 + \frac{a}{3} x_1^3 + C.
\]

次に、大きい面積から小さい面積を引くと、求める面積 \(S\) の値は次のようになります：
\[
\text{面積~$S$} = b(x_2 - x_1) + \frac{a}{3}(x_2^3 - x_1^3).
\]

これが求めた答えです。具体的な数値を代入してみましょう。例えば、\(b = 10\)、\(a = 0.06\)、\(x_2 = 8\)、\(x_1 = 6\) のとき、面積 \(S\) は次のようになります：
\begin{gather*}
10(8 - 6) + \frac{0.06}{3} (8^3 - 6^3) \\
\begin{aligned}
&= 20 + 0.02(512 - 216)    \\
&= 20 + 0.02 \times  296    \\
&= 20 + 5.92     \\
&= 25.92.
\end{aligned}
\end{gather*}

ここで、限界についての結論を記述するためのシンボリックな方法を以下に示します：
\[
\int^{x=x_2}_{x=x_1} y\, dx = y_2 - y_1,
\]
ここで \(y_2\) は \(x_2\) に対応する積分値、そして \(y_1\) は \(x_1\) に対応する積分値です。

限界間のすべての積分はこのような差を求めることで求められます。また、引数を計算する際に追加された定数 \(C\) が消えます。

例題．
(1) 前もって答えを知っている問題で慣れてみましょう。三角形（図53）の面積を求めます。底辺は \(x = 12\)、高さは \(y = 4\) です。明らかに測量から、答えは \(24\) になります。

図53

ここで、「曲線」は傾斜した直線で、その方程式は
\[
y = \frac{x}{3}
\]
となります。

求める面積は次のようになります：
\[
\int^{x=12}_{x=0} y · dx = \int^{x=12}_{x=0} \frac{x}{3} · dx.
\]

積分 $\frac{x}{3}\, dx$（Pageref{diffrule}）を行います。一般解の値を四角括弧内に書き、上部と下部に限界を記します。すると次のようになります。

\begin{align*}
面積 &= \left[ \frac{1}{3} · \frac{1}{2} x^2 \right]^{x=12}_{x=0} + C \\
  &= \left[ \frac{x^2}{6} \right]^{x=12}_{x=0} + C \\
  &= \left[ \frac{12^2}{6} \right] - \left[ \frac{0^2}{6} \right] \\
  &= \frac{144}{6} = 24.\quad \textit{答}.
\end{align*}

この計算法が意外と有用であることを確認するために、単純な例を使って試してみましょう。八分格子紙や十分格子紙を用意し、方程式 $y = \frac{x}{3}$ のグラフを描きます。

座標の値は以下の通りです：
\[
\begin{array} {|c|| *{5}{c|}}
\hline
x & 0 & 3 & 6 & 9 & 12 \\
\hline
y & 0 & 1 & 2 & 3 & 4 \\
\hline
\end{array}
\]

このグラフは \Fig{54} に示されています。

次に、$x = 0$ から $x = 12$ の範囲で曲線の下側にある面積を数格子法で計算します。完全な格子が 18 個あり、さらに 4 個の三角形があり、それぞれの三角形の面積は 1.5 格子分です。合計では 24 格子分となります。したがって、$\frac{x}{3}\, dx$ の $x = 0$ から $x = 12$ までの定積分の数値は 24 です。

さらに、$x = 3$ から $x = 15$ の範囲での同様な積分の値が 36 であることを示してください。

\begin{align*}
面積
  &= \int^{x=x_1}_{x=0} y · dx \\
   &= \int^{x=x_1}_{x=0} \frac{b}{x+a}\, dx  \\
  &= b \left[\log_\epsilon(x + a) \right]^{x_1}_{0} + C  \\
  &= b \left[\log_\epsilon(x_1 + a) - \log_\epsilon(a)\right]  \\
  &= b \log_\epsilon \frac{x_1 + a}{a}.\quad \textit{答}.
\end{align*}

注意すべき点は、定積分の場合、定数 $C$ は引算によって常に消えてしまうことです。

この過程は、大きなものから小さなものを引いて差を求める方法であり、一般的な手法です。平面の輪（図56）の面積を求めたい場合、外径が \(r_2\) で内径が \(r_1\) の輪の面積をどのように求めますか？測量から分かることは、外側の円の面積は \(\pi r_2^2\) であり、次に内側の円の面積 \(\pi r_1^2\) を求め、それを引いて環の面積 \(= \pi(r_2^2 - r_1^2)\) を得ます。これはさらに
\[
\pi(r_2 + r_1)(r_2 - r_1)
\]
と書き換えることができ、これは「輪の平均円周」$\times$「輪の幅」と等しいです。

(3) また別の例は、衰減曲線（節5参照）です。この曲線（図57）の方程式は
\[
y = b \epsilon^{-x}
\]
であり、その間の \(x = 0\) から \(x = a\) の面積を求めます。

積分（3節参照）により
\begin{align*}
\text{面積} &= b \int_{x=0}^{x=a} \epsilon^{-x} \, dx \\
&= b \left[ -\epsilon^{-x} \right]_0^a \\
&= b \left( -\epsilon^{-a} - (-\epsilon^{-0}) \right) \\
&= b (1 - \epsilon^{-a}).
\end{align*}

図57, 58

(4) また別の例は、完璧な気体の等音性曲線（図58）で与えられます。その方程式は \(p v^n = c\) であり、ここで \(p\) は圧力、\(v\) は体積、そして \(n\) の値は1.42です。

突然気体を圧縮する際にした作業量に比例する曲線の下側の面積（体積 \(v_2\) から体積 \(v_1\) まで）を求めます。

ここでは
\begin{align*}
\text{面積} &= \int_{v=v_1}^{v=v_2} c v^{-n} \, dv \\
&= c \left[ \frac{1}{1-n} v^{1-n} \right]_{v_1}^{v_2} \\
&= c \frac{1}{1-n} (v_2^{1-n} - v_1^{1-n}) \\
&= \frac{-c}{0.42} \left( \frac{1}{v_2^{0.42}} - \frac{1}{v_1^{0.42}} \right).
\end{align*}

演習問題。
円の面積公式を証明し、半径 \(R\) の円の面積 \(A\) は \(\pi R^2\) に等しいことを示します。

図59

中心から半径 \(r\) の位置にある幅 \(dr\) の要素的な帯状領域（図59）を考えます。全体の表面はこのような狭い帯状領域で構成され、全体の面積 \(A\) は中心から縁までそのような要素的な帯状領域の積分であり、つまり \(r = 0\) から \(r = R\) までの積分です。

したがって、狭い帯の原始面積 \(dA\) を求めなければなりません。これを幅 \(dr\) の帯とみなし、半径 \(r\) の円の周長である長さ \(2 \pi r\) と考えます。すると、狭い帯の面積は
\[
dA = 2 \pi r \, dr
\]
となります。

したがって、円全体の面積 \(A\) は：
\[
A = \int dA
  = \int_{r=0}^{r=R} 2 \pi r \cdot dr
  = 2 \pi \int_{r=0}^{r=R} r \, dr.
\]

今、一般積分 \(r \, dr\) の結果は \(\frac{1}{2} r^2\) です。従って，
\begin{align*}
A &= 2 \pi \left[ \frac{1}{2} r^2 \right]_{r=0}^{r=R}; \\
\lintertext{または}
A &= 2 \pi \left[ \frac{1}{2} R^2 - \frac{1}{2}(0)^2 \right]; \\
\lintertext{従って}
A &= \pi R^2.
\end{align*}

\subsection{別の練習。}
曲線 \(y = x - x^2\) の正の部分の平均値を求めてみましょう（図60参照）。
図60はこの曲線を示しています。

平均値を求めることには、まず領域 \(OMN\) の面積を求め、その上で基底の長さ \(ON\) で割る必要があります。しかし、面積を求めることができる前に、基底の長さを求めて、積分の限界を知る必要があります。点 \(N\) では \(y = 0\) なので、方程式から \(x\) の値を見つけて \(y = 0\) となる \(x\) を求めます。明らかに、\(x = 0\) のとき \(y = 0\) であり、曲線は原点 \(O\) を通りますが、また \(x = 1\) のときも \(y = 0\) なので、点 \(N\) の位置は \(x = 1\) です。

求めたい面積は
\begin{align*}
  &= \int_{x=0}^{x=1} (x - x^2) \, dx \\
  &= \left[ \frac{1}{2} x^2 - \frac{1}{3} x^3 \right]_0^1 \\
  &= \left( \frac{1}{2} - \frac{1}{3} \right) - (0 - 0) \\
  &= \frac{1}{6}.
\end{align*}

しかし、基底の長さは \(1\) です。

したがって、曲線の平均値 \(= \frac{1}{6}\) です。

[\NB---最大値と最小値に関する微分法による問題として、この曲線の最大値の高さを求めることは美しい簡単な練習になります。それは平均よりも大きいはずです。]

任意の曲線 \(x = 0\) から \(x = x_1\) の範囲での平均 \(y\) 値は以下の式で与えられます：
\[
\text{平均~$y$} = \frac{1}{x_1} \int_{x=0}^{x=x_1} y \, dx.
\]

同様に、回転体の表面積も求めることができます。

\textbf{【例題】} 曲線 \(y = x^2 - 5\) を \(x\) 軸を中心に回転させます。\(x = 0\) から \(x = 6\) の間で生成される面積を求めましょう。

曲線上の点、その縦座標が \( y \) である点は、長さ \( 2\pi y \) の円周を描きます。この点に対応する幅 \( dx \) の狭い帯面の面積は \( 2\pi y\, dx \) です。全体の面積は
\[
2\pi \int^{x=6}_{x=0} y\, dx = 2\pi \int^{x=6}_{x=0} (x^2-5)\, dx = 2\pi \left[\frac{x^3}{3} - 5x\right]^6_0 = 6.28 \times  42 = 263.76
\]
となります。

極座標での面積.

領域の境界の方程式が、固定点 \( O \)（pole）からの距離 \( r \) と、\( r \) が正の水平方向 \( OX \) との角度 \( \theta \) として与えられているとき（図61参照），先に説明した過程は、わずかな修正を加えることで同じように適用できます。面積の帯ではなく、小さな三角形 \( OAB \) を考慮し、頂点 \( O \) の角度が \( d\theta \) であることに注意します。必要な領域を構成するすべての小さな三角形の和を求めます。

このような小さな三角形の面積は大体
\[
\frac{1}{2} AB \times  r = \frac{1}{2} r\, d\theta \times  r
\]
となります；したがって、曲線と角度 \( \theta_1 \) および \( \theta_2 \) に対応する \( r \) の二つの位置間の領域の部分は以下の積分で与えられます：
\[
\frac{1}{2} \int^{\theta=\theta_2}_{\theta=\theta_1} r^2\, d\theta.
\]

例題.

(1) 半径 \( a \) の円周上の 1 ラジアンの扇形の面積を求めよ。

極座標方程式は明らかに \( r = a \) である。面積は
\[
\frac{1}{2} \int^{\theta=\theta_2}_{\theta=\theta_1} a^2\, d\theta = \frac{a^2}{2} \int^{\theta=1}_{\theta=0} d\theta = \frac{a^2}{2}
\]
である。

(2) \( r = a(1 + \cos \theta) \) という極座標方程式を持つ「パスカルの螺線」の第1象限の面積を求めよ。
\[
\text{面積} = \frac{1}{2} \int^{\theta=\frac{\pi}{2}}_{\theta=0} a^2(1 + \cos \theta)^2\, d\theta
  = \frac{a^2}{2} \int^{\theta=\frac{\pi}{2}}_{\theta=0} (1 + 2 \cos \theta + \cos^2 \theta)\, d\theta
  = \frac{a^2}{2} \left[\theta + 2 \sin \theta + \frac{\theta}{2} + \frac{\sin 2 \theta}{4}\right]^{\frac{\pi}{2}}_{0}
  = \frac{a^2(3\pi+8)}{8}.
\]

積分による体積.

小さな立体の断面積に対して行ったことを、もちろん球体の体積についても同様に行います。全体の小さな断面積を加算し、その体積を求めることができます。これは、最終的な図形の面積を求めたときと同じように、小さな部分を加算して求めます。

例題。
(1) 半径 $r$ の球体の体積を求めよ。

薄い球殻の体積は $4\pi x^2\, dx$ である（図59 参照）；球全体を構成する同心的な殻を加算すると、
\[
\text{球体の体積}
  = \int^{x=r}_{x=0} 4\pi x^2\, dx
  = 4\pi \left[\frac{x^3}{3} \right]^r_0
  = \tfrac{4}{3} \pi r^3.
\]

球体の断面を次のようにも求めることができます：厚さ $dx$ の球体の断面は $\pi y^2\, dx$ の体積である（図62 参照）。また、$x$ と $y$ は以下の関係で結ばれています：
\[
y^2 = r^2 - x^2.
\]
したがって
\begin{align*}
\lintertext{\indent これより}
\text{球体の体積}
  &= 2 \int^{x=r}_{x=0} \pi(r^2-x^2)\, dx \\
  &= 2 \pi \left[ \int^{x=r}_{x=0} r^2\, dx - \int^{x=r}_{x=0} x^2\, dx \right] \\
  &= 2 \pi \left[r^2x - \frac{x^3}{3} \right]^r_0 = \frac{4\pi}{3} r^3.
\end{align*}

(2) $y^2=6x$ の曲線が $x$ 軸周りに回転した立体の体積を求めよ。$x=0$ から $x=4$ までの範囲で考えます。

立体の小さな断面の体積は $\pi y^2\, dx$ である。
\begin{align*}
\lintertext{\indent これより}
\text{体積}
  &= \int^{x=4}_{x=0} \pi y^2\, dx = 6\pi \int^{x=4}_{x=0} x\, dx  \\
  &= 6\pi \left[ \frac{x^2}{2} \right]^4_0 = 48\pi = 150.8.
\end{align*}

二次平均について。
物理学の一部、特に交流電流の研究においては、変動する量の「二次平均」を計算することが必要です。これは、「二次平均」として、考慮される限界間の値の平方の平均の平方根を指します。任意の量の二次平均には他の名称もあります：その「仮想的な」値や「R.M.S.」（意味は「ルート平均平方」）の値です。フランス語では「効果的な値」と呼ばれます。$y$ が考慮される関数であり、二次平均を $x=0$ から $x=l$ の範囲で求めることにすると、二次平均は以下の式で表されます：
\[
\sqrt[2] {\frac{1}{l} \int^l_0 y^2\, dx}.
\]

例題。
(1) 関数 $y=ax$ の平方平均を求めよ（図63）。

このとき、積分は $\displaystyle \int_0^l a^2 x^2 \, dx = \frac{1}{3} a^2 l^3$ である。

図2.5in 235a 63

$l$ で割って平方根を取ると、
\[
\text{平方平均} = \frac{1}{\sqrt{3}} al.
\]

この関数の算術平均は $\frac{1}{2}al$；そして平方平均と算術平均の比（これを「形因子」と呼ぶ）は $\dfrac{2}{\sqrt{3}}=1.155$。

(2) 関数 $y=x^a$ の平方平均を求めよ。

このとき、積分は $\displaystyle \int_0^l x^{2a} \, dx = \frac{l^{2a+1}}{2a+1}$ である。
\[
\text{平方平均} = \sqrt{\dfrac{l^{2a}}{2a+1}}.
\]

(3) 関数 $y=a^{\frac{x}{2}}$ の平方平均を求めよ。

このとき、積分は $\displaystyle \int_0^l (a^{\frac{x}{2}})^2 \, dx = \int_0^l a^x \, dx$ である，
\[
\left[ \frac{a^x}{\log_\epsilon a} \right]_0^l,
\]
これは $\dfrac{a^l-1}{\log_\epsilon a}$ である。

したがって、平方平均は $\sqrt{\dfrac{a^l - 1}{l \log_\epsilon a}}$ である。

練習問題 XVIII（解答は Pageref{AnsEx:XVIII} を参照）。
以下の問題を解け：

(1) $y=x^2+x-5$ の曲線における $x=0$ から $x=6$ までの面積と、これらの限界間の平均値を求めよ。

(2) $y=2a\sqrt{x}$ の放物線における $x=0$ から $x=a$ までの面積を求めよ。また、これは極限値とその横座標の矩形の二分の一であることを示せ。

(3) 正弦曲線の正の部分の面積と平均値を求めよ。

(4) $y=\sin^2 x$ の正の部分の面積と平均値を求めよ。

(5) $y=x^2 \pm x^{\frac{5}{2}}$ の二つの枝から $x=0$ まで $x=1$ までの領域の面積を求めよ。また、下側の枝の正の部分の面積も求めよ（図30 参照）。

(6) 底径 $r$、高さ $h$ の円錐体の体積を求めよ。

(7) $y=x^3 - \log_\epsilon x$ の曲線における $x=0$ から $x=1$ までの面積を求めよ。

(8) $y=\sqrt{1+x^2}$ の曲線が $x$ 軸周りに回転したときの体積を求めよ。$x=0$ から $x=4$ までである。

(9) 正弦曲線が $x$ 軸周りに回転したときの体積を求めよ。また、その表面積も求めよ。

\begin{enumerate}
\item{(10)} $xy=a$ の曲線で、$x=1$ から $x = a$ までの部分の面積を求めよ。これらの限界における平均値を求めよ。

\item{(11)} 関数 $y=\sin x$ の $0$ ラジアンから $\pi$ ラジアンまでの二次平均を $\dfrac{\sqrt{2}}{2}$ であることを示せ。また、同じ限界における算術平均を求めよ。さらに、これらの平均値の形式因子が $=1.11$ であることを示せ。

\item{(12)} 関数 $x^2+3x+2$ の $x=0$ から $x=3$ までの算術平均と二次平均を求めよ。

\item{(13)} 関数 $y=A_1 \sin x + A_1 \sin 3x$ の二次平均と算術平均を求めよ。

\item{(14)} 曲線 $y=3.42\epsilon^{0.21x}$ と $x$ 軸との間の領域を求めよ。この点から $x=8$ までの平均値の高さも求めよ。

\item{(15)} 円の半径がその面積が極座標図の面積の二倍であるとき、その円の半径は極座標図におけるすべての $r$ の二次平均に等しいことを示せ。

\item{(16)} 曲線 $y=±\dfrac{x}{6}\sqrt{x(10-x)}$ が $x$ 軸周りに回転したとき生成される体積を求めよ。
\end{enumerate}



\chapter{積分の巧妙法・落とし穴・そして勝利}
\noindent\textit{\large Dodges, Pitfalls, and Triumphs}\\[
2mm]
\noindent\textbf{\small 【章の論理とテーマ】置換積分・部分積分・三角置換・微積分学の高度な実践技巧}\\[4mm]
章{XX}{ Dodge, Pitfalls, and Triumphs}
節{Dodge}
積分を行う際の大部分の労力は、何かしらの形に整えることであり、そのような形に整形することで積分が可能になります。厳格な積分計算に関する教科書には、このような作業のために使用される計画や方法、トリックなどが満載です。以下はその一部です。
節{部分積分}
この名前は、次の公式を持つトリックに与えられます。
\[
\int u \, dx = ux - \int x \, du + C.
\]
これは直接的に扱えない某些のケースで有用であり、もし任意の場合において $\displaystyle \int x \, du$ が見つかった場合、$\displaystyle \int u \, dx$ も見つかることが示されます。この公式は以下の通り導き出せます。式（\pageref{differprod}）から、
\[
d(ux) = u \, dx + x \, du,
\]
を書き換えて
\[
u \, (dx) = d(ux) - x \, du,
\]
と表すことができます。直接積分することで上の式が得られます。

例題．
（1）$\displaystyle \int w \sin w \, dw$ を求めよ。
$u = w$ として、$\sin w \, dw$ を $dx$ とする。すると $du = dw$ であり、$\displaystyle \int \sin w \, dw = -\cos w = x$ となる。

これらの値を公式に代入すると
\begin{align*}
\int w \sin w \, dw &= w(-\cos w) - \int -\cos w \, dw \\
                    &=-w \cos w + \sin w + C.
\end{align*}

（2）$\displaystyle \int x e^x \, dx$ を求めよ。
%

\begin{align*}
\lintertext{\indent $u = x$, そして} \epsilon^x \, dx &= dv; \\
\lintertext{すると} du &= dx, v &= \epsilon^x,
\end{align*}
\begin{gather*}
\lintertext{そして}
\int x e^x \, dx
   = x e^x - \int e^x \, dx \quad \text{(公式により)} \\
   = x e^x - e^x = e^x (x-1) + C.
\end{gather*}%

(3) $\displaystyle \int \cos^2 \theta\, d\theta$を試してみてください。
\begin{align*}
u &= \cos \theta, & \cos \theta\, d\theta &= dv.
\end{align*}
\indent したがって
\begin{align*}
du&= -\sin \theta\, d\theta, & v &=\sin \theta,
\end{align*}
\begin{align*}
\int \cos^2 \theta\, d\theta
  &= \cos \theta \sin \theta+ \int \sin^2 \theta\, d\theta       \\
  &= \frac{2 \cos\theta \sin\theta}{2} +\int(1-\cos^2 \theta)\, d\theta  \\
  &= \frac{\sin 2\theta}{2} + \int d\theta - \int \cos^2 \theta\, d\theta.
\end{align*}
\begin{align*}
\lintertext{したがって}
2 \int \cos^2 \theta\, d\theta
  &= \frac{\sin 2\theta}{2} + \theta \\
\lintertext{および}
\int \cos^2 \theta\, d\theta
  &= \frac{\sin 2\theta}{4} + \frac{\theta}{2} + C.
\end{align*}

(4) $\displaystyle \int x^2 \sin x\, dx$を求めよ。
\begin{align*}
x^2  &= u, & \sin x\, dx &= dv; \\
\lintertext{すると}
du &= 2x\, dx, & v &= -\cos x,
\end{align*}
\[
\int x^2 \sin x\, dx = -x^2 \cos x + 2 \int x \cos x\, dx.
\]

次に$\displaystyle \int x \cos x\, dx$を求めよ（例1のように部分積分を行う）：
\[
\int x \cos x\, dx = x \sin x + \cos x+C.
\]

したがって
\begin{align*}
\int x^2 \sin x\, dx
  &= -x^2 \cos x + 2x \sin x + 2 \cos x + C' \\
  &= 2 \left[ x \sin x + \cos x \left(1 - \frac{x^2}{2}\right) \right] +C'.
\end{align*}

(5) $\displaystyle \int \sqrt{1-x^2}\, dx$を求めよ。
\begin{align*}
\lintertext{書き換えて}
u &= \sqrt{1-x^2},\quad dx=dv;  \\
\lintertext{すると}
du &= -\frac{x\, dx}{\sqrt{1-x^2}}\quad \text{(参照: 第9章，\Pageref{chap:IX})}
\end{align*}
$x=v$なので
\[
\int \sqrt{1-x^2}\, dx=x \sqrt{1-x^2} + \int \frac{x^2\, dx}{\sqrt{1-x^2}}.
\]

ここで少し手を加えることができる。$\displaystyle \sqrt{1-x^2}$の分母に$x^2$を足し引きできるから
\[
\int \sqrt{1-x^2}\, dx
  = \int \frac{(1-x^2)\, dx}{\sqrt{1-x^2}}
  = \int \frac{dx}{\sqrt{1-x^2}} - \int \frac{x^2\, dx}{\sqrt{1-x^2}}.
\]

これらの2つの式を加えると$\displaystyle \int \dfrac{x^2\, dx}{\sqrt{1-x^2}}$が消え、
\[
2 \int \sqrt{1-x^2}\, dx = x\sqrt{1-x^2} + \int \frac{dx}{\sqrt{1-x^2}}
\]
を得る。

$\displaystyle \dfrac {dx}{\sqrt{1-x^2}}$を思い出しているか。これは$y=\arcsin x$を微分することで得られる（参照: \Pageref{intex3}）。したがってその積分は$\arcsin x$であり、
\[
\int \sqrt{1-x^2}\, dx = \frac{x \sqrt{1-x^2}}{2} + \tfrac{1}{2} \arcsin x +C.
\]

自分でいくつかの練習問題を試してみてください。章末にいくつかあります。

\par\medskip\noindent\textbf{置換法。}これは第IX章、\pageref{chap:IX}で説明した手口と同じです。積分におけるその応用をいくつかの例によって示しましょう。

(1) $\displaystyle \int \sqrt{3+x}\, dx$。
\begin{align*}
\lintertext{\indent 令}
3+x &= u,\quad dx = du; \\
\lintertext{置換}
\int u^{\frac{1}{2}}\, du
  &= \tfrac{2}{3} u^{\frac{3}{2}} = \tfrac{2}{3}(3+x)^{\frac{3}{2}}.
\end{align*}

(2) $\displaystyle \int \dfrac{dx}{\epsilon^x+\epsilon^{-x}}$。
\begin{gather*}
\lintertext{\indent 令}
\epsilon^x = u,\quad \frac{du}{dx} = \epsilon^x,\quad\text{および}\quad
dx = \frac{du}{\epsilon^x}; \\
\lintertext{よって}
\int \frac{dx}{\epsilon^x+\epsilon^{-x}}
  = \int \frac{du}{\epsilon^x(\epsilon^x+\epsilon^{-x})}
  = \int \frac{du}{u\left(u + \dfrac{1}{u}\right)}
  = \int \frac{du}{u^2+1}.
\end{gather*}

$\dfrac{du}{1+u^2}$は$x$に関する$\arctan x$の微分結果です。

したがって、積分は$\arctan \epsilon^x$となります。

(3) $\displaystyle \int \dfrac{dx}{x^2+2x+3} = \displaystyle \int \dfrac{dx}{x^2+2x+1+2} = \displaystyle \int \dfrac{dx}{(x+1)^2+(\sqrt 2)^2}$。

\begin{gather*}
\lintertext{\indent 令}
x+1=u,\quad dx=du;
\end{gather*}
すると積分は$\displaystyle \int \dfrac{du}{u^2+(\sqrt2)^2}$に変形します；また、$\dfrac{du}{u^2+a^2}$は$u=\dfrac{1}{a} \arctan \dfrac{u}{a}$の微分結果です。

したがって、与えられた積分の値は最終的に$\dfrac{1}{\sqrt2} \arctan \dfrac{x+1}{\sqrt 2}$となります。

\emph{約分公式}は主に二項式や三角関数を含む積分に適用される特別な形式で、その積分が既知の形になるように簡略化するためのものです。

\textit{有理化}と\textit{分母因数分解}は特定の場合に適用可能な手口ですが、これらには短いまたは一般的な説明は適さない。これらの前処理過程に対する多くの練習が必要です。

次の例では、第XIII章、\pageref{partfracs2}で学んだ部分分数分解の手法が積分における応用方法を示しています。

再度、$\displaystyle \int \dfrac{dx}{x^2+2x+3}$ を考慮する。この積分を部分分数に分解すると、以下のようになります（参照：\Pageref{partfracs3}）：
\[
\dfrac{1}{2\sqrt{-2}} \left[\int \dfrac{dx}{x+1-\sqrt{-2}} - \int \dfrac{dx}{x+1+\sqrt{-2}} \right]
=
\dfrac{1}{2\sqrt{-2}} \log_\epsilon \dfrac{x+1-\sqrt{-2}}{x+1+\sqrt{-2}}.
\]
同じ積分は、しばしば異なる方法で表すことができ、これらは互いに等しい（）。

\par\medskip\noindent\textbf{注意点。} 初学者は、経験豊富な手が避けるべきポイントを無視する可能性があります。例えば、ゼロや無限大と等価な因子の使用、不定形量（例：$\dfrac{0}{0}$）の出現などです。万能のルールはありません；それには練習と知恵が必要です。第XVIII章（\Pageref{chap:XVIII}）で$x^{-1}\, dx$ の積分問題に遭遇した際、このような注意点を避けるために対処する必要がありました。

 triumphs とは、微分積分学が適用され、それまで解けなかった問題を解決する成功を指します。物理的な関係の考慮においては、部分間の相互作用や制御する力に関する法則を表現することがしばしば可能です。そのような表現は自然に微分方程式の形を取ります。すなわち、微分係数が含まれる（または他の代数的量と共に含む）方程式です。そして、このような微分方程式を見つけた後では、積分されるまでさらに進むことはできません。通常、適切な微分方程式を述べることは比較的容易ですが、それを解くのは難しい：実際の問題は、積分を行うときから始まります。ただし、方程式が既知の形式を持ち、その積分が既知である場合、そのような勝利は容易です。微分方程式を積分した結果得られる方程式は、「解」と呼ばれ、



この意味では、実際にそれを解決する結果自体が「解」と呼ばれます。しかし多くの数学者は、フォーシュ教授によれば、「依存変数の値を独立変数の関数として表現できる（既知の関数や積分を通じて）とき、微分方程式は考慮されるべき解決済みである」と言います。そして、多くの場合、解が微分方程式と直接関連しているとは思えないほど異なる形になります。解はしばしば、元の表現から蝶のように異なるように見えます。誰もが、

\[
\dfrac{dy}{dx} = \dfrac{1}{a^2-x^2}
\]

というような無害な式が、

\[
y = \dfrac{1}{2a} \log_\epsilon \dfrac{a+x}{a-x} + C
\]

のように開花するとは予想しなかったでしょう。しかし後者は前者の解です。

最後の例として、一緒に計算してみましょう。

By partial fractions,
\begin{align*}
\frac{1}{a^2-x^2} &= \frac{1}{2a(a+x)} + \frac{1}{2a(a-x)},  \\
dy &= \frac {dx}{2a(a+x)}+ \frac{dx}{2a(a-x)},  \\
y  &= \frac{1}{2a}
       \left( \int \frac{dx}{a+x}
            + \int \frac{dx}{a-x} \right)  \\
   &= \frac{1}{2a} \left(\log_\epsilon (a+x) - \log_\epsilon (a-x) \right)  \\
   &= \frac{1}{2a} \log_\epsilon \frac{a+x}{a-x} + C.
\end{align*}

Not a very difficult metamorphosis!

There are whole treatises, such as Boole's *Differential Equations*, devoted to the subject of thus finding
the ``solutions'' for different original forms.

\par\textbf{【演習 XIX】}\par (See \Pageref{AnsEx:XIX} for Answers.)

\begin{enumerate}
\item[(1)] Find $\displaystyle \int \sqrt {a^2 - x^2}\, dx$.
\item[(2)] Find $\displaystyle \int x \log_\epsilon x\, dx$.

\item[(3)] Find $\displaystyle \int x^a \log_\epsilon x\, dx$.
\item[(4)] Find $\displaystyle \int \epsilon^x \cos \epsilon^x\, dx$.

\item[(5)] Find $\displaystyle \int \dfrac{1}{x} \cos (\log_\epsilon x)\, dx$.
\item[(6)] Find $\displaystyle \int x^2 \epsilon^x\, dx$.

\item[(7)] Find $\displaystyle \int \dfrac{(\log_\epsilon x)^a}{x}\, dx$.
\item[(8)] Find $\displaystyle \int \dfrac{dx}{x \log_\epsilon x}$.

\item[(9)] Find $\displaystyle \int \dfrac{5x+1}{x^2 +x-2}\, dx$.
\item[(10)] Find $\displaystyle \int \dfrac{(x^2 -3)\, dx}{x^3 - 7x+6}$.

\item[(11)] Find $\displaystyle \int \dfrac{b\, dx}{x^2 -a^2}$.
\item[(12)] Find $\displaystyle \int \dfrac{4x\, dx}{x^4 -1}$.

\item[(13)] Find $\displaystyle \int \dfrac{dx}{1-x^4}$.
\item[(14)] Find $\displaystyle \int \dfrac{dx}{x \sqrt {a-bx^2}}$.

\end{enumerate}



\chapter{微分方程式の解法}
\noindent\textit{\large Finding some Solutions}\\[
2mm]
\noindent\textbf{\small 【章の論理とテーマ】変数分離形微分方程式・物理現象や工学問題の数理モデルの完全解}\\[4mm]
この章では、前章で示した過程を使用して、いくつか重要な微分方程式の解を求める作業に取り組みます。
初学者は、現在のところこれらの過程の大半がいかに簡単であるかを知っているため、ここから始めて積分が\emph{芸術}であることを実感するでしょう。すべての芸術と同様に、この芸術においても、熟練度を得るには怠らない定期的な練習が必要です。その熟練度を手に入れるためには、微分積分学の正規の教科書で見つかる例題を解き、さらに多くの問題を解く必要があります。
例1．以下の微分方程式の解を求めよ：
\[
ay + b \frac{dy}{dx} = 0.
\]

移項すると
\[
b \frac{dy}{dx} = -ay.
\]

この関係式を見れば、$\dfrac{dy}{dx}$が$y$に比例していることがわかります。$y$を$x$の関数として表す曲線を考えると、その任意の点での傾きはその点における縦座標に比例し、$y$が正であれば負の傾きとなります。したがって、明らかにこの曲線は衰減曲線（\Pageref{section:5}）であり、解には$\epsilon^{-x}$が因子として含まれます。しかし、この洞察力を前提にせずに、作業を始めましょう。

$y$と$dy$が方程式の両側に現れ、一方は$dx$の反対側にあるため、これらを片方に集め、もう片方に$dx$を移すことはできません。これを達成するには、通常 inseparable な $dy$ と $dx$ を分離する必要があります。
\[
\frac{dy}{y} = - \frac{a}{b}\, dx.
\]

Having 作業を終えたら、両辺が積分可能な形になっていることが確認できます。なぜなら、$\dfrac{dy}{y}$ や $\dfrac{1}{y}\, dy$ を対数微分の際に遭遇した差分と認識できるからです。そこで、次の積分手順を即座に記述できます：
\[
\int \frac{dy}{y} = \int -\frac{a}{b}\, dx;
\]
両方の積分を行いますと、
\[
\log_\epsilon y = -\frac{a}{b} x + \log_\epsilon C,
\]
となります。ここで $\log_\epsilon C$ は未定義の定数で、積分定数の形としてここでは $\log_\epsilon C$ を採用しています。他の項も対数と扱われているため、後々の計算を簡略化します。次に、デログアライズ（逆対数）を行います：
\[
y = C \epsilon^{-\frac{a}{b} x},
\]
これが求められた解です。この解は元の微分方程式とは全く異なる形に見えますが、専門的な数学家には両者は $y$ が $x$ にどのように依存するかという情報を同じように伝えています。

次に $C$ の意味についてですが、これは $y$ の初期値によって定義されます。つまり、$x = 0$ を代入して $y$ の値を確認すると、$y = C \epsilon^{-0}$ となります。そして $\epsilon^{-0} = 1$ より、$C$ は $y$ の初期値そのものであり、これを $y_0$ と呼びます。したがって解を以下のように書き換えることができます：
\[
y = y_0 \epsilon^{-\frac{a}{b} x}.
\]

例題2：

例として、次の方程式を考えます：
\[
ay + b \frac{dy}{dx} = g,
\]
ここで $g$ は定数です。方程式を見ると、(1) 何かしら $\epsilon^x$ が解に含まれる可能性があるし、(2) 曲線 $y$ の一部で最大値や最小値が現れ、$\frac{dy}{dx} = 0$ となる場合、$y$ の値は $= \frac{g}{a}$ になることが推測されます。しかし、以前と同じように、微分を分離し、積分可能な形に変換してみましょう。
\begin{align*}
b\frac{dy}{dx} &= g - ay; \\
\frac{dy}{dx} &= \frac{a}{b}\left(\frac{g}{a} - y\right); \\
\frac{dy}{y - \frac{g}{a}} &= -\frac{a}{b} dx.
\end{align*}

ここで、$y$ と $dy$ のみが一方に、$dx$ のみが他方に配置されました。しかし、左辺の結果は積分可能でしょうか？

これは \Pageref{differlog} の結果と同じ形式です。因此，積分の指示を書き出すと次のようになります：
\[
\int{\frac{dy}{y - \frac{g}{a}}} = - \int{\frac{a}{b} dx};
\]
そして、積分を行い、適切な定数を加えると，
\begin{align*}
\log_\epsilon\left(y - \frac{g}{a}\right) &= -\frac{a}{b}x + \log_\epsilon C; \\
\lintertext{従って } y - \frac{g}{a} &= C\epsilon^{-\frac{a}{b}x}; \\
\lintertext{そして最終的に } y &= \frac{g}{a} + C\epsilon^{-\frac{a}{b}x},
\end{align*}
これが \emph{解} です。

条件が $y = 0$ のとき $x = 0$ であると定めると、$C$ を求めることができます。そのとき指数関数は $= 1$ になりますので，
\begin{align*}
0 &= \frac{g}{a} + C, \\
\lintertext{従って } C &= -\frac{g}{a}.
\end{align*}

この値を代入すると、解は
\[
y = \frac{g}{a} (1 - \epsilon^{-\frac{a}{b} x})
\]
となります。

さらに $x$ が無限大に増加する場合、$y$ は最大値に収束します。$x = \infty$ のとき指数関数は $= 0$ となり、$y_{\text{max.}} = \frac{g}{a}$ となります。これを代入すると最終的に
\[
y = y_{\text{max.}} (1 - \epsilon^{-\frac{a}{b} x})
\]
となります。

この結果は物理学でも重要な意味を持ちます。

例題3．

\begin{gather*}
\lintertext{\indent 令} ay + b \frac{dy}{dt} = g \sin 2\pi nt.
\end{gather*}

これは前の問題よりも扱いが難しいです。まず $b$ で割ります。
\[
\frac{dy}{dt} + \frac{a}{b} y = \frac{g}{b} \sin 2\pi nt.
\]

今、左辺は積分可能ではありません。
しかし、全ての項に $\epsilon^{\frac{a}{b}t}$ を掛けるという技巧によって、そのような形にすることができます。そしてこれは、技術と練習が示す計画です：
\[
\frac{dy}{dt} \epsilon^{\frac{a}{b} t} + \frac{a}{b} y \epsilon^{\frac{a}{b} t} = \frac{g}{b} \epsilon^{\frac{a}{b} t} \sin 2 \pi nt,
\]
これは
\[
\frac{dy}{dt} \epsilon^{\frac{a}{b} t} + y \frac{d(\epsilon^{\frac{a}{b} t})}{dt} = \frac{g}{b} \epsilon^{\frac{a}{b} t} \sin 2 \pi nt;
\]
と等価であり、これは完全な微分であるため、次のように積分できます：---もし $u = y\epsilon^{\frac{a}{b} t}$ なら、$\dfrac{du}{dt} = \dfrac{dy}{dt} \epsilon^{\frac{a}{b} t} + y \dfrac{d(\epsilon^{\frac{a}{b} t})}{dt}$ より
\begin{align*}
  y \epsilon^{\frac{a}{b} t}
  &= \frac{g}{b} \int \epsilon^{\frac{a}{b} t} \sin 2 \pi nt \, dt + C, \\
\lintertext{または}
y &= \frac{g}{b} \epsilon^{-\frac{a}{b} t}
     \int \epsilon^{ \frac{a}{b} t} \sin 2\pi nt \, dt
       + C\epsilon^{-\frac{a}{b} t}.
\tag*{[\textsc{a}]}
\end{align*}

最後の項は明らかに $t$ が増加するにつれて消滅する項であり、省略することができます。現在の問題は、現れる因子としての積分を見つけ出すことです。これを扱うために、部分積分（参照：\Pageref{intparts}）という装置を使用します。一般的な公式は $\displaystyle  \int u dv = uv - \int v du$ です。この目的のために
\begin{align*}
&\left\{
\begin{aligned}
 u &= \epsilon^{\frac{a}{b} t}; \\
dv &= \sin 2\pi nt \, dt.
\end{aligned}
\right.  \\
\intertext{とすると}
\&\left\{
\begin{aligned}
du &= \epsilon^{\frac{a}{b} t} \times \frac{a}{b} \, dt; \\
v &= - \frac{1}{2\pi n} \cos 2\pi nt.
\end{aligned}
\right.
\end{align*}

これらの値を代入すると、求める積分は次のようになります：
\begin{align*}
\int \epsilon^{\frac{a}{b} t} &{} \sin 2 \pi n t \, dt \\
&= -\frac{1}{2 \pi n} \epsilon^{\frac{a}{b} t} \cos 2 \pi nt
   -\int -\frac{1}{2\pi n} \cos 2 \pi nt \epsilon^{\frac{a}{b} t} \frac{a}{b} \, dt \\
&= -\frac{1}{2 \pi n} \epsilon^{\frac{a}{b} t} \cos 2 \pi nt
   +\frac{a}{2 \pi nb} \int \epsilon^{\frac{a}{b} t} \cos 2 \pi nt \, dt.
\tag*{[\textsc{b}]}
\end{align*}

最後の積分はまだ簡約化できません。困難を避けるため、左辺の積分を逆方向に書き換えて部分積分を繰り返します：
\[
\begin{aligned}
u &= \sin 2 \pi n t; \\
dv &= \epsilon^{\frac{a}{b} t} \, dt;
\end{aligned}
\]
したがって
\[
\begin{aligned}
du &= 2 \pi n \cos 2 \pi n t \, dt; \\
v &= \frac{b}{a} \epsilon^{\frac{a}{b} t}.
\end{aligned}
\]

これらを代入すると、次のようになります：
\begin{align*}
\int \epsilon^{\frac{a}{b} t} &{} \cdot \sin 2 \pi n t \, dt \\
&= \frac{b}{a} \epsilon^{\frac{a}{b} t} \sin 2 \pi n t - \frac{2 \pi n b}{a} \int \epsilon^{\frac{a}{b} t} \cos 2 \pi n t \, dt. \tag*{[\textsc{c}]}
\end{align*}

式 [\textsc{c}] の最終的な未解決の積分は [\textsc{b}] と同一であることに注意すると、 [\textsc{b}] を $\frac{2 \pi n b}{a}$ で掛けて [\textsc{c}] を $\frac{a}{2 \pi n b}$ で掛けることでこれを消去できます。

結果は次のようになります：
\begin{align*}
\int \epsilon^{\frac{a}{b} t} &{} \cdot \sin 2 \pi n t \, dt \\
&= \epsilon^{\frac{a}{b} t} \left\{\frac{ab \sin 2 \pi n t - 2 \pi n b^2 \cos 2 \pi n t}{a^2 + 4 \pi^2 n^2 b^2}\right\}. \tag*{[\textsc{d}]}
\end{align*}

 [\textsc{a}] にこの値を代入すると、次のようになります：
\[
y = g \left\{\frac{a \sin 2 \pi n t - 2 \pi n b \cos 2 \pi n t}{a^2 + 4 \pi^2 n^2 b^2}\right\}.
\]

さらに簡略化するために、$\tan \phi = \frac{2 \pi n b}{a}$ のような角度 $\phi$ を想像します。
\[
\begin{aligned}
\sin \phi &= \frac{2 \pi nb}{\sqrt{a^2 + 4 \pi^2 n^2 b^2}}, \\
\cos \phi &= \frac{a}{\sqrt{a^2 + 4 \pi^2 n^2 b^2}}.
\end{aligned}
\]
これらを代入すると、次のようになります：
\[
y = g \frac{\cos \phi \sin 2 \pi n t - \sin \phi \cos 2 \pi n t}{\sqrt{a^2 + 4 \pi^2 n^2 b^2}},
\]
これは
\[
y = g \frac{\sin(2 \pi n t - \phi)}{\sqrt{a^2 + 4 \pi^2 n^2 b^2}}
\]
と書き換えることができます。これが求めた解です。

この式は実際には、電流の交差する電圧を表しており、$g$ は電動勢の振幅、$n$ は周波数、$a$ は抵抗、$b$ は回路の自己誘導係数、$\phi$ は遅延角を表します。

例4．

\begin{gather*}
\lintertext{\indent 假如}
M\, dx + N\, dy = 0.
\end{gather*}

この式を直接積分することは、$M$が$x$のみの関数であり、$N$が$y$のみの関数である場合に可能ですが、もし両方とも$M$と$N$が$x$と$y$の両方に依存する関数であれば、どのように積分すればよいでしょうか？この式は正確な微分形でしょうか？すなわち、$M$と$N$が共通の関数$U$から部分微分によって形成されたかどうかです。もしそうであれば、

\[
  \begin{aligned}
    \frac{\partial U}{\partial x} &= M, \\
    \frac{\partial U}{\partial y} &= N.
  \end{aligned}
\]

そして、このような共通の関数が存在すれば、

\[
\frac{\partial U}{\partial x}\, dx + \frac{\partial U}{\partial y}\, dy
\]

は正確な微分形となります（参照：第\pageref{partialdiff}節）。

今、問題の本質は次の通りです。この式が正確な微分形である場合、以下が成り立つ必要があります：

\begin{align*}
        \frac{\partial M}{\partial y} &= \frac{\partial N}{\partial x}; \\
\lintertext{すなわち}
\frac{\partial (dU)}{\partial x\, dy} &= \frac{\partial (dU)}{\partial y\, dx},\\
\lintertext{これは必ず成り立ちます。}
\end{align*}

例として、方程式

\[
(1 + 3xy)\, dx + x^2\, dy = 0.
\]

を考えましょう。

この式が正確な微分形であるかどうかをテストします。
\[
\left\{
  \begin{aligned}
    \frac{\partial (1 + 3xy)}{\partial y} &= 3x, \\
    \dfrac{\partial (x^2)}{\partial x} &= 2x,
  \end{aligned}
\right.
\]

これらは一致していません。したがって、この式は正確な微分形ではなく、関数$1+3xy$と$x^2$は共通の原始関数から形成されていません。

このような場合でも、積分因子を見つけることが可能です。すなわち、両方をこの因子で乗じることによって、式が正確な微分形になるような因子です。このような積分因子を見つけるための一贯したルールはありませんが、経験は通常適切なものを示唆します。現在の例では$2x$がその役割を果たします。$2x$で乗じると、

\[
(2x + 6x^2y)\, dx + 2x^3\, dy = 0.
\]

Now apply the test to this.
\[
\left\{
  \begin{aligned}
   \frac{d(2x + 6x^2y)}{dy} &= 6x^2, \\
   \dfrac{d(2x^3)}{dx} &= 6x^2,
  \end{aligned}
\right.
\]
which agrees. Hence this is an exact differential, and
may be integrated. Now, if \( w = 2x^3y \),
\begin{gather*}
dw = 6x^2y \, dx + 2x^3 \, dy. \\
\lintertext{\indent Hence} \int 6x^2y \, dx + \int 2x^3 \, dy = w = 2x^3y; \\
\lintertext{so that we get} U = x^2 + 2x^3y + C.
\end{gather*}

\par\textbf{【例題 5.】}\par Let \(\frac{d^2 y}{dt^2} + n^2 y = 0\).

In this case we have a differential equation of the
second degree, in which \(y\) appears in the form of
a second differential coefficient, as well as in person.

Transposing, we have \(\frac{d^2 y}{dt^2} = - n^2 y\).

From this it is clear that we are dealing with a
function such that its second differential coefficient is
proportional to itself, but with reversed sign. In
Chapter\textasciitilde{}XV., we found that there was such a function---namely,
the \emph{sine} (or the \emph{cosine} also) which
possessed this property. So, without further ado,
we may infer that the solution will be of the form
\(y = A \sin (pt + q)\). However, let us go to work.

Multiply both sides of the original equation by \(2\frac{dy}{dt}\)
and integrate, giving us \(2\frac{d^2 y}{dt^2} \, \frac{dy}{dt} + 2x^2 y \frac{dy}{dt} = 0\), and, as
\[
2 \frac{d^2y}{dt^2} \, \frac{dy}{dt}
  = \frac{d \left(\frac{dy}{dt}\right)^2}{dt},\quad
\left(\frac{dy}{dt}\right)^2 + n^2 (y^2 - C^2) = 0,
\]
\(C\) being a constant. Then, taking the square roots,
\[
\frac{dy}{dt} = -n \sqrt{ y^2 - C^2}\quad \text{and}\quad
\frac{dy}{\sqrt{C^2 - y^2}} = n \cdot dt.
\]

But it can be shown that (see \Pageref{intex3})
\[
\frac{1}{\sqrt{C^2 - y^2}} = \frac{d (\arcsin \frac{y}{C})}{dy};
\]
whence, passing from angles to sines,
\[
\arcsin \frac{y}{C} = nt + C_1\quad \text{and}\quad y = C \sin (nt + C_1),
\]

where \(C_1\) is a constant angle that comes in by integration.

Or, preferably, this may be written
\[
y = A \sin nt + B \cos nt, \text{ which is the solution.}
\]

\par\textbf{【例題 6.】}\par \hfil \(\frac{d^2 y}{dt^2} - n^2 y = 0\).

ここでは明らかに、二階微分係数が自分自身と比例する関数 \( y \) を扱う必要があります。このような性質を持つ唯一の関数は指数関数（\pageref{unchanged} を参照）であり、従って方程式の解はその形になるでしょう。

以前と同じように、両辺に \( 2 \frac{dy}{dx} \) を掛けて積分すると、
\[
 2 \frac{d^2 y}{dx^2} \frac{dy}{dx} - 2x^2 y \frac{dy}{dx} = 0, 
\]
\begin{gather*}
\lintertext{そして}
2 \frac{d^2 y}{dx^2} \frac{dy}{dx} 
  = \frac{d \left( \frac{dy}{dx} \right)^2}{dx}, \quad
\left( \frac{dy}{dx} \right)^2 - n^2 (y^2 + c^2) = 0, \\
\frac{dy}{dx} - n \sqrt{y^2 + c^2} = 0,
\end{gather*}
ここで \( c \) は定数であり、
\[
 \frac{dy}{\sqrt{y^2 + c^2}} = n \, dx. 
\]

今、もし
\[
 w = \log_\epsilon ( y + \sqrt{y^2 + c^2}) = \log_\epsilon u, 
\]
\begin{gather*}
\frac{dw}{du} = \frac{1}{u}, \quad \frac{du}{dy} = 1 + \frac{y}{\sqrt{y^2 + c^2}} = \frac{y + \sqrt{ y^2 + c^2}}{\sqrt{y^2 + c^2}}, \\
\lintertext{そして}
\frac{dw}{dy} = \frac{1}{\sqrt{ y^2 + c^2}}.
\end{gather*}

従って、積分すると
\begin{gather*}
\log_\epsilon (y + \sqrt{y^2 + c^2}) = nx + \log_\epsilon C, \\
y + \sqrt{y^2 + c^2} = C \epsilon^{nx}.
\tag*{(1)}  \\
\lintertext{\indent 今}
( y + \sqrt{y^2 + c^2} ) \times  ( -y + \sqrt{y^2 + c^2} ) = c^2 ;    \\
\lintertext{従って}
- y + \sqrt{y^2 + c^2} = \frac{c^2}{C} \epsilon^{-nx}.
\tag*{(2)}
\end{gather*}

(1) から (2) を引いて 2 で割ると、
\[
 y = \frac{1}{2} C \epsilon^{nx} - \frac{1}{2} \frac{c^2}{C} \epsilon^{-nx}, 
\]
これはより便利に書き換えると
\[
 y = A \epsilon^{nx} + B \epsilon^{-nx}. 
\]
または、最初の見かけでは元の方程式とは関係がないように思われる解ですが、\( y \) は二つの項から成り立っており、一つは \( x \) の増加とともに対数的に成長し、もう一つは \( x \) の増加とともに消滅します。

例 7.
\begin{align*}
\lintertext{\indent 令}
b \frac{d^2y}{dt^2} + a \frac{dy}{dt} + gy &= 0.
\end{align*}

この式を検討すると、$b = 0$ の場合、それは例1の形になります。その解は負指数でした。一方で、$a = 0$ の場合、その形式は例6と同じになり、その解は正と負の指数の和です。したがって、現在の例の解が
\begin{align*}
y &= (\epsilon^{-mt})(A \epsilon^{nt} + B \epsilon^{-nt}), \\
m &= \frac{a}{2b}\quad \text{および}\quad
n  = \sqrt{\frac{a^2}{4b^2}} - \frac{g}{b}.
\end{align*}
であることが見当たらないわけではありません。

この解が得られる手順はここでは述べられていません；詳細は高度な文献で見つけることができます。

例8．
\[
\frac{d^2y}{dt^2} = a^2 \frac{d^2y}{dx^2}.
\]

これは（例4参照）元の式
\[
y = F(x+at) + f(x-at),
\]
から導かれたものであり、$F$ と $f$ は任意の $t$ の関数でした。

別の方法として、変数の変換により次の形に変形することができます：
\[
\frac{d^2y}{du \cdot dv} = 0,
\]
ここで $u = x + at$, $v = x - at$ であり、これは同じ一般解を導きます。もし $F$ が消滅する場合を考えると、単に
\[
y = f(x-at)
\]
となります。これは $t=0$ のとき $y$ は $x$ の特定の関数であることを述べるだけであり、$y$ と $x$ の関係の曲線が特定の形状を持っていることを示しています。その後の時間 $t$ の任意の変化は単に $x$ の原点を変更することに相当します。つまり、関数の形を保つ一方で、$x$ 方向に一様な速度 $a$ で伝播していることを示しています。したがって、特定の時間 $t_0$ における特定の点 $x_0$ の縦座標 $y$ の任意の値は、その後の時間 $t_1$ において、$x_0 + a(t_1 - t_0)$ の横座標を持つ点で同じ値を示します。この場合、簡略化された方程式は任意の形状の波が $x$ 方向に一様な速度で伝播することを表しています。

微分方程式が
\[
m \frac{d^2y}{dt^2} = k \, \frac{d^2y}{dx^2}
\]
と書かれていた場合、解は同じでしたが、伝播速度は
\[
a = \sqrt{\frac{k}{m}}
\]
となります。

你现在已被亲自引领進入了这片神奇的土地。为了让你方便查阅主要结果，作者在告辞之际请求为你提供一张护照的形式，即一份实用的標准形式集合（参见\pageref{stdforms1}～\pageref{stdforms2})。
中央列列举了一些最常见的函数。它们的导数写在左侧，积分结果写在右侧。
希望対你有所帮助！



\chapter{エピローグ：エピローグと寓話}

\noindent\textbf{\large 【概要】読者への激励・数学の門をくぐり抜けた旅人へ}\\[
4mm]
Epilogue and Apologue
この「微分積分の簡単な説明」が専門的な数学者の手に渡ったとき、彼らは（怠け者でなければ）一斉に立ち上がり、これを見下し、完璧に悪い本だと非難するでしょう。その点から見れば、絶対的に疑いの余地はありません。この著作にはいくつか深刻かつ悲劇的な誤りがあります。
まず第一に、微分積分の多くの操作がいかに驚くほど簡単であるかを示しています。
第二に、数々の業界秘密を明かします。あなたに「一つの愚か者ができることが他にもう一度できる」ということを示すことで、数学的な権威者たちが微分積分という非常に難解な主題をマスターしたと自慢している理由は大きくないことを見せます。彼らはその科目がいかに難しくて恐ろしいものであるかとあなたを信じ込ませたいし、その迷信が粗暴に消え去ることを望んでいません。
第三に、「So Easy」についての彼らの恐ろしい言葉には、著者がいくつかの方法の妥当性を厳格かつ満足するほど完全な証明を行わなかったという点があります。彼はこれらの方法を単純な形で提示し、甚至は問題を解決するためにそれらを使用したにもかかわらず、その正当性を示すことを恐れずに挑戦しました！しかし、なぜ彼がそれをすべきでないでしょうか？誰も作る方法を知らないからといって時計を使うことを禁止しますか？自作のバイオリンで演奏することを阻止しますか？言葉を使いこなせる前に文法のルールを教えることはありません。微分積分の初学者に対して一般的な厳格な証明を要求することは同等に荒唐無稽です。
この完全に悪い悪質な書籍について、専門の数学者たちはもう一つ言います：その理由がなぜ \emph{こんなに簡単} であるかは、著者が本当に難しいことをすべて省いたからです。そして、この批判の恐ろしい事実は---\emph{それが真実なのです！} つまり、本書は、微分積分学の要素を習得するのにこれまで intimid されていた無防備な人々のために書かれました。その教え方がほとんど常に不適切であるためにです。どんな科目も困難がたくさんあると提示すれば repellent になります。この本の目標は、初心者がその言語を学び、愛すべき単純さに親しみ、問題解決の強力な方法を理解するためであり、実用的な数学者が好む複雑で偏った（そして大抵無関係な）数学的トレーニングを通じて労力を費やすことを強制されることなくです。
若手の技術者の中には、\emph{何一つ愚か者ができることがあれば、他の愚か者もできる} という格言が耳に馴染んでいるものがいます。彼らは著者の秘密を漏らすな、また数学者たちに自分がどれほど愚かであるかを告げないよう懇願します。
\chapter{付録1：標準公式一覧表}
\noindent\textbf{\large 【概要】すべての微積分公式の体系的便覧}\\[4mm]
\begin{center}
\begin{tabular}{|c|c|c|}
\hline
$\dfrac{dy}{dx}$ & $\longleftarrow\quad y\quad\longrightarrow$ & $\displaystyle\int y\, dx$ \\
\hline
\multicolumn{3}{|c|}{\textbf{代数関数（Algebraic）}} \\
\hline
$1$ & $x$ & $\frac{1}{2} x^2 + C$ \\
$0$ & $a$ & $ax + C$ \\
$1$ & $x \pm a$ & $\frac{1}{2} x^2 \pm ax + C$ \\
$a$ & $ax$ & $\frac{1}{2} ax^2 + C$ \\
$2x$ & $x^2$ & $\frac{1}{3} x^3 + C$ \\
$nx^{n-1}$ & $x^n$ & $\dfrac{1}{n+1} x^{n+1} + C$ \\
$-x^{-2}$ & $x^{-1}$ & $\log_\epsilon x + C$ \\
$\dfrac{du}{dx} \pm \dfrac{dv}{dx} \pm \dfrac{dw}{dx}$ & $u \pm v \pm w$ & $\int u\, dx \pm \int v\, dx \pm \int w\, dx$ \\
$u\, \dfrac{dv}{dx} + v\, \dfrac{du}{dx}$ & $uv$ & 一般公式なし \\
$\dfrac{v\, \dfrac{du}{dx} - u\, \dfrac{dv}{dx}}{v^2}$ & $\dfrac{u}{v}$ & 一般公式なし \\
$\dfrac{du}{dx}$ & $u$ & $ux - \int x\, du + C$ \\
\hline
\multicolumn{3}{|c|}{\textbf{指数関数・対数関数（Exponential and Logarithmic）}} \\
\hline
$\epsilon^x$ & $\epsilon^x$ & $\epsilon^x + C$ \\
$x^{-1}$ & $\log_\epsilon x$ & $x(\log_\epsilon x - 1) + C$ \\
$0.4343 \times x^{-1}$ & $\log_{10} x$ & $0.4343x (\log_\epsilon x - 1) + C$ \\
$a^x \log_\epsilon a$ & $a^x$ & $\dfrac{a^x}{\log_\epsilon a} + C$ \\
\hline
\multicolumn{3}{|c|}{\textbf{三角関数（Trigonometrical）}} \\
\hline
$\cos x$ & $\sin x$ & $-\cos x + C$ \\
$-\sin x$ & $\cos x$ & $\sin x + C$ \\
$\sec^2 x$ & $\tan x$ & $-\log_\epsilon \cos x + C$ \\
\hline
\multicolumn{3}{|c|}{\textbf{逆三角関数（Circular / Inverse）}} \\
\hline
$\dfrac{1}{\sqrt{1-x^2}}$ & $\arcsin x$ & $x \cdot \arcsin x + \sqrt{1 - x^2} + C$ \\
$-\dfrac{1}{\sqrt{1-x^2}}$ & $\arccos x$ & $x \cdot \arccos x - \sqrt{1 - x^2} + C$ \\
$\dfrac{1}{1+x^2}$ & $\arctan x$ & $x \cdot \arctan x - \frac{1}{2} \log_\epsilon (1 + x^2) + C$ \\
\hline
\multicolumn{3}{|c|}{\textbf{双曲線関数（Hyperbolic）}} \\
\hline
$\cosh x$ & $\sinh x$ & $\cosh x + C$ \\
$\sinh x$ & $\cosh x$ & $\sinh x + C$ \\
$\sech^2 x$ & $\tanh x$ & $\log_\epsilon \cosh x + C$ \\
\hline
\multicolumn{3}{|c|}{\textbf{その他の有用な標準公式（Miscellaneous）}} \\
\hline
$-\dfrac{1}{(x + a)^2}$ & $\dfrac{1}{x + a}$ & $\log_\epsilon (x+a) + C$ \\
$-\dfrac{x}{(a^2 + x^2)^{\frac{3}{2}}}$ & $\dfrac{1}{\sqrt{a^2 + x^2}}$ & $\log_\epsilon (x + \sqrt{a^2 + x^2}) + C$ \\
$\mp \dfrac{b}{(a \pm bx)^2}$ & $\dfrac{1}{a \pm bx}$ & $\pm \dfrac{1}{b} \log_\epsilon (a \pm bx) + C$ \\
$-\dfrac{3a^2x}{(a^2 + x^2)^{\frac{5}{2}}}$ & $\dfrac{a^2}{(a^2 + x^2)^{\frac{3}{2}}}$ & $\dfrac{x}{\sqrt{a^2 + x^2}} + C$ \\
$a \cos ax$ & $\sin ax$ & $-\dfrac{1}{a} \cos ax + C$ \\
$-a \sin ax$ & $\cos ax$ & $\dfrac{1}{a} \sin ax + C$ \\
$a \sec^2 ax$ & $\tan ax$ & $-\dfrac{1}{a} \log_\epsilon \cos ax + C$ \\
$\sin 2x$ & $\sin^2 x$ & $\dfrac{x}{2} - \dfrac{\sin 2x}{4} + C$ \\
$-\sin 2x$ & $\cos^2 x$ & $\dfrac{x}{2} + \dfrac{\sin 2x}{4} + C$ \\
$n \sin^{n-1} x \cos x$ & $\sin^n x$ & $-\dfrac{\cos x}{n} \sin^{n-1} x + \dfrac{n-1}{n} \int \sin^{n-2} x\, dx + C$ \\
$-\dfrac{\cos x}{\sin^2 x}$ & $\dfrac{1}{\sin x}$ & $\log_\epsilon \tan \dfrac{x}{2} + C$ \\
$-\dfrac{\sin 2x}{\sin^4 x}$ & $\dfrac{1}{\sin^2 x}$ & $-\cotan x + C$ \\
$\dfrac{\sin^2 x - \cos^2 x}{\sin^2 x \cos^2 x}$ & $\dfrac{1}{\sin x \cos x}$ & $\log_\epsilon \tan x + C$ \\
$2a \sin 2ax$ & $\sin^2 ax$ & $\dfrac{x}{2} - \dfrac{\sin 2ax}{4a} + C$ \\
$-2a \sin 2ax$ & $\cos^2 ax$ & $\dfrac{x}{2} + \dfrac{\sin 2ax}{4a} + C$ \\
\hline
\end{tabular}
\end{center}
\chapter{付録2：全演習問題解答}
\noindent\textbf{\large 【概要】原著全章の全問完全解答集}\\[4mm]
\par\bigskip\noindent{\Large\bfseries [演習 I の解答]}\par\smallskip
\begin{enumerate}
\item[(1)] $\dfrac{dy}{dx} = 13x^{12}$.
\item[(2)] $\dfrac{dy}{dx} = - \dfrac{3}{2} x^{-\frac{5}{2}}$.
\item[(3)] $\dfrac{dy}{dx} = 2ax^{(2a-1)}$.
\item[(4)] $\dfrac{du}{dt} = 2.4t^{1.4}$.
\item[(5)] $\dfrac{dz}{du} = \dfrac{1}{3} u^{-\frac{2}{3}}$.
\item[(6)] $\dfrac{dy}{dx} = -\dfrac{5}{3}x^{-\frac{8}{3}}$.
\item[(7)] $\dfrac{du}{dx} = -\dfrac{8}{5}x^{-\frac{13}{5}}$.
\item[(8)] $\dfrac{dy}{dx} = 2ax^{a-1}$.
\item[(9)] $\dfrac{dy}{dx} = \dfrac{3}{q} x^{\frac{3-q}{q}}$.
\item[(10)] $\dfrac{dy}{dx} = -\dfrac{m}{n} x^{-\frac{m+n}{n}}$.
\par\smallskip
\end{enumerate}
\par\bigskip\noindent{\Large\bfseries [演習 II の解答]}\par\smallskip
\begin{enumerate}
\item[(1)] $\dfrac{dy}{dx} = 3ax^2$.
\item[(2)] $\dfrac{dy}{dx} = 13 \times  \frac{3}{2}x^{\frac{1}{2}}$.
\item[(3)] $\dfrac{dy}{dx} = 6x^{-\frac{1}{2}}$.
\item[(4)] $\dfrac{dy}{dx} = \dfrac{1}{2}c^{\frac{1}{2}} x^{-\frac{1}{2}}$.
\item[(5)] $\dfrac{du}{dz} = \dfrac{an}{c} z^{n-1}$.
\item[(6)] $\dfrac{dy}{dt} = 2.36t$.
\item[(7)] $\dfrac{dl_t}{dt} = 0.000012\times l_0$.
\item[(8)] $\dfrac{dC}{dV} = abV^{b-1}$, $0.98$, $3.00$ and $7.47$\textasciitilde{}candle power per volt respectively.
\item[(9)] $\begin{aligned}[t]
  \dfrac{dn}{dD} &= -\dfrac{1}{LD^2} \sqrt{\dfrac{gT}{\pi \sigma}}, &
  \dfrac{dn}{dL} &= -\dfrac{1}{DL^2} \sqrt{\dfrac{gT}{\pi \sigma}}, \\
%
  \dfrac{dn}{d \sigma}
  &= -\dfrac{1}{2DL} \sqrt{\dfrac{gT}{\pi \sigma^3}}, &
  \dfrac{dn}{dT} &=  \dfrac{1}{2DL} \sqrt{\dfrac{g}{\pi \sigma T}}.
\end{aligned}$
\item[(10)] $\dfrac{\text{Rate of change of~$P$ when $t$\textasciitilde{}varies}}
            {\text{Rate of change of\textasciitilde{}$P$ when $D$\textasciitilde{}varies}}
  = - \dfrac{D}{t}$.
\item[(11)] $2\pi$, $2\pi r$, $\pi l$, $\frac{2}{3}\pi rh$, $8\pi r$, $4\pi r^2$. \hfil
(12)\textasciitilde{}$\dfrac{dD}{dT} = \dfrac{0.000012l_t}{\pi}$.
\par\smallskip
\end{enumerate}
\par\bigskip\noindent{\Large\bfseries [演習 III の解答]}\par\smallskip
\begin{enumerate}
\item[(1)] (\textit{a}) $1 + x + \dfrac{x^2}{2} + \dfrac{x^3}{6} + \dfrac{x^4}{24} + \ldots$ \qquad
    (\textit{b}) $2ax + b$. \qquad (\textit{c}) $2x + 2a$.
    (\textit{d}) $3x^2 + 6ax + 3a^2$.
\item[(2)] $\dfrac{dw}{dt} = a - bt$.
\item[(3)] $\dfrac{dy}{dx} = 2x$.
\item[(4)] $14110x^4 - 65404x^3 - 2244x^2 + 8192x + 1379$.
\item[(5)] $\dfrac{dx}{dy} = 2y + 8$.
\item[(6)] $185.9022654x^2 + 154.36334$.
\item (7) \(\dfrac{-5}{(3x + 2)^2}\).
\item (8) \(\dfrac{6x^4 + 6x^3 + 9x^2}{(1 + x + 2x^2)^2}\).
\item (9) \(\dfrac{ad - bc}{(cx + d)^2}\).
\item (10) \(\dfrac{anx^{-n-1} + bnx^{n-1} + 2nx^{-1}}{(x^{-n} + b)^2}\).
\item (11) \(b + 2ct\).
\item (12) \(R_0(a + 2bt)\), \quad \(R_0 \left(a + \dfrac{b}{2\sqrt{t}}\right)\), \quad \(-\dfrac{R_0(a + 2bt)}{(1 + at + bt^2)^2}\) \quad or \quad \(\dfrac{R^2 (a + 2bt)}{R_0}\).
\item (13) \(1.4340(0.000014t - 0.001024)\), \quad \(-0.00117\), \quad \(-0.00107\), \quad \(-0.00097\).
\item (14) \(\dfrac{dE}{dl} = b + \dfrac{k}{i}\), \quad \(\dfrac{dE}{di} = -\dfrac{c + kl}{i^2}\).
\end{enumerate}
\par\bigskip\noindent{\Large\bfseries [演習 IV の解答]}\par\smallskip
\begin{enumerate}
\item (1) \(17 + 24x\); \quad \(24\).
\item (2) \(\dfrac{x^2 + 2ax - a}{(x + a)^2}\); \quad \(\dfrac{2a(a + 1)}{(x + a)^3}\).
\item (3) \(1 + x + \dfrac{x^2}{1 \times  2} + \dfrac{x^3}{1 \times  2 \times  3}\); \quad \(1 + x + \dfrac{x^2}{1 \times  2}\).
\item (4) (\textit{Exercises III.}):
\begin{itemize}
\item[(1)] (\textit{a}) \(\dfrac{d^2 y}{dx^2} = \dfrac{d^3 y}{dx^3} = 1 + x + \frac{1}{2}x^2 + \frac{1}{6} x^3 + \ldots\). \\
  (\textit{b}) \(2a\), \(0\). \hfil
  (\textit{c}) \(2\), \(0\). \hfil
  (\textit{d}) \(6x + 6a\), \(6\).
\item[(2)] \(-b\), \(0\). \hfil (3) \(2\), \(0\). \hfil
\item[(4)] \(\begin{gathered}[t]
    56440x^3 - 196212x^2 - 4488x + 8192. \\
    169320x^2 - 392424x - 4488.
    \end{gathered}\)
\item[(5)] \(2\), \(0\). \hfil (6) \(371.80453x\), \(371.80453\). \hfil
\item[(7)] \(\dfrac{30}{(3x + 2)^3}\), \quad \(-\dfrac{270}{(3x + 2)^4}\).
\end{itemize}
\par\medskip\noindent\textbf{\normalfont(\textit{Examples}, p. 268):}
\begin{itemize}
\item[(1)] \(\dfrac{6a}{b^2} x\), \quad \(\dfrac{6a}{b^2}\). \hfil
(2) \(\dfrac{3a \sqrt{b}} {2 \sqrt{x}} - \dfrac{6b \sqrt[3]{a}}{x^3}\), \quad
\(\dfrac{18b \sqrt[3]{a}}{x^4} - \dfrac{3a \sqrt{b}}{4 \sqrt{x^3}}\).
\item[(3)] \(\dfrac{2}{\sqrt[3]{\theta^8}} - \dfrac{1.056}{\sqrt[5]{\theta^{11}}}\), \quad
\(\dfrac{2.3232}{\sqrt[5]{\theta^{16}}} - \dfrac{16}{3 \sqrt[3]{\theta^{11}}}\).
\item[(4)] \(\begin{gathered}[t]
  810t^4 - 648t^3 + 479.52t^2 - 139.968t + 26.64. \\
  3240t^3 - 1944t^2 + 959.04t - 139.968.
  \end{gathered}\)
\item[(5)] \(12x + 2\), \(12\). \hfil
(6) \(6x^2 - 9x\), \quad \(12x - 9\). \hfil
\end{itemize}
\item[(7)]
\[
\frac{3}{4} \left( \frac{1}{\sqrt{\theta}} + \frac{1}{\sqrt{\theta^5}} \right)
+ \frac{1}{4} \left( \frac{15}{\sqrt{\theta^7}} - \frac{1}{\sqrt{\theta^3}} \right).
\]
\[
\frac{3}{8} \left( \frac{1}{\sqrt{\theta^5}} - \frac{1}{\sqrt{\theta^3}} \right)
- \frac{15}{8} \left( \frac{7}{\sqrt{\theta^9}} + \frac{1}{\sqrt{\theta^7}} \right).
\]

\item[(2)] 64; 147.2; and 0.32 フィート每秒.

\item[(3)] \( x = a - gt \); \(\ddot{x} = -g\).

\item[(4)] 45.1 フィート每秒.

\item[(5)] 12.4 フィート每秒每秒.\quad はい.

\item[(6)] 角速度 \(= 11.2\) ラジアン每秒; 角加速度 \(= 9.6\) ラジアン每秒每秒.

\item[(7)] \( v = 20.4t^2 - 10.8 \).\quad \( a = 40.8t \).\quad 172.8インチ毎秒, 122.4インチ毎秒平方.

\item[(8)] \( v = \frac{1}{30 \sqrt[3]{(t - 125)^2}} \),\quad \( a = - \frac{1}{45 \sqrt[3]{(t - 125)^5}} \).

\item[(9)] \( v = 0.8 - \frac{8t}{(4 + t^2)^2} \),\quad \( a = \frac{24t^2 - 32}{(4 + t^2)^3} \),\quad 0.7926 および 0.00211.

\item[(10)] \( n = 2 \), \( n = 11 \).

\par\smallskip

\end{enumerate}
\par\bigskip\noindent{\Large\bfseries [演習 VI の解答]}\par\smallskip

\begin{enumerate}
\item[(1)]  \(\frac{x}{\sqrt{ x^2 + 1}}\).
\item[(2)]  \(\frac{x}{\sqrt{ x^2 + a^2}}\).
\item[(3)]  \(- \frac{1}{2 \sqrt{(a + x)^3}}\).

\item[(4)]  \(\frac{ax}{\sqrt{(a - x^2)^3}}\).
\item[(5)]  \(\frac{2a^2 - x^2}{x^3 \sqrt{ x^2 - a^2}}\).

\item[(6)]  \(\frac{\frac{3}{2} x^2 \left[ \frac{8}{9} x \left( x^3 + a \right) - \left( x^4 + a \right) \right]}{(x^4 + a)^{\frac{2}{3}} (x^3 + a)^{\frac{3}{2}}}\)
\item[(7)]  \(\frac{2a \left(x - a \right)}{(x + a)^3}\).

\item[(8)]  \(\frac{5}{2} y^3\).
\item[(9)]  \(\frac{1}{(1 - \theta) \sqrt{1 - \theta^2}}\).
\par\smallskip

\end{enumerate}
\par\bigskip\noindent{\Large\bfseries [演習 VII の解答]}\par\smallskip

\begin{enumerate}
\item[(1)]  \(\frac{dw}{dx} = \frac{3x^2 \left( 3 + 3x^3 \right)} {27 \left(\frac{1}{2} x^3 + \frac{1}{4} x^6 \right)^3}\).

\item[(2)]  \(\frac{dv}{dx} = - \frac{12x}{\sqrt{1 + \sqrt{2} + 3x^2} \left(\sqrt{3} + 4 \sqrt{1 + \sqrt{2} + 3x^2}\right)^2}\).

\item[(3)]  \(\frac{du}{dx} = - \frac{x^2 \left(\sqrt{3} + x^3 \right)} {\sqrt{ \left[ 1 + \left( 1 + \dfrac{x^3}{\sqrt{3}} \right) ^2 \right]^3}} \)
\par\smallskip

\end{enumerate}
\par\bigskip\noindent{\Large\bfseries [演習 VIII の解答]}\par\smallskip

\begin{enumerate}
\item[(2)] 1.44.
\par\smallskip

\item[(4)] $\dfrac{dy}{dx} = 3x^2 + 3$; 数値は：$3$, $3 \frac{3}{4}$, $6$, $15$。

\item[(5)] $± \sqrt{2}$。

\item[(6)] $\dfrac{dy}{dx} = - \dfrac{4}{9} \dfrac{x}{y}$. 接線の傾きは$x = 0$で$0$であり、$x = 1$では$\mp \dfrac{1}{3 \sqrt{2}}$。

\item[(7)] $m = 4$, $n = -3$。

\item[(8)] $x = 1$, $x = -3$での交点。角度は$153^\circ\;26'$, $2^\circ\;28'$。

\item[(9)] $x = 3.57$, $y = 3.50$での交点。角度は$16^\circ\;16'$。

\item[(10)] $x = \frac{1}{3}$, $y = 2 \frac{1}{3}$, $b = -\frac{5}{3}$。
\par\smallskip

\end{enumerate}
\par\bigskip\noindent{\Large\bfseries [演習 IX の解答]}\par\smallskip

\begin{enumerate}
\item[(1)] 最小値：$x = 0$, $y = 0$; 最大値：$x = -2$, $y = -4$。

\item[(2)] $x = a$。

\item[(4)] 面積は$25 \sqrt{3}$平方インチ。

\item[(5)] $\dfrac{dy}{dx} = - \dfrac{10}{x^2} + \dfrac{10}{(8 - x)^2}$; $x = 4$; $y = 5$。

\item[(6)] 最大値は$x = -1$で、最小値は$x = 1$である。

\item[(7)] 四つの辺の中点を結ぶ。

\item[(8)] $r = \frac{2}{3} R$, $r = \dfrac{R}{2}$, 最大値なし。

\item[(9)] $r = R \sqrt{\dfrac{2}{3}}$, $r = \dfrac{R}{\sqrt{2}}$, $r = 0.8506R$。

\item[(10)] 毎秒$\dfrac{8}{r}$平方フィートの速度で変化する。

\item[(11)] $r = \dfrac{R \sqrt{8}}{3}$。
\item[(12)] $n = \sqrt{\dfrac{NR}{r}}$。
\par\smallskip

\end{enumerate}
\par\bigskip\noindent{\Large\bfseries [演習 X の解答]}\par\smallskip

\begin{enumerate}
\item[(1)] 最大値：$x = -2.19$, $y = 24.19$; 最小値：$x = 1.52$, $y = -1.38$。

\item[(2)] $\dfrac{dy}{dx} = \dfrac{b}{a} - 2cx$; $\dfrac{d^2 y}{dx^2} = -2c$; $x = \dfrac{b}{2ac}$ (最大値)。

\item[(3)] (\textit{a}) 一つの最大値と二つの最小値。 \\
(\textit{b}) 一つの最大値。 ($x = 0$; 他の点は実数でない)。

\item[(4)] 最小値：$x = 1.71$, $y = 6.14$。

\item[(5)] 最大値：$x = -0.5$, $y = 4$。

\item[(6)] 最大値：$x = 1.414$, $y = 1.7675$. \\
最小値：$x = -1.414$, $y = 1.7675$。

\item[(7)] 最大値：$x = -3.565$, $y = 2.12$. \\
最小値：$x = +3.565$, $y = 7.88$。

\item[(8)] $0.4N$, $0.6N$。

\item[(9)] $x = \sqrt{\dfrac{a}{c}}$。

\item[(10)] 速度は$8.66$海里/時間。所要時間は$115.47$時間。 \\
最小費用は$p$$112ポンド.\textasciitilde{}$12シリング。

\item[(11)] $x = 7.5$での最大値と最小値：$y = ±5.414$。 (例題10参照)。

\item[(12)] 最小値：$x = \frac{1}{2}$, $y= 0.25$; 最大値：$x = - \frac{1}{3}$, $y= 1.408$。
\par\smallskip

\item[(1)] $\dfrac{2}{ x - 3} + \dfrac{1}{ x + 4}$.

\item[(2)] $\dfrac{1}{ x - 1} + \dfrac{2}{ x - 2}$.

\item[(3)] $\dfrac{2}{ x - 3} + \dfrac{1}{ x + 4}$.

\item[(4)] $\dfrac{5}{ x - 4} - \dfrac{4}{ x - 3}$.

\item[(5)] $\dfrac{19}{13(2x + 3)} - \dfrac{22}{13(3x - 2)}$.

\item[(6)] $\dfrac{2}{ x - 2} + \dfrac{4}{ x - 3} - \dfrac{5}{ x - 4}$.

\item[(7)] $\dfrac{1}{6(x - 1)} + \dfrac{11}{15(x + 2)} + \dfrac{1}{10(x - 3)}$.

\item[(8)] $\dfrac{7}{9(3x + 1)} + \dfrac{71}{63(3x - 2)} - \dfrac{5}{7(2x + 1)}$.

\item[(9)] $\dfrac{1}{3(x - 1)} + \dfrac{2x + 1}{3(x^2 + x + 1)}$.

\item[(10)] $x + \dfrac{2}{3(x + 1)} + \dfrac{1 - 2x}{3(x^2 - x + 1)}$.

\item[(11)] $\dfrac{3}{(x + 1)} + \dfrac{2x + 1}{x^2 + x + 1}$.

\item[(12)] $\dfrac{1}{ x - 1} - \dfrac{1}{ x - 2} + \dfrac{2}{(x - 2)^2}$.

\item[(13)] $\dfrac{1}{4(x - 1)} - \dfrac{1}{4(x + 1)} + \dfrac{1}{2(x + 1)^2}$.

\item[(14)] $\dfrac{4}{9(x - 1)} - \dfrac{4}{9(x + 2)} - \dfrac{1}{3(x + 2)^2}$.

\item[(15)] $\dfrac{1}{ x + 2} - \dfrac{x - 1}{ x^2 + x + 1} - \dfrac{1}{(x^2 + x + 1)^2}$.

\item[(16)] $\dfrac{5}{ x + 4} -\dfrac{32}{(x + 4)^2} + \dfrac{36}{(x + 4)^3}$.

\item[(17)] $\dfrac{7}{9(3x - 2)^2} + \dfrac{55}{9(3x - 2)^3} + \dfrac{73}{9(3x - 2)^4}$.

\item[(18)] $\dfrac{1}{6(x - 2)} + \dfrac{1}{3(x - 2)^2} - \dfrac{x}{6(x^2 + 2x + 4)}$.

\end{enumerate}
\par\bigskip\noindent{\Large\bfseries [演習 XII の解答]}\par\smallskip

\begin{enumerate}
\item[(1)] $ab(\epsilon^{ax} + \epsilon^{-ax})$.

\item[(2)] $2at + \dfrac{2}{t}$.

\item[(3)] $\log_\epsilon n$.

\item[(5)] $npv^{n-1}$.

\item[(6)] $\dfrac{n}{x}$.

\item[(7)] $\dfrac{3\epsilon^{- \frac{x}{x-1}}}{(x - 1)^2}$.

\item[(8)] $6x \epsilon^{-5x} - 5(3x^2 + 1)\epsilon^{-5x}$.

\item[(9)] $\dfrac{ax^{a-1}}{x^a + a}$.

\item[(10)] $\left(\dfrac{6x}{3x^2-1} + \dfrac{1}{2\left(\sqrt x + x\right)}\right) \left(3x^2-1\right)\left(\sqrt x + 1\right)$.

\item[(11)] $\dfrac{1 - \log_\epsilon \left(x + 3\right)}{\left(x + 3\right)^2}$.

\item[(12)] $a^x\left(ax^{a-1} + x^a \log_\epsilon a\right)$.

\item[(14)] Min.: $y = 0.7$ for $x = 0.694$.

\item[(15)] $\dfrac{1 + x}{x}$.

\item[(16)] $\dfrac{3}{x} (\log_\epsilon ax)^2$.
\par\smallskip

\end{enumerate}
\par\bigskip\noindent{\Large\bfseries [演習 XIII の解答]}\par\smallskip

\begin{enumerate}
\item[(1)]  Let $\dfrac{t}{T} = x$ ($\therefore t = 8x$), and use the Table on \Pageref[page]{littletable}.

\item[(2)]  $T = 34.627$; $159.46$ minutes.

\item[(3)]  Take $2t = x$; and use the Table on \Pageref[page]{littletable}.

\item[(5)]  (\textit{a}) $x^x (1 + \log_\epsilon x)$;\quad
(\textit{b}) $2x (\epsilon^x)^x$;\quad
(\textit{c}) $\epsilon^{x^x} \times  x^x (1 + \log_\epsilon x)$.

\item[(6)]  $0.14$ second.

\item[(7)]  (\textit{a}) $1.642$;\quad (\textit{b}) $15.58$.

\item[(8)]  $\mu = 0.00037$, $31^m \frac{1}{4}$. %[** Time units]

\item[(9)]  $i$ is $63.4$\% of\textasciitilde{}$i_0$, $220$\textasciitilde{}kilometres.

\item[(10)]  $0.133$, $0.145$, $0.155$, mean $0.144$; $-10.2$\%, $-0.9$\%, $+77.2$\%.

\item[(11)]  Min.\ for $x = \dfrac{1}{\epsilon}$.

\item[(12)]  Max.\ for $x = \epsilon$.

\item[(13)]  Min.\ for $x = \log_\epsilon a$.
\par\smallskip

\end{enumerate}
\par\bigskip\noindent{\Large\bfseries [演習 XIV の解答]}\par\smallskip

\begin{enumerate}
\item[(1)] (i) $\dfrac{dy}{d\theta} = A \cos \left( \theta - \dfrac{\pi}{2} \right)$;

(ii) $\dfrac{dy}{d\theta} = 2\sin\theta \cos\theta = \sin2\theta$ and $\dfrac{dy}{d\theta} = 2\cos2\theta$;

(iii) $\dfrac{dy}{d\theta} = 3\sin^2 \theta \cos\theta$ and $\dfrac{dy}{d\theta} = 3\cos3\theta$.

\item[(2)]  $\theta = 45^\circ$ or $\dfrac{\pi}{4}$ radians.

\item[(3)]  $\dfrac{dy}{dt} = -n \sin 2\pi nt$.

\item[(4)]  $a^x \log_\epsilon a \cos a^x$.

\item[(5)]  $\dfrac{\cos x}{\sin x} = \cotan x$

\item[(6)]  $18.2 \cos \left(x + 26^\circ \right)$.

\item[(7)] { The slope is $\dfrac{dy}{d\theta} = 100\cos\left(\theta - 15^\circ \right)$, which is a maximum
 when $(\theta -15^\circ) = 0$, or $\theta = 15^\circ$; the value of the slope
 being then ${}= 100$. When $\theta = 75^\circ$ the slope is
 $100\cos(75^\circ  - 15^\circ) = 100\cos 60^\circ = 100 \times  \frac{1}{2} = 50$.}

\item[(8)] $\begin{aligned}[t]
  \cos\theta \sin2\theta + 2\cos2\theta \sin\theta
  &= 2\sin\theta\left(\cos^2 \theta + \cos2\theta\right) \\
  &= 2\sin\theta\left(3\cos^2 \theta - 1\right).
  \end{aligned}$

\item[(9)] $amn\theta^{n-1} \tan^{m-1}\left(\theta^n\right)\sec^2 \theta^n$.

\item[(10)] $\epsilon^x \left(\sin^2 x + \sin2x\right)$;\quad $\epsilon^x \left(\sin^2 x + 2\sin2x + 2\cos2x\right)$.

\item (11) \begin{enumerate}[label=(\roman*)]
    \item $\dfrac{dy}{dx} = \dfrac{ab}{(x + b)^2}$;
    \item $\dfrac{a}{b} \epsilon^{-\frac{x}{b}}$;
    \item $\dfrac{1}{90^\circ} \times \dfrac{ab}{b^2 + x^2}$.
\end{enumerate}

\item (12) \begin{enumerate}[label=(\roman*)]
    \item $\dfrac{dy}{dx} = \sec x \tan x$;
    \item $\dfrac{dy}{dx} = - \dfrac{1}{\sqrt{1 - x^2}}$;
    \item $\dfrac{dy}{dx} = \dfrac{1}{1 + x^2}$;
    \item $\dfrac{dy}{dx} = \dfrac{1}{x \sqrt{x^2 - 1}}$;
    \item $\dfrac{dy}{dx} = \dfrac{\sqrt{3 \sec x} (3 \sec^2 x - 1)}{2}$.
\end{enumerate}

\item (13) $\dfrac{dy}{d\theta} = 4.6 (2\theta + 3)^{1.3} \cos (2\theta + 3)^{2.3}$.

\item (14) $\dfrac{dy}{d\theta} = 3\theta^2 + 3 \cos (\theta + 3) - \log_\epsilon 3 (\cos \theta \times 3^{\sin \theta} + 3\theta)$.

\item (15) $\theta = \cot \theta$; $\theta = \pm 0.86$; $+\theta$ のとき最大、$-\theta$ のとき最小。

\item (1) $x^3 - 6x^2 y - 2y^2$; $\dfrac{1}{3} - 2x^3 - 4xy$.

\item (2) $2xyz + y^2 z + z^2 y + 2xy^2 z^2$;
     $2xyz + x^2 z + xz^2 + 2x^2 yz^2$;
     $2xyz + x^2 y + xy^2 + 2x^2 y^2 z$.

\item (3) $\dfrac{1}{r} \left( (x - a) + (y - b) + (z - c) \right) = \dfrac{(x + y + z) - (a + b + c)}{r}$; $\dfrac{3}{r}$.

\item (4) $dy = vu^{v-1} du + u^v \log_\epsilon u dv$.

\item (5) $dy = 3 \sin v u^2 du + u^3 \cos v dv$,\\
     $dy = u \sin x^{u-1} \cos x dx + (\sin x)^u \log_\epsilon \sin x du$,\\
     $dy = \dfrac{1}{v} \dfrac{1}{u} du - \log_\epsilon u \dfrac{1}{v^2} dv$.

\item (7) 最小値は $x = y = -\dfrac{1}{2}$ で達成される。

\item (8) (\textit{a}) 長さ $2$ フィート、幅 = 深さ = $1$ フィート、体積 = $2$ 立方フィート。\\
     (\textit{b}) 半径 = $\dfrac{2}{\pi}$ フィート = $7.46$ インチ、長さ = $2$ フィート、体積 = $2.54$。

\item (9) 三つの部分は等しく、積が最大となる。

\item (10) 最小値は $x = y = 1$ で達成される。

\item (11) 最小値は $x = \dfrac{1}{2}$ および $y = 2$ で達成される。

\item (12) 頂点の角度 = $90^\circ$; 等しい辺 = 長さ = $\sqrt[3]{2V}$。
\end{enumerate}
\par\bigskip\noindent{\Large\bfseries [演習 XVII の解答]}\par\smallskip

\begin{enumerate}
\item[(1)] \(\dfrac{4\sqrt{a} x^{\frac{3}{2}}}{3} + C\).

\item[(2)] \(-\dfrac{1}{x^3} + C\).

\item[(3)] \(\dfrac{x^4}{4a} + C\).

\item[(4)] \(\tfrac{1}{3} x^3 + ax + C\).

\item[(5)] \(-2x^{-\frac{5}{2}} + C\).

\item[(6)] \(x^4 + x^3 + x^2 + x + C\).

\item[(7)] \(\dfrac{ax^2}{4} + \dfrac{bx^3}{9} + \dfrac{cx^4}{16} + C\).

\item[(8)] \(\dfrac{x^2 + a}{x + a} = x - a + \dfrac{a^2 + a}{x + a}\) と除算により得られる。従って、答えは \(\dfrac{x^2}{2} - ax + (a^2 + a)\log_\epsilon (x + a) + C\) である。（参照ページ：\pageref{intex1} および\textasciitilde{}\pageref{intex2}）

\item[(9)] \(\dfrac{x^4}{4} + 3x^3 + \dfrac{27}{2} x^2 + 27x + C\).

\item[(10)] \(\dfrac{x^3}{3} + \dfrac{2 - a}{2} x^2 - 2ax + C\).

\item[(11)] \(a^2(2x^{\frac{3}{2}} + \tfrac{9}{4} x^{\frac{4}{3}}) + C\).

\item[(12)] \(-\tfrac{1}{3} \cos\theta - \tfrac{1}{6} \theta + C\).

\item[(13)] \(\dfrac{\theta}{2} + \dfrac{\sin 2a\theta}{4a} + C\).

\item[(14)] \(\dfrac{\theta}{2} - \dfrac{\sin 2\theta}{4} + C\).

\item[(15)] \(\dfrac{\theta}{2} - \dfrac{\sin 2a\theta}{4a} + C\).

\item[(16)] \(\tfrac{1}{3} \epsilon^{3x}\). % [F1: +C?]

\item[(17)] \(\log(1 + x) + C\).

\item[(18)] \(-\log_\epsilon (1 - x) + C\).
\par\smallskip

\end{enumerate}
\par\bigskip\noindent{\Large\bfseries [演習 XVIII の解答]}\par\smallskip

\begin{enumerate}
\item[(1)] 面積 = 60; 平均曲線値 = 10.

\item[(2)] 面積 = \(a \times  2a \sqrt{a}\) の \(\frac{2}{3}\).

\item[(3)] 面積 = 2; 平均曲線値 = \(\dfrac{2}{\pi} = 0.637\).

\item[(4)] 面積 = 1.57; 平均曲線値 = 0.5.

\item[(5)] 0.572, 0.0476.

\item[(6)] 体積 = \(\pi r^2 \dfrac{h}{3}\).

\item[(7)] 1.25.

\item[(8)] 79.4.

\item[(9)] 体積 = 4.9348; 表面積の面積 = 12.57（0から \(\pi\) まで）.

\item[(10)] \(a\log_\epsilon a\), \(\dfrac{a}{a - 1} \log_\epsilon a\).

\item[(12)] 算術平均 = 9.5; 平方根平均 = 10.85.

\item[(13)] 平方根平均 = \(\dfrac{1}{\sqrt{2}} \sqrt{A_1^2 + A_3^2}\); 算術平均 = 0.

最初のものは、少し複雑な積分に関連しています。二次平均は次のように定義されます：
\[
\sqrt{\dfrac{1}{2\pi} \int_0^{2\pi} (A_1 \sin x + A_3 \sin 3x)^2\, dx}.
\]
ここで、次の積分
\[
\int (A_1^2 \sin^2 x + 2A_1 A_3 \sin x \sin 3x + A_3^2 \sin^2 3x)\, dx
\]
は、$\sin^2 x$ を
\[
\dfrac{1 - \cos 2x}{2}
\]
と書き換えることでより容易に計算できます。また、$2\sin x \sin 3x$ は $\cos 2x - \cos 4x$ と書き換えられ、$\sin^2 3x$ は
\[
\dfrac{1 - \cos 6x}{2}
\]
と書き換えることができます。

これらの置換を行い、積分を計算すると（参照：\Pageref{cosax}）次のようになります：
\[
\dfrac{A_1^2}{2} \left( x - \dfrac{\sin 2x}{2} \right)
 + A_1 A_3 \left( \dfrac{\sin 2x}{2} - \dfrac{\sin 4x}{4} \right)
 + \dfrac{A_3^2}{2} \left( x - \dfrac{\sin 6x}{6} \right).
\]

下限において $x = 0$ を代入すると、すべてが消え、上限において $x = 2\pi$ を代入すると $A_1^2 \pi + A_3^2 \pi$ となります。したがって、答えは得られます。

(14) 面積は62.6平方単位です。平均縦座標は10.42です。

(16) 436.3（この立体は梨型です）。

解答：

(1) $\dfrac{x\sqrt{a^2 - x^2}}{2} + \dfrac{a^2}{2} \sin^{-1} \dfrac{x}{a} + C$。

(2) $\dfrac{x^2}{2}(\log_\epsilon x - \tfrac{1}{2}) + C$。

(3) $\dfrac{x^{a+1}}{a + 1} \left(\log_\epsilon x - \dfrac{1}{a + 1}\right) + C$。

(4) $\sin \epsilon^x + C$。

(5) $\sin(\log_\epsilon x) + C$。

(6) $\epsilon^x (x^2 - 2x + 2) + C$。

(7) $\dfrac{1}{a + 1} (\log_\epsilon x)^{a+1} + C$。
 

(8) $\log_\epsilon(\log_\epsilon x) + C$。

(9) $2\log_\epsilon(x - 1) + 3\log_\epsilon(x + 2) + C$。

(10) $\frac{1}{2} \log_\epsilon(x - 1) + \frac{1}{5} \log_\epsilon(x - 2) + \frac{3}{10} \log_\epsilon(x + 3) + C$。

(11) $\dfrac{b}{2a} \log_\epsilon \dfrac{x - a}{x + a} + C$。

(12) $\log_\epsilon \dfrac{x^2 - 1}{x^2 + 1} + C$。

(13) $\frac{1}{4} \log_\epsilon \dfrac{1 + x}{1 - x} + \frac{1}{2} \arctan x + C$。

(14) $\dfrac{1}{\sqrt{a}} \log_\epsilon \dfrac{\sqrt{a} - \sqrt{a - bx^2}}{x\sqrt{a}}$。（$\dfrac{1}{x} = v$ と置き換えて、結果では $\sqrt{v^2 - \dfrac{b}{a}} = v - u$ とする）

あなたは現在の解答と、元の表現に戻って確認するべきです。各段階で自己用の更多例子应自行制造，以测试自己的能力。在积分时，可以通过対答案求导来验证是否能回到原始表达式。

多くの書籍には練習のために例題が含まれています。ここでは二つだけ挙げます：R. G. Blaine の『微積分及其应用』と F. M. Saxelby の『実用数学教程』。

---

グラスゴー：ユニバーシティ・プリスンズ・バイ・ロバート・マクレホース・アンド・コ・リミテッドで印刷されました。

---

以下は選択された数学作品の目録です：

---

**目録**

1. 『微積分入门』グラフィカル方法为基础。由 G. A. Gibson 博士，M.A., LL.D 编著。3先令6便士。
2. 『初等微积分教程』附有几何、力学和物理的插图。由 G. A. Gibson 博士，M.A., LL.D 编著。7先令6便士。
3. 『微分入门』由 J. Edwards 博士，M.A 编著。4先令6便士。
4. 『积分入门』附有微分方程研究介绍。由 Joseph Edwards 博士，M.A 编著。4先令6便士。
5. 『微积分简易教程』被称作微分和积分的美丽計算方法的一种非常简单的介绍。由 F. R. S 编著。2先令净价。新版，附有许多例题。
6. 『初等微积分教程』由 W. F. Osgood 博士，Ph.D 编著。8先令6便士（净价）。
7. 『工程师用实用积分法』由 A. S. Percival 博士，M.A 编著。2先令6便士（净价）。
8. 『微分教程』附有应用和大量例题。由 Joseph Edwards 博士，M.A 编著。14先令。
9. 『技术学校与学院用微积分教程』由 P. A. Lambert 博士，M.A 编著。7先令6便士。
10. 『微积分及其应用』由 Sir A. G. Greenhill, F.R.S 编著。10先令6便士。

---

伦敦：麦克米伦有限公司出版。

---

原著書籍のリストです。

1. 《積分学の Treatise》by I. Todhunter, F.R.S. 10s. 6d. キー版 by H. St. J. Hunter, M.A. 10s. 6d.
2. 《微分学と積分学の要素の Treatise》with numerous examples by I. Todhunter, F.R.S. 10s. 6d. キー版 by H. St. J. Hunter, M.A. 10s. 6d.
3. 《常微分方程式》An Elementary Text-book by James Morris Page, Ph.D. 6s. 6d.
4. 《現代理論の導入》Equations by Prof. F. Cajori, Ph.D. 7s. 6d. net.
5. 《微分方程式の Treatise》by Andrew Russell Forsyth, Sc.D., LL.D. Fourth Edition. 14s. net.
6. 《微分方程式の短いコース》by Prof. Donald F. Campbell, Ph.D. 4s. net.
7. 《四元数の Manual》by C. J. Joly, M.A., D.Sc., F.R.S. 10s. net.
8. 《行列式理論の歴史的順序での発展》Vol. I. Part I. General Determinants, up to 1841. Part II. Special Determinants, up to 1841. 17s. net. Vol. II. The Period 1841 to 1860. 17s. net. by T. Muir, M.A., LL.D., F.R.S.
9. 《無限級数の理論への入門》by T. J. I'a Bromwich, M.A., F.R.S. 15s. net.
10. 《Fourier級数と積分、および熱伝導の数学的理論への入門》by Prof. H. S. Carslaw, M.A., D.Sc., F.R.S.E. 14s. net.

---

ロンドン：マクミラン・アンド・カンパニー・リミテッド

---

更新版が前の版を置き換えるとします—古い版は再命名されます。

\end{enumerate}




\backmatter

\chapter{訳者・数理校閲後記}

本書は、微積分学教育の歴史的名著であるシルバヌス・P・トンプソンの『誰にもわかる微積分』（\textit{Calculus Made Easy}, 1914年改訂第2版）の全21章および序文、プロローグ、エピローグ、標準公式一覧、全演習問題解答に至る全26セクションを、Project Gutenberg（\#33283）の公式原著TeXソースを底本として完全直訳・数理校閲したものである。

微積分という数学の道具が本来持っていた「変化を捉える驚くべき直観の力」と「誰にでも使える道具としての親しみやすさ」を、トンプソンの軽妙で力強い筆致そのままに現代の日本語読者へ届けることができた。

\begin{flushright}
2026年10月\\
古典数理復刻プロジェクト
\end{flushright}

\end{document}
